\documentclass[11pt]{article}

\usepackage{amsmath,amssymb,amsthm,amsfonts,latexsym,xspace,graphicx,float,mathtools,braket,}
\usepackage[letterpaper,margin=1in]{geometry}
\usepackage[backref,colorlinks,citecolor=blue,,linkcolor=magenta,bookmarks=true]{hyperref}
\usepackage{enumitem}
\usepackage{aliascnt}
\usepackage[nameinlink,capitalize]{cleveref}

\newtheorem{theorem}{Theorem}[section]
\newcommand{\newsharedtheorem}[3]{%
  \newaliascnt{#1}{theorem}%
  \newtheorem{#1}[#1]{#2}%
  \aliascntresetthe{#1}%
  \crefname{#1}{#2}{#3}%
  \Crefname{#1}{#2}{#3}%
}
\newtheorem*{namedtheorem}{\theoremname}
\newcommand{\theoremname}{testing}
\newenvironment{named}[1]{ \renewcommand{\theoremname}{#1} \begin{namedtheorem}} {\end{namedtheorem}}
\newsharedtheorem{lemma}{Lemma}{Lemmas}
\newsharedtheorem{claim}{Claim}{Claims}
\newsharedtheorem{proposition}{Proposition}{Propositions}
\newsharedtheorem{fact}{Fact}{Facts}
\newsharedtheorem{corollary}{Corollary}{Corollaries}
\newsharedtheorem{conjecture}{Conjecture}{Conjectures}
\newsharedtheorem{infthm}{Informal Theorem}{Informal Theorems}

\theoremstyle{definition}
\newsharedtheorem{definition}{Definition}{Definitions}
\newsharedtheorem{infdef}{Informal Definition}{Informal Definitions}
\newsharedtheorem{remark}{Remark}{Remarks}
\newsharedtheorem{notation}{Notation}{Notations}
\newsharedtheorem{example}{Example}{Examples}
\newsharedtheorem{examples}{Examples}{Examples}
\renewcommand{\Pr}{\mathop{\bf Pr\/}}
\newcommand{\E}{\mathop{\bf E\/}}
\newcommand{\tr}{\mathrm{tr}} \newcommand{\Tr}{\tr}
\newcommand{\poly}{\mathrm{poly}}
\newcommand{\supp}{\mathrm{supp}}
\newcommand{\C}{\mathbb C}
\newcommand{\F}{\mathbb F}
\newcommand{\NP}{\mathsf{NP}}
\newcommand{\NEXP}{\mathsf{NEXP}}
\newcommand{\NEEXP}{\mathsf{NEEXP}}
\newcommand{\MIP}{\mathsf{MIP}}
\newcommand{\QMA}{\mathsf{QMA}}
\newcommand{\eps}{\epsilon}
\newcommand{\ot}{\otimes}
\newcommand{\calA}{\mathcal{A}}
\newcommand{\calB}{\mathcal{B}}
\newcommand{\calD}{\mathcal{D}}
\newcommand{\calH}{\mathcal{H}}
\newcommand{\calO}{\mathcal{O}}
\newcommand{\calP}{\mathcal{P}}
\newcommand{\calX}{\mathcal{X}}
\newcommand{\bone}{\boldsymbol{1}}
\newcommand{\ba}{\boldsymbol{a}}
\newcommand{\bb}{\boldsymbol{b}}
\newcommand{\boldf}{\boldsymbol{f}}
\newcommand{\bg}{\boldsymbol{g}}
\newcommand{\bi}{\boldsymbol{i}}
\newcommand{\bell}{\boldsymbol{\ell}}
\newcommand{\br}{\boldsymbol{r}}
\newcommand{\bs}{\boldsymbol{s}}
\newcommand{\bt}{\boldsymbol{t}}
\newcommand{\bu}{\boldsymbol{u}}
\newcommand{\bv}{\boldsymbol{v}}
\newcommand{\bw}{\boldsymbol{w}}
\newcommand{\bx}{{\boldsymbol{x}}}
\newcommand{\by}{\boldsymbol{y}}
\newcommand{\abs}[1]{\lvert #1 \rvert}
\newcommand{\polyfunc}[3]{\calP(#1, #2, #3)}
\newcommand{\polymeas}[3]{\mathrm{PolyMeas}(#1, #2, #3)}
\newcommand{\polysub}[3]{\mathrm{PolySub}(#1, #2, #3)}


\title{Quantum soundness of the classical low individual degree test}

\author{
Zhengfeng Ji\thanks{zhengfeng.ji@uts.edu.au}\\
\small{\sl University of Technology Sydney}
\and Anand Natarajan\thanks{anandn@mit.edu.  Most of this work performed while affiliated with
the California Institute of Technology.}\\
 \small{\sl Massachusetts Institute of Technology}
 \and Thomas Vidick\thanks{vidick@caltech.edu}\\
 \small{\sl California Institute of Technology}
 \and John Wright\thanks{wright@cs.utexas.edu. Most of this work performed while also affiliated with
the California Institute of Technology.}\\
 \small{\sl University of Texas at Austin}
 \and Henry Yuen\thanks{hyen@cs.toronto.edu}\\
 \small{\sl University of Toronto}\vspace*{10pt}
}

\date{}

\begin{document}

\maketitle

\begin{abstract}
Low degree tests play an important role in classical complexity theory,
serving as basic ingredients in foundational results such as $\MIP = \NEXP$~\cite{BFL91}
and the PCP theorem~\cite{AS98,ALM+98}.
Over the last ten years,
versions of these tests which are sound against quantum provers
have found increasing applications to the study of nonlocal games and the complexity class~$\MIP^*$.
The culmination of this line of work is the result $\MIP^* = \mathsf{RE}$~\cite{JNV+20}.

One of the key ingredients in the first reported proof of  $\MIP^* = \mathsf{RE}$ is a two-prover variant of the low degree test, initially
shown to be sound against multiple quantum provers in~\cite{Vid16}.
Unfortunately a mistake was recently discovered in the latter result, invalidating the main result of~\cite{Vid16} as well as its use in subsequent works, including~\cite{JNV+20}.

We analyze a variant of the low degree test called the low individual degree test.
We give a simplified proof that the two-player version of this test is
sound against quantum provers.  The quantitative bound proved here is
$21000m^4\max\{1,d^2\}(\eps^{1/64}+(d/q)^{1/64})$.
The historical motivation is to repair applications of low-degree
testing to~$\MIP^*$; this exposition does not derive those applications
from its weaker dependence on~$d$.
\end{abstract}

\newpage

\tableofcontents
\vfill
\thispagestyle{empty}
\newpage

\section{Introduction}\label{sec:intro}

An $m$-variate polynomial over the finite field~$\F_q$ is a function $g:\F_q^m \rightarrow \F_q$ of the form
\begin{equation*}
g(x_1, \ldots, x_m) = \sum_{i_1, \ldots, i_m} c_{i_1, \ldots, i_m} \cdot x_1^{i_1} \cdots x_{m}^{i_m},
\end{equation*}
where each coefficient~$c_{i_1, \ldots, i_m}$ is an element of~$\F_q$.
We say that~$g$ has \emph{total degree $d$}
(or \emph{degree $d$}, for short) if $i_1 + \cdots + i_m \leq d$ for each nonzero coefficient $c_{i_1, \ldots, i_m}$,
and \emph{individual degree $d$} if $i_1, \ldots, i_m \leq d$ for each nonzero coefficient $c_{i_1, \ldots, i_m}$.
Low-degree polynomials have a variety of properties which make them useful in theoretical computer science,
chief among which is their \emph{distance}:
by the Schwartz-Zippel lemma, two nonequal degree $d$ polynomials~$g$ and~$h$ agree on at most a $d/q$ fraction of the points in~$\F_q^m$.

\emph{Low (individual) degree testing} refers to the task
of verifying that an unknown function $g:\F_q^m \rightarrow \F_q$
is representable as a polynomial of (individual) degree~$d$
by querying~$g$ on a small number of points $u \in \F_q^m$.
There is a pair of canonical tests for doing so
known as the \emph{surface-versus-point low-degree test}
and the \emph{low individual degree test}.
It is common to frame these tests
as games between a referee and two provers.
In this setting, the surface-versus-point low degree test, parameterized by an integer $1\leq k\leq m$, is performed by the verifier as follows.
\begin{enumerate}
\item Select~$\bu \sim \F_q^m$ uniformly at random. Give it to Prover~$\mathrm{A}$. They respond with a value~$\ba \in \F_q$.
\item Select a uniformly random $k$-dimensional affine surface~$\bs$ in~$\F_q^m$ containing~$\bu$.
	Give it to Prover~$\mathrm{B}$. They respond with a degree-$d$ $k$-variate polynomial $\boldf:\bs\rightarrow \F_q$.
\item Accept if $\boldf(\bu) = \ba$.
\end{enumerate}
This test is motivated by the fact that a degree-$d$ polynomial
$g:\F_q^m\rightarrow\F_q$ restricts to a degree-$d$ polynomial
on every affine $k$-dimensional surface~$s$.
Hence, if~$g$ is degree-$d$,
then the provers can win with probability~$1$ by always replying with~$\ba = g(\bu)$ and $\boldf = g|_{\bs}$.
\emph{Soundness} of the low-degree test refers to the converse statement,
namely that players who succeed with high probability
must be responding based on a low-degree polynomial.
This is formalized as follows.

\begin{theorem}[Raz-Safra~\cite{RS97}]\label{thm:raz-safra}
Suppose deterministic classical Provers~$\mathrm{A}$ and~$\mathrm{B}$ pass the $k = 2$ surface-versus-point low-degree test with probability~$1-\eps$.
Then there exists a degree-$d$ polynomial $g:\F_q^m \rightarrow \F_q$ such that
\begin{equation*}
\Pr_{\bu \sim \F_q^m}[g(\bu) = \ba] \geq 1 - \eps - \poly(m) \cdot \poly(d/q).
\end{equation*}
\end{theorem}

A similar ``local characterization" of low individual degree polynomials states that a polynomial~$g$ has individual degree $d$
if and only if $g|_{\ell}$ is a univariate degree-$d$ polynomial for all axis-parallel lines~$\ell$.
An axis-parallel line is a line of the form
$\ell = \{u+t e_i \mid t \in \F_q\}$, for $u \in \F_q^m$ and
$i \in \{1, \ldots, m\}$.
Motivated by this, the low individual degree test follows the same outline as the surface-versus-point low degree test
except with the second step substituted with the following.
\begin{enumerate}
\setcounter{enumi}{1}
\item Select a uniformly random axis-parallel line $\bell$ in~$\F_q^m$ containing~$\bu$.
	Give it to Prover~$\mathrm{B}$. They respond with a degree-$d$ univariate polynomial $\boldf:\bell\rightarrow \F_q$.
\end{enumerate}
When $d = 1$, the low individual degree test
is called the \emph{multilinearity test} because a polynomial with individual degree $d = 1$ is a multilinear polynomial.
The multilinearity  and low individual degree tests were
first introduced and proven sound by Babai, Fortnow, and Lund in~\cite{BFL91}.
The analysis of its soundness was then improved by~\cite{AS98} and then further sharpened by~\cite{FHS94}.
The best bound follows from the work of Polishchuk and Spielman~\cite{PS94};
their work considers only the bivariate $m = 2$ case,
but extending it to the multivariate case yields the following result.
\begin{theorem}[Polishchuk-Spielman~\cite{PS94}]\label{thm:classical-test-soundness}
Suppose deterministic classical Provers~$\mathrm{A}$ and~$\mathrm{B}$ pass the low individual degree test with probability~$1-\eps$.
Then there exists a polynomial $g:\F_q^m \rightarrow \F_q$ with individual degree~$d$ such that
\begin{equation*}
\Pr_{\bu \sim \F_q^m}[g(\bu) = \ba] \geq 1 -  \poly(m) \cdot (\poly(\eps) + \poly(d/q)).
\end{equation*}
\end{theorem}
We note that the soundness error the low individual test gives is actually worse than the low degree test,
because the low individual degree function~$g$ is only $\poly(m) \cdot (\poly(\eps) + \poly(d/q))$ close to Player~$\mathrm{A}$'s strategy,
rather than $\eps + \poly(m) \cdot \poly(d/q)$.
We will discuss this weakness of the low individual degree test below.

The multilinearity test, low individual degree test, and low-degree test
form a sequence
in which each test generally enables more applications than the previous one.
The multilinearity test can be used to show that $\MIP = \NEXP$ using a polynomial number of rounds~\cite{BFL91},
the low individual degree test can reduce the number of rounds to~$1$,
and the low degree test can be used to ``scale this result down''
and prove the PCP theorem, i.e.\ $\NP = \MIP[O(\log(n)), O(1)]$~\cite{AS98,ALM+98}.

\subsection{Quantum soundness of the low degree tests}

The work of Ito and Vidick~\cite{IV12}
initiated a program of studying these tests
in the case when the players are quantum,
as a means of proving bounds on the complexity class~$\MIP^*$.
Because the provers are quantum,
they are allowed to share an entangled state,
a resource which could potentially allow them to ``cheat" the test
and win without using a low-degree polynomial.
The goal of this program is to show that this is not possible.
In other words, the goal is to show that these tests are \emph{quantum sound},
which means that provers who succeed with high success probability
must answer their questions according to a low (individual) degree polynomial,
even if they are allowed to share quantum entanglement.\footnote{We note that \emph{quantum soundness of the low-degree tests}, which is the focus of this work,
is distinct from \emph{soundness of the quantum low-degree test}.
The ``quantum low-degree test'' is a particular test introduced by Natarajan and Vidick in~\cite{NV18a}
which gets its name from the prominent role that the low-degree test plays as a subroutine,
and its ``soundness'' is simply the result that they prove about it.
}

Correctly formalizing the notion of quantum soundness is a subtle task,
as quantum provers can in fact ace these tests
using a broader class of strategies than their classical counterparts.
For example, the two provers can use their quantum state~$\ket{\psi}$ to simulate shared randomness,
which they can use to sample a random low-degree polynomial~$\bg$  to answer their questions with;
what makes this still acceptable is that~$\bg$ depends only on their shared randomness and not their questions.
The correct formalization of quantum soundness was identified by Ito and Vidick~\cite{IV12},
which states the following:
suppose Provers~$\mathrm{A}$ and~$\mathrm{B}$ pass the low-degree test with probability close to~$1$.
For each point question~$u \in \F_q^m$,
let $A^{u} = \{A^{u}_a\}$ be the measurement that Prover~$\mathrm{A}$ applies to
their share of~$\ket{\psi}$ to produce the answer $\ba \in \F_q$.
Then the test being quantum sound
means that there should be a measurement $G = \{G_g\}$, independent of~$u \in \F_q^m$,
which outputs degree-$d$ polynomials~$g$
and ``acts like~$A$''.
In other words, rather than measuring~$A^{u}$ to produce the outcome~$\ba \in \F_q$,
Prover~$\mathrm{A}$ could have simply measured~$G$, received the polynomial~$\bg$,
and outputted its evaluation at~$\bu$, i.e.\ the value $\bg(\bu)$.
We will measure the similarity between~$A$ and~$G$
by considering the experiment where Prover~$\mathrm{A}$ measures with~$A^{\bu}$ to produce~$\ba$,
Prover~$\mathrm{B}$ measures with~$G$ to produce~$\bg$, and we check if $\bg(\bu) = \ba$.
This entails studying the quantity
\begin{equation*}
\E_{\bu \sim \F_q^m} \sum_{a \in \F_q} \sum_{g:g(\bu) = a}\bra{\psi} A^{\bu}_a \ot G_g \ket{\psi},
\end{equation*}
which we aim to show is as close to~$1$ as possible.
In this way, the provers' quantum advantage is limited to their ability to select a low-degree polynomial~$g$.

One additional quirk of the quantum setting
is that it has been historically useful to consider variants of these tests which feature more than two provers.
This allows one to use monogamy of entanglement to reduce the power that entanglement gives to the provers,
making it easier to show that a given test is quantum sound.
Low degree test results with fewer provers are more difficult to show and have more applications.

The program of showing that these tests are quantum sound was carried out for the $3$-prover multilinearity test by Ito and Vidick~\cite{IV12}
and for the $3$-prover low-degree test by Vidick~\cite{Vid16},
which was later improved to $2$-provers by Natarajan and Vidick~\cite{NV18b}.
This latter result led to a sequence of works
which culminated in the proof that $\mathsf{MIP}^* = \mathsf{RE}$
and the refutation of the Connes embedding conjecture in~\cite{JNV+20}.
We summarize this line of research in \Cref{fig:research}.

{
\floatstyle{boxed}
\restylefloat{figure}
\begin{figure}
\begin{tabular}{l c c c}
& Test shown & Complexity-theoretic  & Number \\
&  quantum-sound &  consequence & of provers  \\
			& &&\\
1. \cite{IV12}: & multilinearity test & $\NEXP \subseteq \mathsf{MIP}^*$ & 3\\[1ex]
2. \cite{Vid16}: & low-degree test & $\NP \subseteq \mathsf{MIP}^*[O(\log(n)), O(1)]$  &3\\[1ex]
3. \cite{NV18b}: & low-degree test & $\NP \subseteq \mathsf{MIP}^*[O(\log(n)), O(1)]$ & 2\\
			& &&\\
& Consequences of \cite{NV18b}: & &\\
			& & &\\
&\multicolumn{1}{l}{\quad\qquad(a) \cite{NV18a}:}   & $\QMA \subseteq \mathsf{MIP}^*[O(\log(n)), O(1)]$ &7\\
&    & (under randomized reductions) &\\[1ex]
&\multicolumn{1}{l}{\quad\qquad (b) \cite{NW19}:}   & $\NEEXP \subseteq \mathsf{MIP}^*$ &2\\[1ex]
&\multicolumn{1}{l}{\quad\qquad (c) \cite{JNV+20}:}   & $\MIP^* = \mathsf{RE}$&2
\end{tabular}
	\caption{Prior work on quantum-sound low degree tests and their complexity-theoretic consequences.
			The first three works showed a quantum-sound test and an $\MIP^*$ bound,
			both involving the same number of provers indicated in the final column.
			The last three works use the low-degree test from~\cite{NV18b}
			to show the indicated $\MIP^*$ bound.}
\label{fig:research}
\end{figure}
}

Subsequent to the initial posting of~\cite{JNV+20} on the arXiv,
an error was discovered in the analysis of the quantum-sound low degree test
contained in~\cite{Vid16} which was propagated to~\cite{NV18b}. The error affects the proof in a manner that appears difficult to fix. As such, we currently do not know if the low-degree test is quantum-sound for any number of provers.
The invalidation of this analysis affects every result in \Cref{fig:research} except for~\cite{IV12}.

The purpose of this exposition is to give a simplified soundness
analysis for the low individual degree test, with the quantitative
parameters stated explicitly below.
We do so by revisiting the three-player quantum-sound multilinearity test of~\cite{IV12}
and improving this result in two ways.
First, we generalize it to hold for the degree-$d$ low individual degree test,
of which the multilinearity test is the $d=1$ special case.
Second, using techniques introduced in~\cite{Vid16,NV18b},
we reduce the number of provers from~$3$ to~$2$.
Our main result is as follows.

\begin{theorem}[Main theorem, informal]\label{thm:main-informal}
Suppose Provers~$\mathrm{A}$ and~$\mathrm{B}$
pass the two-prover degree-$d$ low individual degree test with probability $1-\eps$
using a finite-dimensional projective strategy.
Let $A = \{A^u_a\}$ be the measurement the provers perform when they are given the point $u \in \F_q^m$
to produce a value $a \in \F_q$.
Then there exists a projective measurement $G = \{G_g\}$
whose outcomes~$g$ are polynomials of individual degree~$d$
such that
\begin{equation*}
\E_{\bu \sim \F_q^m}\sum_{a \in \F_q} \sum_{g:g(\bu) = a} \bra{\psi} A^{\bu}_{a} \ot G_g \ket{\psi}
\geq 1 - \poly(m,d) \cdot (\poly(\eps) + \poly(d/q)).
\end{equation*}
In other words, if Prover~$\mathrm{A}$ measures according to~$A^{\bu}$ to produce $\ba$ and Prover~$\mathrm{B}$ measures according to~$G$ to produce~$\bg$,
then $\bg(\bu) = \ba$ except with probability $\poly(m,d) \cdot (\poly(\eps) + \poly(d/q))$.
\end{theorem}
\noindent
Thus, we are able to extend \Cref{thm:classical-test-soundness}
to the case of two quantum provers (with an additional polynomial dependence on $d$; see \Cref{thm:main-formal} below for the formal statement of \Cref{thm:main-informal}).

The historical motivation for the original low individual degree
analysis was to recover $\NEEXP\subseteq\MIP^*$ from~\cite{NW19},
$\MIP^*=\mathsf{RE}$ from~\cite{JNV+20}, and the self-testing
applications of~\cite{NV18a}.  The simplified proof in this document
has an additional $\max\{1,d^2\}$ prefactor.  We state that dependence
explicitly and make no claim here that this quantitative bound
suffices for all of those applications.

\subsection{Total degree versus individual degree}

We now contrast the low degree test with the individual degree test
and explain why we are only able to prove the latter quantum sound.
We begin by explaining why the low individual degree test,
unlike the low degree test,
requires a $\poly(m) \cdot \poly(\eps)$ dependence in the soundness error.

\begin{example}\label{ex:bad-individual-degree-example}
Assume $d+1<q$, and consider the degree-$(d+1)$ polynomial
$h(x_1, \ldots, x_m) = x_1^{d+1}$,
and suppose that Players~$\mathrm{A}$ and~$\mathrm{B}$ play according to the following classical strategy.
\begin{itemize}
\item[$\circ$] (Player~$\mathrm{A}$): given $\bu \in \F_q^m$, return the value $\ba = h(\bu)$.
\item[$\circ$] (Player~$\mathrm{B}$): given the axis parallel line~$\bell = \{\bu + x \cdot e_{\bi} \mid x \in \F_q\}$, act as follows.
				If $\bi > 1$, then $h$ is a constant function along~$\bell$, and so return $\boldf = h|_{\bell}$.
				Otherwise, if $\bi = 1$, then $h$ is a degree-$(d+1)$ polynomial along~$\bell$.
				As the verifier expects a degree-$d$ polynomial,
				simply give up and return $\boldf \equiv 0$.
\end{itemize}
Using this strategy, $\boldf(\bu) = h|_{\bell}(\bu) = h(\bu) = \ba$ whenever $\bi \neq 1$, which occurs with probability $1- \frac{1}{m}$.
Hence, the players pass the degree-$d$ low individual degree test with probability at least $1 -\eps$ for $\eps = \frac{1}{m}$.
However, Player~$\mathrm{A}$ is responding to their questions using the polynomial~$h$
which is degree-$(d+1)$ but not degree-$d$,
and any individual-degree-$d$ polynomial~$g$ agrees with~$h$ on at
most a $\frac{d+1}{q}$ fraction of all inputs.  Indeed, after fixing
$x_2,\ldots,x_m$, the difference $h-g$ is a nonzero univariate
polynomial in $x_1$ of degree $d+1$, so Schwartz--Zippel applies.
This means that the agreement between Player~$\mathrm{A}$'s strategy and any individual-degree-$d$ polynomial~$g$ is at most
\begin{equation*}
\Pr_{\bu \sim \F_q^m}[g(\bu) = \ba]
\leq
1 - m \cdot \eps + \frac{d+1}{q}.
\end{equation*}
\end{example}

\Cref{ex:bad-individual-degree-example} shows why a dependence on~$m$
is necessary in the soundness error; it does not address the additional
$d$ dependence in \Cref{thm:main-informal}.
This reveals a weakness with the low individual degree test:
one can only conclude that the players are using a low individual degree strategy
when their failure probability $\eps$ is tiny---on the order of $\frac{1}{m}$ or smaller.
The low (total) degree test is alluring because it has the potential to avoid this dependence on~$m$.

Soundness of the low total and individual degree tests
is typically proven by induction on~$m$.
For the low individual degree test,
our simplified induction incurs errors polynomial in both $m$ and $d$.
The recurrence keeps the accumulated error of the form
$\poly(m,d)\cdot(\poly(\eps)+\poly(d/q))$,
exactly as in \Cref{thm:main-informal}.
For the low degree test,
on the other hand, each step of the induction only incurs an error of $\poly(\eps) + \poly(d/q)$.
However, if summed over all~$m$ steps,
this still gives a total error of $m \cdot (\poly(\eps) + \poly(d/q))$,
which is too large.

To account for this,
the soundness proofs in \cite{Vid16, NV18b}
introduce a technique at the end of each induction step
called \emph{Consolidation}
in which the growing error is ``reset''
down to an error $\poly(\eps) + \poly(d/q)$
which remains fixed across all levels of the induction.
This allows them to conclude with an error that was independent of the dimension~$m$.
Consolidation works as follows:
if the error of the projective measurement $G = \{G_g\}$
ever grows past this fixed error, the Consolidation step argues that on some portion of the Hilbert space,
the measurement~$G$ must be performing much worse than expected;
it then corrects this by inductively calling the low-degree test soundness to produce a better measurement on this portion of the Hilbert space.
Ultimately, however, this creates a cascade of Consolidation steps calling each other with increasing error,
which at some point grows so large that the low-degree soundness can no longer be applied.
This is the source of the bug.

Fortunately, this work shows that the techniques of these prior works,
aside from Consolidation,
are still sound.
So we can show soundness bounds which grow as a function of~$m$,
even if we cannot yet show bounds independent of~$m$.
In the end, this ``weakness" of the low individual degree test
is precisely what allows us to show that it, and not the low degree test, is quantum sound.

That said, we do believe, although we have not rigorously verified,
that related techniques could prove the original dimension-only
error form for the low degree test,
i.e.\ a soundness bound of $\poly(m) \cdot (\poly(\eps) + \poly(d/q))$ rather than $\eps + \poly(m) \cdot \poly(d/q)$.
This ``weak'' quantum soundness does not rule out the possibility that quantum provers can significantly outperform their classical counterparts.
It is, however, still sufficient for the same applications as the low individual degree test,
and it hints towards the possibility of a full quantum soundness of the low degree test.

\subsection{Conclusion and open problems}

In recent years,
the classical low-degree test has played a critical role in the study of nonlocal games
and the complexity class~$\MIP^*$.
In addition to correcting previous proofs of soundness,
we hope that this new exposition
will invite new researchers to engage with this beautiful area.
We conclude with a short list of open problems for future work.

\begin{enumerate}
\item Is the classical low-degree test quantum sound? Can the Consolidation step be fixed?
\item Even if the classical low-degree test is shown quantum sound,
	there are still interesting questions to be answered about the low individual degree test.
	For example, can the diagonal lines test be removed? (See \Cref{sec:the-test} for a description of this subtest.)
	Doing so would likely simplify the proof of $\MIP^* = \mathsf{RE}$~\cite{JNV+20},
	as one could replace the complicated ``conditional linear functions'' used in the proof
	with a simpler subclass known as ``coordinate deletion functions''.
	However, as discussed at the end of \Cref{sec:technical-overview},
	we know of an example that requires the diagonal lines test
	for the low individual degree test with parameters $m = 2$, $d = 2$, and $q = 4$.
	Can we find similar examples for larger~$q$?
\item Classically,  the low degree test
	generalized in various directions,
	including tests for affine invariant properties~\cite{KS08,Sud11}
	and tests for tensor product codes (see, for example, \cite{CMS17}).
	Does quantum soundness hold for these tests as well?
\item This work continues the trend of showing that well-studied classical property testers
	are also quantum sound.
	Prior to this, the work of~\cite{IV12} (see also \cite[Chapter 2]{Vid11}) showed that the linearity tester of Blum, Luby, and Rubinfeld~\cite{BLR93}
	is also quantum sound.
	The proofs of these results have so far been case-by-case adaptations of the classical proofs to the quantum setting;
	in the case of this paper, the proof is highly nontrivial
	and involves many ad hoc calculations which fortunately go in our favor.
	Is there a more conceptual reason why quantum soundness holds for these testers?
	Perhaps a reduction from the quantum case to the classical case?
\end{enumerate}

\section{Technical overview}\label{sec:technical-overview}

At a high level, our proof follows the approach of~\cite{BFL91}, and
it is useful to start by summarizing their analysis, which applies to
classical, \emph{deterministic} strategies. In this version of the test, the provers' strategy is described
by a ``points function,'' assigning a value in $\F_q$ to each point in
$\F_q^m$, and a
``lines function,'' assigning a low-degree polynomial to each
line in $\F_q^{m}$ queried in the test. From the assumption that the
points and lines functions agree at a randomly chosen point with high
probability, we would like to construct a global low-degree polynomial that has high agreement with
the points function. This is done by inductively constructing
``subspace functions'' defined on affine axis-aligned subspaces of increasing
dimension $k$. The base case is $k=1$, and is supplied by the lines
function. At each step, to construct the subspace function for a
subspace $S$ of dimension $k+1$, we pick $d+1$ parallel subspaces of
dimension $k$ that lie within $S$, and compute the unique degree-$d$
polynomial that interpolates them. The analysis shows that at each
stage, the function constructed through interpolation has high
agreement on average with the points function and lines function. In
the end, when we reach $k=m$ the ambient dimension, the subspace
function we construct is the desired global function.

For the sake of simplicity, it suffices to consider the case $d=1$,
i.e. multilinear functions. In this case, whenever we perform
interpolation, we need to combine $2$ parallel subspaces.

\paragraph{The classical zero-error case}

To start building intuition, it is useful to think about how to carry
out the above program in a highly simplified setting: the classical
zero-error case, for $m=3$. In this case, we assume that we have access to a
points function $f$ and lines function $g$ that \emph{perfectly} pass the BFL
test. Moreover, we will focus on the \emph{final} step of the
induction: thus, we assume that we have already constructed a set of
planes functions defined for every axis-parallel plane, that are
perfectly consistent with the line and point functions. In the final
step of the induction, our goal is to combine these planes functions
to create a single global function $h$
over all of $\F_q^{3}$ that is consistent with the points and lines functions.

To do this, we will interpolate the planes as follows. Let us label
the 3 coordinates in the space $x, y, $and $z$, and consider planes
parallel to the $(x,y)$-plane. Each plane $S_z$ is specified by a
value of the $z$ coordinate:
\[ S_z = \{(x,y,z): (x, y) \in \F_q^2\}. \]
By the induction hypothesis, for every such plane $S_z$ there exists a
bilinear function $g_{z}: \F_q^{2} \to \F_q$ that agrees with the
points function. To construct a global multilinear function $h:
\F_q^3 \to \F_q$, we pick two distinct values $z_1 \neq z_2$, and ``paste''
the two plane functions $g_{z_1} $ and
$g_{z_2}$ together using \emph{polynomial
  interpolation}. Specifically, we define $h$ to be the unique
multilinear polynomial that interpolates between $g_{z_1}$ on the
plane $S_{z_1}$ and $g_{z_2}$ on the plane $S_{z_2}$.
\[ h(x,y,z) = \mathrm{interpolate}_{z_1, z_2}(g_{z_1}, g_{z_2}) = \left(\frac{z - z_2}{z_1 - z_2}\right) g_{z_1}(x,y) +
  \left(\frac{z - z_1}{z_2 - z_1}\right) g_{z_2}(x,y). \]

This procedure defines a global function $h$, and by the assumption
that $g_{z_1}$ and $g_{z_2}$ are multilinear, it follows that $h$ is
also a multilinear function. But why is $h$ consistent with the points
function? To show this, we need to consider the lines function on
lines parallel to the $z$ axis. Given a point $(x,y,z)$, let $\ell$ be
the line parallel to the $z$ axis through this point, and let $g_\ell$
be the associated lines function. By construction, $h$ agrees with
$g_\ell$ at the two points $z_1$ and $z_2$. But $g_\ell$ and the restriction
$h_{|\ell}$ of $h$ to $\ell$ are both linear functions, and hence if
they agree at two points, they must agree \emph{everywhere}. (This is
a special case of the \emph{Schwartz-Zippel lemma}, which says that if
two degree-$d$ polynomials agree at $d+1$ points, they must be equal.)
Thus, $h$
agrees with $g_\ell$ at the original point $z$ as well. By success in
the test, $g_\ell$ in turn agrees with the points function $f$ at
$(x,y,z)$, and thus, $h(x,y,z) = f(x,y,z)$. Thus, we have shown that
the global function $h$ is both multilinear and agrees with the points
function $f$ exactly.

\paragraph{Dealing with errors}
To extend the sketch above to the general case, with nonzero error,
requires some modifications. At the most basic level, we may consider
what happens when we allow for deterministic classical strategies that
succeed with probability less than $1$ in the test. Such strategies
may have ``mislabeling'' error: the points function
$f$ may be imagined to be a multilinear function that has been
corrupted at a small fraction of the points. This type of error is
handled by the analysis in~\cite{BFL91}. The main modification to the zero-error
sketch above is a careful analysis of the probability that the pasting
step produces a ``good'' interpolated polynomial for a randomly chosen
pair of planes $S_{z_1}, S_{z_2}$. This analysis makes use of the
Schwartz-Zippel lemma together with combinatorial properties of the
point-line test itself (e.g. the expansion properties of the question
graph associated with the test).

At the next level of generality, we could consider classical
\emph{randomized} strategies. Suppose we are given a randomized
strategy that succeeds in the test with probability $1 - \eps$. Any
randomized strategy can be modeled by first sampling a random seed,
and then playing a deterministic strategy conditioned on the value of
the seed. A success probability of $1 - \eps$ could have two
qualitatively different underlying causes: (1) on $O(\eps)$ fraction
of the seeds, the strategy uses a function which is totally corrupted,
and (2) on a large fraction of the seeds, the strategy uses functions
which are only $\eps$-corrupted. An analysis of randomized
strategies could naturally proceed in a ``seed-by-seed'' fashion, applying the
deterministic analysis of~\cite{BFL91} to the large fraction of ``good'' seeds (which
are each only $\eps$-corrupted), while giving up entirely on the
``bad'' seeds.

In this document, we consider quantum strategies, which
have much richer possibilities for error. Nevertheless, we are able to
preserve some intuition from the randomized case, by working with
\emph{sub-measurements}: quantum measurements that do not always yield
an outcome. Working with a sub-measurement allows us to distinguish two
kinds of error: \emph{consistency error} (the probability that a
sub-measurement returns a wrong outcome) and \emph{completeness error}
(the probability that the sub-measurement fails to return an outcome at
all). Roughly speaking, the completeness error corresponds to the
probability of obtaining a ``bad'' seed in the randomized case, while
the consistency error corresponds to how well the strategies do on
``good'' seeds.

The technique of using sub-measurements and managing the two types of
error separately goes back to~\cite{IV12}. That work developed a
crucial tool to convert between these two types of error called the
\emph{self-improvement lemma}, and in our analysis we make extensive
use of a refined version of this lemma (\Cref{thm:self-improvement-in-induction-section}), building on~\cite{Vid16,NV18b}. Essentially, the lemma says the
following: suppose that (a) the provers pass the test with probability
$1 - \eps$, and (b) there is a complete measurement $G_{g}$ whose
outcomes are low-degree polynomials that has
consistency error $\nu$: that is, $G$ always returns an outcome, but
has probability $\nu$ of producing an outcome $g$ that disagrees with the points measurement $A^u_a$ at a
random point $u$. Then there exists an
``improved'' sub-measurement  $H_{h}$ with consistency error $\zeta$
depending on the test error and the parameters $m,d/q$, independently
of $\nu$, and with completeness error
(i.e. probability of not producing an outcome at all) of $\nu
+\zeta$. Essentially, this lemma says we can always ``reset'' the
consistency error of any measurement we construct at intermediate
points in the analysis to a universal
function $\zeta$ of the test error and the fixed test parameters,
at the cost of introducing some amount completeness error. Intuitively, one may
think of the action of the lemma as correcting $G$ on the portions of
Hilbert space where it is only mildly corrupted, while ``cutting out'' the
portions of Hilbert space where $G$ is too corrupted to be
correctable. In some sense, this lemma is the quantum analog of the
idea of identifying ``good'' and ``bad'' random seeds in the classical
randomized case. The proof of the lemma uses a primal--dual semidefinite program for measurements
together with the combinatorial facts used in~\cite{BFL91}
(specifically, expansion of the question graph of the test, and
the Schwartz-Zippel lemma).

Armed with the self-improvement lemma, we set up the following
quantum version of the BFL
induction loop: for $r$ running from $1$ to $m$, we construct a family
of subspace measurements $G^{S}_{g}$ that returns a low-degree
polynomial $g$ for every axis-aligned affine subspace $S$ of dimension $r$.
\begin{enumerate}
  \item By the induction hypothesis, we know that there exists a
    measurement  $G^{S}_{g}$ for every $r$-dimensional subspace, which
    has consistency error $\delta(r)$ with the points
    measurement. (For the base case $r = 1$, this is the lines
    measurement from the provers' strategy).
  \item We apply the self-improvement lemma to these measurements, yielding
    sub-measurements $\hat{G}^{S}_{g}$ that have consistency error
    $\zeta$ independent of $\delta(r)$, and completeness error $\kappa(r) = \delta(r) + \zeta$.
    We dilate all of these sub-measurements on a common enlarged space
    to obtain projective slices with the same correlations.
  \item For each subspace $S$ of dimension $r+1$, we construct a
    pasted sub-measurement, by performing a quantum version of the
    classical interpolation argument: we define the pasted
    sub-measurement by completing the measurements with a rejection
    outcome, sampling parallel $r$-dimensional subspaces, and
    interpolating only the first $d+1$ successful outcomes at distinct
    slice coordinates. This pasted sub-measurement has consistency slightly worse than $\zeta$, and completeness error which
    is slightly worse than $\kappa(r)$. It is at this step that it is crucial to treat the two types
    of error separately: in particular, we need the consistency error
    to be low to ensure that the interpolation produces a good result.
  \item We convert the resulting sub-measurement into a full
    measurement, by assigning an arbitrary fixed outcome whenever the
    sub-measurement fails to yield an outcome. This measurement will
    have consistency error $\delta(r+1)$ which is larger than $\delta(r)$
    by some additive factor.  Compressing this measurement back to the
    original space preserves its consistency with the point measurements.
  \end{enumerate}
  At the end of the loop, when $r = m$, we obtain a single measurement that returns
  a global polynomial as desired.

The pasting analysis uses one positive contraction, the average of the
projections onto matching outcomes of the completed slice measurements.
Its powers control independent random sequences of questions.  A
spectral inequality then bounds both the cost of moving a successful
slice past an entire prefix and the change in marginal success
probabilities.  A union bound proves consistency of the interpolant,
and the expected number of successes controls completeness.  No
independence between the outcomes is needed.  In \Cref{sec:first-construction}, boundedness turns the same
estimates into a bound on $\bra\psi(G-G^2)\otimes I\ket\psi$ and proves
that the original construction with exactly $d+1$ slices also works.
The induction uses the first-success construction for its sharper
quantitative bound.
\paragraph{The diagonal lines test}
An important element of the test we analyze in this document, which is
unnecessary in the classical case, is the diagonal lines test. The purpose
of this test is to certify that the points measurements used by the provers
approximately commute on average over all pairs of points $(x, y)$ in
$\F_q^m$---something which is automatically true in the classical case. This is done by asking one prover for a polynomial defined
on the line going through $x$ and $y$, and the other for the function
evaluation at either $x$ or $y$. The line through $x$ and
$y$ will not in general be axis-parallel; we refer to these general
lines as ``diagonal.''

The commutation guarantee plays an important role in our analysis of
the test, and it is an interesting question whether it is truly
necessary to test it directly with the diagonal lines test; might it
not automatically follow from success in the axis-parallel lines test? An interesting contrast can be drawn to
the Magic Square
game~\cite{mermin1990simple,peres1990incompatible,aravind2002simple},
in which questions are either cells or axis-parallel lines in a $3 \times 3$ square grid.
This has the same question distribution as
the axis-parallel line-point test over $\F_3^2$ (although the answers in the Magic Square game are strings in~$\F_2$ rather than~$\F_3$).
For the Magic Square game,
``points'' measurements along the same axis-parallel ``line'' commute,
but points that are not axis-aligned do \emph{not} commute: indeed,
for the perfect strategy, they anticommute. Despite this example,
we know that commutation between all pairs of points \emph{can} be deduced from the
axis-parallel lines test alone, rendering the diagonal lines test unnecessary,
at least in the case of the bivariate ($m = 2$) multilinearity test (the low individual degree test when $d = 1$).
On the other hand, we know of a quantum strategy using noncommuting measurements
which succeeds with probability~$1$ in the $m = 2$, $d = 2$, $q = 4$ low individual degree test.
Whether this counterexample can be extended to larger~$q$ remains an open question.

\paragraph{Organization}
The rest of this document is organized as follows. In
\Cref{sec:prelims} we review some preliminaries concerning finite
fields, polynomials, and quantum measurements. In
\Cref{sec:making-measurements-projective}, we prove the rounding
lemma used once at the end to obtain projective output on the original
local spaces. In \Cref{sec:induction}, we give the inductive argument
proving the main theorem. In \Cref{sec:expansion,sec:variance} we show
some properties of the hypercube graph and families of measurements
indexed by points on the hypercube, which are used in
\Cref{sec:self-improvement} to prove the self-improvement lemma. In
\Cref{sec:commutativity-points,sec:g-comm}, we show commutativity
properties of the measurements constructed in the induction, and
finally in \Cref{sec:ld-pasting} we analyze the pasting step of the induction.

\paragraph{Acknowledgments}

We thank Lewis Bowen for pointing out typos and a few minor errors in a previous version.
We also thank Madhu Sudan for help with references on the classical low individual degree test.

\section{The test}\label{sec:the-test}

\begin{definition}[Roles]
The low individual degree test will be played between two provers and a verifier.
The two provers are named Player~$\mathrm{A}$ and Player~$\mathrm{B}$.
A \emph{role} $r$ is an element of the set $\{\mathrm{A}, \mathrm{B}\}$.
Given a role~$r$, we write $\overline{r}$ for the other element of the set $\{\mathrm{A}, \mathrm{B}\}$.
\end{definition}

\begin{definition}[Low individual degree test]
Let $m\geq1$ and $d\geq0$ be integers.
Let $q$ be a prime power.
Then the \emph{$(m,q,d)$-low individual degree test}
is stated in \Cref{fig:test}.
\end{definition}

{
\floatstyle{boxed}
\restylefloat{figure}
\begin{figure}
With probability~$\tfrac{1}{3}$ each, perform one of the following three tests.
\begin{enumerate}
	\item \textbf{Axis-parallel lines test:}
		Pick a uniformly random role $\br \sim \{\mathrm{A},\mathrm{B}\}$.
		Let $\bu \sim \F_q^m$ be a uniformly random point.
		Select $\bi \sim \{1, \ldots, m\}$ uniformly at random,
		and let $\bell = \{\bu + t \cdot e_{\bi} \mid t \in \F_q\}$
		be the axis-parallel line which passes through~$\bu$ in the $\bi$-th direction.
		\begin{itemize}
		\item[$\circ$] Player~$\br$:
			Give $\bell$;
			receive the univariate degree-$d$ polynomial $\boldf:\bell \rightarrow \F_q$.
		\item[$\circ$] Player~$\overline{\br}$:
			Give $\bu$;
			receive $\ba \in \F_q$.
		\end{itemize}
		Accept if $\boldf(\bu) = \ba$.
	\item \textbf{Self-consistency test:}
		Let $\bu \sim \F_q^m$ be a uniformly random point.
		\begin{itemize}
		\item[$\circ$] Player~$\mathrm{A}$:
			Give $\bu$;
			receive $\ba \in \F_q$.
		\item[$\circ$] Player~$\mathrm{B}$:
			Give $\bu$;
			receive $\bb \in \F_q$.
		\end{itemize}
		Accept if $\ba = \bb$.
	\item \textbf{Diagonal lines test:}
		Pick a uniformly random role $\br \sim \{\mathrm{A},\mathrm{B}\}$.
		Let $\bu \sim \F_q^m$ be a uniformly random point.
		Select $\bi \sim \{1, \ldots, m\}$ uniformly at random,
		and let $\bv \in \F_q^m$ be a uniformly random point whose last $m - \bi$ coordinates are~$0$.
		Finally,
		let $\bell = \{\bu + t \cdot \bv \mid t \in \F_q\}$
		be the line which passes through~$\bu$ in direction~$\bv$.
		When $\bv=0$, this is the singleton $\{\bu\}$; throughout the
		paper such a singleton is regarded as a permitted degenerate
		diagonal line.
		\begin{itemize}
		\item[$\circ$] Player~$\br$:
			Give $\bell$;
			receive the univariate degree-$md$ polynomial $\boldf:\bell \rightarrow \F_q$.
		\item[$\circ$] Player~$\overline{\br}$:
			Give $\bu$;
			receive $\ba \in \F_q$.
		\end{itemize}
		Accept if $\boldf(\bu) = \ba$.
	\end{enumerate}
	\caption{The $(m,q,d)$-low individual degree test.\label{fig:test}}
\end{figure}
}

An important class of strategies are those which are \emph{symmetric}.
In this case, the bipartite state $\ket{\psi}$ Alice and Bob share is symmetric.
Furthermore, for any question Alice and Bob receive, they apply the same measurement to their share of the state.
This means that rather than, for example, keeping track of a separate points measurement for Alice and Bob,
we can use a single measurement to refer to both of their strategies,
and similarly for the axis-parallel lines and diagonal lines measurements.
This is formalized in the following definition, where we also consider strategies which are projective.

We make one harmless relaxation that is important for the induction.
For the diagonal-lines measurement $L^\ell$, we allow every function
$f:\ell\to\F_q$ as an outcome, rather than only degree-$md$
polynomials.  Given an ordinary strategy, first coarse-grain together
all polynomial labels that induce the same function on~$\ell$, and set
$L^\ell_f=0$ for every remaining function.  Coarse-graining a
projective measurement is projective, and the acceptance probability
is unchanged.  Conversely, the proof below uses $L$ only through
projectivity and the coarse-grained evaluations
$L^\ell_{[f(u)=a]}$, so it proves soundness even for this larger answer
alphabet.  This relaxation is closed under restricting to an affine
slice; the degree-$md$ answer alphabet alone is not.

\begin{definition}[Symmetric, projective strategy]
A \emph{symmetric, projective strategy} for the $(m,q, d)$-low individual degree test
is a tuple $(\psi, A, B, L)$ defined as follows.
\begin{itemize}
\item[$\circ$] $\calH$ is finite-dimensional and
$\ket{\psi} \in \calH \ot \calH$ is a bipartite, permutation-invariant state.
\item[$\circ$] $A = \{A^u_a\}$ contains a matrix for each $u \in \F_q^m$ and $a \in \F_q$.
			For each $u$, $A^u$ is a projective measurement on~$\calH$.
\item[$\circ$] $B = \{B^{\ell}_f\}$ contains a matrix for each axis-parallel line~$\ell$ in $\F_q^m$
			and univariate degree-$d$ polynomial $f:\ell\rightarrow \F_q$.
			For each $\ell$, $B^{\ell}$ is a projective measurement on~$\calH$.
\item[$\circ$] $L = \{L^{\ell}_f\}$ contains a matrix for each line~$\ell$ in $\F_q^m$
			and each function $f:\ell\rightarrow \F_q$.
			For each $\ell$, $L^{\ell}$ is a projective measurement on~$\calH$.
\end{itemize}
\end{definition}

We can also consider more general strategies which are no longer assumed to be symmetric.
In this case, Players~$\mathrm{A}$ and~$\mathrm{B}$ each have their own versions of the measurements~$A$, $B$, and~$L$.

\begin{definition}[General projective strategy]
A \emph{projective strategy} for the $(m,q, d)$-low individual degree test
is a tuple $(\psi, A^{\mathrm{A}},B^{\mathrm{A}},L^{\mathrm{A}},A^{\mathrm{B}},B^{\mathrm{B}},L^{\mathrm{B}})$ defined as follows.
\begin{itemize}
\item[$\circ$] $\calH_{\mathrm A}$ and $\calH_{\mathrm B}$ are
finite-dimensional, and
$\ket{\psi} \in \calH_{\mathrm{A}} \ot \calH_{\mathrm{B}}$ is a bipartite state.
\end{itemize}
Furthermore, for each~$w \in \{\mathrm{A}, \mathrm{B}\}$:
\begin{itemize}
\item[$\circ$] $A^{w} = \{A^{w,u}_a\}$ contains a matrix for each $u \in \F_q^m$ and $a \in \F_q$.
			For each $u$, $A^{w,u}$ is a projective measurement on~$\calH_{w}$.
\item[$\circ$] $B^{w} = \{B^{w,\ell}_f\}$ contains a matrix for each axis-parallel line~$\ell$ in $\F_q^m$
			and univariate degree-$d$ polynomial $f:\ell\rightarrow \F_q$.
			For each $\ell$, $B^{w,\ell}$ is a projective measurement on~$\calH_{w}$.
\item[$\circ$] $L^w = \{L^{w,\ell}_f\}$ contains a matrix for each line~$\ell$ in $\F_q^m$
			and each function $f:\ell\rightarrow \F_q$.
			For each $\ell$, $L^{w,\ell}$ is a projective measurement on~$\calH_{w}$.
\end{itemize}
\end{definition}

\begin{remark}
Throughout this work, we will only consider projective strategies.
So we will henceforth use the terms ``strategy" and ``symmetric strategy'' to refer exclusively to projective strategies and symmetric, projective strategies, respectively.

In addition, we will spend the vast majority of this work dealing solely with symmetric strategies, as they are notationally simpler to work with.
Much of this work will focus on proving~\Cref{thm:main-induction},
a variant of our main theorem for symmetric strategies.
Our main theorem for general strategies, \Cref{thm:main-formal} below, will be proven in \Cref{sec:induction} by a standard reduction to the symmetric case.
We will only see these more cumbersome-to-write general strategies
in this section, where we state \Cref{thm:main-formal},
and in \Cref{sec:induction}, where we carry out the reduction.
\end{remark}

\begin{definition}[Good strategy]
A strategy is $(\eps, \delta, \gamma)$-good if it
passes the axis-parallel lines test with probability at least $1-\eps$,
the self consistency test with probability at least~$1-\delta$,
and the diagonal lines test with probability at least~$1-\gamma$.
\end{definition}

\begin{remark}\label{rem:good-strat-characterization}
Using notation which will be introduced in \Cref{sec:comparing-measurements} below,
a symmetric strategy is $(\eps, \delta, \gamma)$-good
if and only if it satisfies the following three conditions.
For $\bell$ and $\bu$ as in the axis-parallel lines test,
\begin{equation*}
A^u_a \ot I \simeq_{\eps} I \ot B^{\ell}_{[f(u)=a]},
\quad
\text{and}
\quad
A^u_a \ot I \simeq_{\delta} I \ot A^u_a.
\end{equation*}
And for $\bell$ and $\bu$ as in the diagonal lines test,
\begin{equation*}
A^u_a \ot I \simeq_{\gamma} I \ot L^{\ell}_{[f(u)=a]}.
\end{equation*}
\end{remark}

\begin{notation}\label{not:conditioned-on-last-direction}
Our proof will be via induction,
i.e.\ proving soundness of the $(m+1,q,d)$-low individual degree test
using the soundness of the $(m,q,d)$-low individual degree test.
To do this, we will frequently use the axis-parallel line test in the specific case of $i = m+1$.
Thus, it will be convenient to introduce the following notation.
Let $(\psi, A, B, L)$ be a symmetric strategy for the $(m+1, q, d)$-low individual degree test,
Then for each $u \in \F_q^m$ we will write $B^u_f$ as shorthand for $B^{\ell}_f$,
where $\ell = \{(u, x) \mid x \in \F_q\}$.
For a function $f:\ell \rightarrow \F_q$,
we will also sometimes write $f(x)$ as shorthand for $f(u,x)$.
\end{notation}

\begin{definition}
Consider the $(m,q,d)$-low individual degree test.
For $j \in \{1, \ldots, m\}$,
we refer to the \emph{$j$-restricted diagonal lines test}
as the diagonal lines test conditioned on $\bi = j$.
For example, in the $m$-restricted diagonal lines test, one samples
$u,v\in\F_q^m$ independently and uniformly and asks the line
$\ell=\{u+tv:t\in\F_q\}$.  Thus the singleton $\{u\}$ has
probability $q^{-m}$, while each nondegenerate line through~$u$ has
probability $(q-1)q^{-m}$.
\end{definition}

Now we state our main theorem using notation which will be introduced in \Cref{sec:prelims} below.
This is the formal version of \Cref{thm:main-informal}.

\begin{theorem}[Main theorem; quantum soundness of the low individual degree test]\label{thm:main-formal}
  Let $m\geq1$.  Consider a projective strategy $(\psi, A^{\mathrm{A}},B^{\mathrm{A}},L^{\mathrm{A}},A^{\mathrm{B}},B^{\mathrm{B}},L^{\mathrm{B}})$
   which passes the  $(m,q,d)$-low individual degree test with probability at least $1-\eps$.
  Set
  \[
  K_{m,d}=
  \begin{cases}
    m^4,&d=0,\\
    d^2m^4,&d\geq1,
  \end{cases}
  \qquad
  \nu = 21000 K_{m,d}
  \Big(\eps^{1/64} + (d/q)^{1/64}\Big).
  \]
Then there exist projective measurements $G^{\mathrm{A}}, G^{\mathrm{B}} \in \polymeas{m}{q}{d}$ with the following properties:
\begin{enumerate}
  \item (Consistency with~$A$):  On average over $\bu \sim \F_q^{m}$,
	  \begin{align*}
	  A^{\mathrm{A},u}_a \otimes I
		&\simeq_{\nu} I \otimes G^{\mathrm{B}}_{[g(u)=a]},\\
	I \otimes A^{\mathrm{B},u}_a
		&\simeq_{\nu}  G^{\mathrm{A}}_{[g(u)=a]} \ot I.
	  \end{align*}
\item (Self-consistency):
	\begin{equation*}
	G^{\mathrm{A}}_g \ot I \simeq_{\nu} I \ot G^{\mathrm{B}}_g.
	\end{equation*}
\end{enumerate}
\end{theorem}

\section{Preliminaries}\label{sec:prelims}

We use \textbf{boldface font} to denote random variables.
For two complex numbers $\alpha, \beta \in \C$, we write
$\alpha \approx_{\eps} \beta$ if
\begin{equation*}
|\alpha - \beta| \leq \eps.
\end{equation*}
We note the following triangle inequality for numbers, which we will use repeatedly:
\begin{equation}\label{eq:triangle-inequality-for-numbers}
\text{if $\alpha \approx_{\eps} \beta$ and $\beta \approx_{\delta} \gamma$, then $\alpha \approx_{\eps+\delta}\gamma$.}
\end{equation}

\subsection{Finite fields}

A \emph{finite field} is a field with a finite number of elements.
There is a unique finite field~$\F_q$ of~$q$ elements for each prime power $q = p^t$,
and there are no other finite fields.

\subsection{Polynomials over finite fields}

\begin{definition}[Polynomials of low individual degree]
Let $q$ be a prime power,
and let $m$ and $d$ be nonnegative integers.
We define $\polyfunc{m}{q}{d}$ to be the set of polynomials in $\F_q^m$
with individual degree~$d$.
\end{definition}

\begin{remark}
Note that as defined in \Cref{sec:intro},
in this work a polynomial with individual degree~$d$ is one in which the degree~$d_i$ of each coordinate~$i$
is \emph{at most~$d$}.
This allows us to say ``individual degree~$d$'' rather than the wordier ``individual degree at most~$d$''.
For example, under this definition, $\polyfunc{m}{q}{d}$ is contained in $\polyfunc{m}{q}{d+1}$ for each~$d$.
\end{remark}

The most important fact about low-degree polynomials is that they have large distance from each other.
This is shown by the following lemma.

\begin{lemma}[Schwartz-Zippel lemma~\cite{Sch80,Zip79}]\label{lem:schwartz-zippel-total-degree}
Let $g,h\in\F_q[X_1,\ldots,X_m]$ be formal polynomials such that
$g-h$ is nonzero and has total degree at most~$d$.
Then
\begin{equation*}
\Pr_{\bx \sim \F_q^m}[g(\bx) = h(\bx)] \leq \frac{d}{q}.
\end{equation*}
\end{lemma}

Since any polynomial with individual degree~$d$ has total degree~$md$,
\Cref{lem:schwartz-zippel-total-degree} implies the following corollary.

\begin{corollary}[Schwartz-Zippel for individual degree]
Let $g, h \in \polyfunc{m}{q}{d}$ be distinct. Then
\begin{equation*}
\Pr_{\bx \sim \F_q^m}[g(\bx) = h(\bx)] \leq \frac{md}{q}.
\end{equation*}
\end{corollary}
\begin{proof}
Apply \Cref{lem:schwartz-zippel-total-degree} to the nonzero
polynomial $g-h$, whose total degree is at most~$md$.
\end{proof}

\subsection{Measurements}

All measurement outcome sets in this paper are finite.

\begin{definition}[Measurements and sub-measurements]
Let $\calH$ be a Hilbert space and $\calA$ be a set of outcomes.
A \emph{sub-measurement}
is a set of Hermitian, positive-semidefinite matrices
$A = \{A_a\}_{a \in \calA}$
acting on~$\calH$
such that $\sum_a A_a \leq I$.
The sub-measurement is \emph{projective} if $(A_a)^2 = A_a$ for each~$a$.
It is a \emph{measurement}
if it satisfies the stronger condition $\sum_a A_a = I$.
\end{definition}

\begin{proposition}[Orthogonality of projective sub-measurements]
\label{prop:projective-submeasurement-orthogonal}
If $P=\{P_a\}$ is a projective sub-measurement, then
$P_aP_b=0$ whenever $a\ne b$.  Consequently
$P=\sum_aP_a$ is itself a projection.
\end{proposition}
\begin{proof}
Fix~$a$ and put $S=\sum_bP_b\leq I$.  Since $I-S\geq0$,
\[
0\leq P_a(I-S)P_a=-\sum_{b\ne a}P_aP_bP_a.
\]
Each summand $P_aP_bP_a=(P_bP_a)^\dagger(P_bP_a)$ is positive
semidefinite.  They must therefore all vanish, so $P_bP_a=0$ for
$b\ne a$; taking adjoints gives $P_aP_b=0$.  Hence $P^2=P$.
\end{proof}

An important class of sub-measurements
are those that output polynomials~$g$ of individual degree~$d$.
These are defined as follows.

\begin{definition}[Low-degree polynomial measurements]
We write $\polysub{m}{q}{d}$ for the set of sub-measurements $G=\{G_g\}$
with outcomes $g \in \polyfunc{m}{q}{d}$.
We write $\polymeas{m}{q}{d}$ for the subset of $\polysub{m}{q}{d}$
containing only those $G = \{G_g\}$ which are measurements.
\end{definition}

\begin{definition}[Post-processing measurements]\label{def:post-processing}
Let $A = \{A_a\}_{a \in \calA}$ be a set of matrices,
and let $f:\calA \rightarrow \calB$ be a function.
Then for each $b \in \calB$,
we define the matrix
\begin{equation*}
A_{[f(a)=b]} = \sum_{a:f(a)=b} A_a.
\end{equation*}
\end{definition}

\begin{remark}
We note that \Cref{def:post-processing}
agrees with the notation for post-processing measurements
used in \cite{NW19},
but it disagrees with the notation used in \cite{JNV+20}.
That work uses the notation ``$A_{[f(\cdot)=b]}$'' rather
than the notation ``$A_{[f(a)=b]}$'' that we use in this work.
\end{remark}

An easy-to-prove fact is that measurements remain measurements after post-processing their outcomes.

\begin{proposition}
Let $A = \{A_a\}_{a \in \calA}$ be a set of matrices,
and let $f:\calA \rightarrow \calB$ be a function.
Then
\begin{equation*}
\sum_a A_a = \sum_b A_{[f(a)=b]}.
\end{equation*}
Thus, if~$\{A_a\}$ is a sub-measurement (respectively, measurement),
then~$\{A_{[f(a)=b]}\}$ is also a sub-measurement (respectively, measurement).
\end{proposition}
\begin{proof}
The fibers of~$f$ partition~$\calA$, so
$\sum_bA_{[f(a)=b]}=\sum_b\sum_{a:f(a)=b}A_a=\sum_aA_a$.
\end{proof}

\begin{remark}
We note that there is some ambiguity in the notation $A_{[f(a)=b]}$
because it requires one to know which of~$f$ or~$a$ is the measurement outcome
and which is the function applied to it.
For example, given a sub-measurement $G = \{G_g\} \in \polysub{m}{q}{d}$,
we will often consider evaluating its outputs at a point $u \in \F_q^m$.
This entails looking at the sub-measurement
\begin{equation*}
\{G_{[g(u)=a]}\}_{a \in \F_q},
\quad
\text{where}
\quad
G_{[g(u)=a]} = \sum_{g:g(u) = a} G_g.
\end{equation*}
Here, $g$ is the measurement outcome of~$G$,
and the function being applied maps it to its evaluation on the point~$u$, i.e.\ $g(u)$.
In general, it should always be clear from context
what the measurement outcome and the function being applied to it are.
\end{remark}

\begin{notation}[Measurements indexed by questions]
We will frequently encounter sets of sub-measurements $A^x= \{A^x_{a}\}$ indexed by elements~$x$ from a set of ``questions" $\calX$. We will write $A = \{A^x_a\}$ for this set.
We will typically refer to this set~$A$ as a sub-measurement,
and we will refer to it as a measurement if each~$A^x$ is a measurement.
In addition, we will refer to it as projective if each~$A^x$ is projective.
\end{notation}

\begin{notation}[Complete part of sub-measurement]
Given a sub-measurement~$A = \{A_a\}$, we will write $A = \sum_a A_a$.
Similarly, given a sub-measurement~$A = \{A^x_a\}$, we will write $A^x = \sum_a A^x_a$ and $A = \E_{\bx} A^{\bx}$.
\end{notation}

The complete part of a sub-measurement
contrasts with its \emph{incomplete part},
which is the matrix $I - A^x$.
We will sometimes
view~$A$ as a measurement
by throwing in its incomplete part as an additional POVM element.
This is formalized in the following definition.

\begin{definition}[Completing a sub-measurement]\label{def:measurement-completion}
Let $A = \{A^x_a\}_{a \in \calA}$ be a sub-measurement.
We define the \emph{completion of~$A$},
denoted $\mathrm{completion}(A)$,
to be the measurement $\widehat{A} = \{\widehat{A}^x_a\}$
with outcome set $\widehat{\calA} = \calA \cup \{\bot\}$
such that for each $a \in \widehat{\calA}$,
\begin{equation*}
\widehat{A}^x_a
= \left\{\begin{array}{cl}
		A^x_a & \text{if $a \in \calA$,}\\
		I - A^x & \text{if $a = \bot$}.
		\end{array}\right.
\end{equation*}
\end{definition}

\subsection{Comparing measurements}
\label{sec:comparing-measurements}

A central problem in this paper is recognizing when two measurements are close to each other.
We will survey two methods of doing so,
using the \emph{consistency} and the \emph{state dependent distance}.
This section largely mirrors Sections 4.4 and 4.5 of \cite{NW19},
which prove numerous properties of these two distances.
However, we are unable to cite their results directly
because they are mostly stated and proven only for measurements,
whereas we will need to apply them to sub-measurements as well.

\subsubsection{Consistency between measurements}

The most basic notion of similarity between two measurements is given by their \emph{consistency}.

\begin{definition}[Consistency]\label{def:simeq}
Let $\ket{\psi}$ be a state in $\calH_{\mathrm{A}} \ot \calH_{\mathrm{B}}$.
Let $A = \{A^x_a\}$ be a sub-measurement acting on $\calH_{\mathrm{A}}$
and $B = \{B^x_a\}$ be a sub-measurement acting on $\calH_{\mathrm{B}}$.
Finally, let $\calD$ be a distribution on the question set~$\calX$.
Then we say that
 \begin{equation*}
 A^x_a \ot I \simeq_{\delta} I \ot B^x_a
 \end{equation*}
  on state $\ket{\psi}$ and distribution $\calD$ if
  \begin{equation}\label{eq:no-big-Oh}
  \E_{\bx \sim \calD} \sum_{a \neq b} \bra{\psi} A^{\bx}_a \ot B^{\bx}_b
    \ket{\psi} \leq \delta.
    \end{equation}
\end{definition}

This is the joint probability that both sub-measurements return an
outcome and the two outcomes differ; it is not conditioned on success.

\begin{notation}[Simplifying notation]
Because the state $\ket{\psi}$ and distribution $\calD$ are typically clear from context,
we will often write ``$ A^x_a \ot I \simeq_{\delta} I \ot B^x_a$"
as shorthand for ``$ A^x_a \ot I \simeq_{\delta} I \ot B^x_a$ on state $\ket{\psi}$ and distribution $\calD$".
In the case when $\calD$ is not clear by context,
we might specify it implicitly in terms of a random variable~$\bx$ distributed according to~$\calD$.
For example, suppose the distribution~$\calD$ is supposed to be uniform on the question set $\F_q$.
Then we might say ``on average over $\bx \sim \F_q$,
\begin{equation*}
A^x_a \ot I \simeq_{\delta} I \ot B^x_a"
\end{equation*}
as shorthand for ``$ A^x_a \ot I \simeq_{\delta} I \ot B^x_a$ on distribution $\calD$".
\end{notation}

We note that there are two differences between the definition of consistency in \Cref{def:simeq}
and the original definition of consistency given in \cite[Definition 4.11]{NW19}.
The first of these is that the~\cite{NW19} definition
allows the right-hand side of \Cref{eq:no-big-Oh} to be $O(\delta)$
rather than strictly~$\delta$.
The benefit of this is that they do not need to keep track of constant prefactors in their proofs;
we have elected to use this more concrete definition
to make our proofs more easily verifiable, at the expense of tracking these constant prefactors.
The second difference is that they define their consistency as the quantity
\begin{equation*}
1 - \E_{\bx} \sum_a \bra{\psi} A^{\bx}_a \ot B^{\bx}_a \ket{\psi}.
\end{equation*}
As we show below in \Cref{prop:simeq-for-measurements},
this agrees with \Cref{def:simeq}
when~$A$ and~$B$ are both measurements.
However, these two definitions disagree
when~$A$ and~$B$ are \emph{sub}-measurements,
which is a case we will frequently encounter throughout this paper.

\begin{proposition}[Consistency for measurements]\label{prop:simeq-for-measurements}
Let $A = \{A^x_a\}$ and $B = \{B^x_a\}$ be two measurements.
Then
\begin{equation*}
A^x_a \ot I \simeq_{\delta} I \ot B^x_a
\end{equation*}
if and only if
\begin{equation*}
\E_{\bx} \sum_a \bra{\psi} A^{\bx}_a \ot B^{\bx}_a \ket{\psi} \geq 1 - \delta.
\end{equation*}
\end{proposition}
\begin{proof}
Completeness gives the identity
\[
 \E_x\sum_{a\ne b}\bra\psi A_a^x\otimes B_b^x\ket\psi
 =1-\E_x\sum_a\bra\psi A_a^x\otimes B_a^x\ket\psi.
\]
The two stated inequalities are therefore equivalent.
\end{proof}

\subsubsection{The state-dependent distance}

To transfer consistency between nearby measurements, we use the
state-dependent distance.  It applies to arbitrary operator families,
including products that need not be measurements or Hermitian.

\begin{definition}[The state-dependent distance]\label{def:approx_delta}
Let $\ket{\psi}$ be a state in $\calH$.
Let $A = \{A^x_a\}$ and $B = \{B^x_a\}$ be
sets of matrices acting on $\calH$.
Finally, let $\calD$ be a distribution on the question set~$\calX$.
Then we say that
  \begin{equation*}
  A^x_a \approx_{\delta}  B^x_a
  \end{equation*}
 on state $\ket{\psi}$ and distribution $\calD$ if
  \begin{equation*}
  \E_{\bx \sim \calD} \sum_{a} \Vert (A^{\bx}_a - B^{\bx}_a) \ket{\psi} \Vert^2 \leq \delta.
    \end{equation*}
\end{definition}

\begin{lemma}[Direct-sum estimates]\label{lem:direct-sum-estimates}
For vector families indexed by a question $x$ and an outcome $a$, write
\[
 \langle v,w\rangle_\oplus=\E_x\sum_a\langle v_a^x,w_a^x\rangle,
 \qquad \|v\|_\oplus^2=\langle v,v\rangle_\oplus.
\]
Cauchy--Schwarz and the triangle inequality hold for this norm.  If
$\sum_b(C_{a,b}^x)^\dagger C_{a,b}^x\leq I$ for every $x,a$, then
\[
 \|(C_{a,b}^xv_a^x)_{x,a,b}\|_\oplus\leq\|v\|_\oplus.
\]
For any contraction $M$ on the direct sum,
\begin{equation}\label{eq:direct-sum-quadratic-form}
 |\langle v,Mv\rangle_\oplus-\langle w,Mw\rangle_\oplus|
 \leq(\|v\|_\oplus+\|w\|_\oplus)\|v-w\|_\oplus.
\end{equation}
\end{lemma}
\begin{proof}
These are Hilbert-space inequalities after weighting each question
component by the square root of its probability.  For the contraction
family, sum $\langle v_a^x,\sum_b(C_{a,b}^x)^\dagger C_{a,b}^xv_a^x\rangle$.
For the last estimate, expand the difference as
$\langle v-w,Mv\rangle_\oplus+\langle w,M(v-w)\rangle_\oplus$.
\end{proof}

All direct-sum norms below include the displayed question averages;
we omit the subscript $\oplus$ when the indexing is clear.

\begin{proposition}[Swap invariance of the state-dependent distance]
\label{prop:approx-swap-invariance}
If $\ket\psi\in\calH\otimes\calH$ is permutation invariant, then
for any two families of matrices,
\[
X^x_a\otimes I\approx_\eta Y^x_a\otimes I
\quad\Longleftrightarrow\quad
I\otimes X^x_a\approx_\eta I\otimes Y^x_a.
\]
\end{proposition}
\begin{proof}
Apply the tensor-swap unitary, which fixes~$\ket\psi$, to every
vector in the squared-distance sum.
\end{proof}

The next proposition transfers consistency even when the other
family is only a sub-measurement.

\begin{proposition}[Transferring ``$\simeq$" using ``$\approx$"]
Let $\{A^x_a\}$ and $\{B^x_a\}$ be measurements,
and let $\{C^x_a\}$ be a sub-measurement.
Suppose that $A^x_a \ot I \simeq_{\delta} I \ot C^x_a$
and $A^x_a \ot I \approx_{\eps} B^x_a \ot I$.
Then $B^x_a \ot I \simeq_{\delta + \sqrt{\eps}} I \ot C^x_a$.
  \label{prop:triangle-sub}
\end{proposition}

\begin{proof}
Put $C^x=\sum_aC_a^x$.  Since $A$ and $B$ are complete, their
inconsistencies with $C$ are respectively
\[
 \E_x\bra\psi I\otimes C^x\ket\psi
 -\E_x\sum_a\bra\psi A_a^x\otimes C_a^x\ket\psi
\]
and the same expression with $A$ replaced by $B$.  Their difference
has modulus at most
\[
 \left|\E_x\sum_a\bra\psi(A_a^x-B_a^x)\otimes C_a^x\ket\psi\right|
 \leq\sqrt\eps,
\]
by direct-sum Cauchy--Schwarz and $\sum_a(C_a^x)^2\leq I$.
\end{proof}

\begin{proposition}[Rounding inside a consistency relation]
\label{prop:round-right-submeasurement}
Let $\ket\psi\in\calH_{\mathrm A}\otimes\calH_{\mathrm B}$, let
$A=\{A^x_a\}_{a\in\calA}$ be a measurement on $\calH_{\mathrm A}$,
and let $X=\{X^x_b\}_{b\in\calB}$ and
$P=\{P^x_b\}_{b\in\calB}$ be, respectively, a sub-measurement and a
projective sub-measurement on $\calH_{\mathrm B}$.  For each~$x$, let
$f_x:\calB\to\calA$.  If
\[
A^x_a\otimes I\simeq_\delta I\otimes X^x_{[f_x(b)=a]},
\qquad
I\otimes X^x_b\approx_\eta I\otimes P^x_b,
\]
then
\[
A^x_a\otimes I
\simeq_{(\sqrt\delta+\sqrt\eta)^2}
I\otimes P^x_{[f_x(b)=a]}.
\]
If $\calH_{\mathrm A}=\calH_{\mathrm B}$ and $\ket\psi$ is
permutation invariant, the second hypothesis may instead be given on
the left tensor factor.  There is also an analogous statement with
the two tensor factors exchanged.
\end{proposition}
\begin{proof}
Since each $P^x_b$ is a projection, direct-sum Minkowski gives
\begin{align*}
&\left(\E_\bx\sum_{a,b:a\ne f_\bx(b)}
\bra\psi A^{\bx}_a\otimes P^{\bx}_b\ket\psi\right)^{1/2}\\
&\quad=\left(\E_\bx\sum_{a,b:a\ne f_\bx(b)}
\|((A^{\bx}_a)^{1/2}\otimes P^{\bx}_b)\ket\psi\|^2\right)^{1/2}\\
&\quad\leq
\left(\E_\bx\sum_{a,b:a\ne f_\bx(b)}
\|((A^{\bx}_a)^{1/2}\otimes(P^{\bx}_b-X^{\bx}_b))\ket\psi\|^2\right)^{1/2}\\
&\qquad+
\left(\E_\bx\sum_{a,b:a\ne f_\bx(b)}
\|((A^{\bx}_a)^{1/2}\otimes X^{\bx}_b)\ket\psi\|^2\right)^{1/2}.
\end{align*}
The square of the first term is at most~$\eta$, because
$\sum_{a:a\ne f_x(b)}A^x_a\leq I$ for every~$(x,b)$.  The square of
the second term is at most
$\delta$, because $(X^x_b)^2\leq X^x_b$ and the first hypothesis
bounds the inconsistency.  Squaring proves the claim.
The permutation-invariant version follows from
\Cref{prop:approx-swap-invariance}; exchanging the tensor factors
proves the final assertion.
\end{proof}

Next, we show that in the case of measurements,
the state-dependent distance is a weakening of the consistency.

\begin{proposition}[``$\simeq$" implies ``$\approx$" for measurements]\label{prop:simeq-to-approx}
Let $A = \{A^x_a\}$ and $B = \{B^x_a\}$ be two measurements such that
\begin{equation*}
A^x_a \ot I \simeq_{\delta} I \ot B^x_a.
\end{equation*}
Then
\begin{equation*}
A^x_a \ot I \approx_{2\delta} I \ot B^x_a.
\end{equation*}
This is an ``if and only if'' if~$A$ and~$B$ are both projective measurements.
\end{proposition}
\begin{proof}
Expanding the squared distance and using $A_a^2\leq A_a$,
$B_a^2\leq B_a$, and completeness gives
\[
 \E_x\sum_a\|(A_a^x\otimes I-I\otimes B_a^x)\psi\|^2
 \leq2-2\E_x\sum_a\bra\psi A_a^x\otimes B_a^x\ket\psi
 \leq2\delta.
\]
For projective measurements the first inequality is an equality,
which also proves the converse.
\end{proof}

\begin{remark}
We note that \Cref{prop:simeq-to-approx} certainly does \emph{not} hold for sub-measurements.
For example, if $A^x_a = 0$ for all~$a$,
then $A_a^x \ot I \simeq_{0} I \ot B_a^x$,
but
\begin{equation*}
\E_{\bx} \sum_a \Vert (A^{\bx}_a \ot I - I \ot B^{\bx}_a)\ket{\psi} \Vert^2
= \E_{\bx} \sum_a \Vert (I \ot B^{\bx}_a)\ket{\psi} \Vert^2,
\end{equation*}
which is nonzero unless $(I \ot B^x_a) \ket{\psi} = 0$ for all~$x$ and~$a$.
\end{remark}

\subsubsection{Deriving consistency relations from the state-dependent distance}
\label{sec:consistency-from-state-dependent-distance}

The next two rules justify replacing a factor inside an expectation
or an operator product.  Their normalization hypotheses make the
remaining factors contractions in the direct-sum norm.

\begin{proposition}
\label{prop:closeness-of-ip}
	Let $\{A^x_a\}$, $\{B^x_a\}$, and $\{C^x_{a,b} \}$ be matrices.
	Suppose that  $A^x_a \approx_\gamma B^x_a$ and that for all~$x$, $\sum_a (\sum_b C^x_{a,b}) (\sum_b C^x_{a,b})^\dagger  \leq I$. Then
	\begin{equation}
	\E_{\bx} \sum_{a,b} \bra{\psi} C^{\bx}_{a,b} A^{\bx}_a  \ket{\psi} \approx_{\sqrt{\gamma}} \E_{\bx} \sum_{a,b} \bra{\psi} C^{\bx}_{a,b} B^{\bx}_a  \ket{\psi} \,. \label{eq:closeness3}
	\end{equation}
	Similarly, suppose that  $(A^x_a)^\dagger \approx_\gamma (B^x_a)^\dagger$ and that for all~$x$, $\sum_a (\sum_b C^x_{a,b})^\dagger (\sum_b C^x_{a,b}) \leq I$. Then
	\begin{equation}
	\E_{\bx} \sum_{a,b} \bra{\psi} A^{\bx}_a  C^{\bx}_{a,b} \ket{\psi} \approx_{\sqrt{\gamma}} \E_{\bx} \sum_{a,b} \bra{\psi} B^{\bx}_a  C^{\bx}_{a,b}  \ket{\psi} \,. \label{eq:closeness4}
	\end{equation}
\end{proposition}
\begin{proof}
Set $D_a^x=\sum_b C_{a,b}^x$.  The difference in
\Cref{eq:closeness3} is the direct-sum inner product of
$((D_a^x)^\dagger\psi)_{x,a}$ and $((A_a^x-B_a^x)\psi)_{x,a}$.
The first family has norm at most one by the normalization
hypothesis, and the second has norm at most $\sqrt\gamma$.
Cauchy--Schwarz proves the bound.  Taking adjoints proves
\Cref{eq:closeness4} with its stated normalization.
\end{proof}

The product version is \cite[Fact~4.20]{NW19}.
    \begin{proposition}
      \label{prop:cab-approx-delta}
      Let $\{A^x_a\}, \{B^x_a\},$ and $\{C^x_{a,b}\}$ be
      matrices. Suppose that $A^{x}_a \approx_\delta B^{x}_a$ and that
      for all $x$ and $a$, $\sum_b (C^{x}_{a,b})^\dagger
      (C^{x}_{a,b}) \leq I$. Then
      \[ C^{x}_{a,b} A^x_{a} \approx_{\delta} C^{x}_{a,b} B^{x}_a. \]
    \end{proposition}
    \begin{proof}
Apply the contraction-family estimate in
\Cref{lem:direct-sum-estimates} to $v_a^x=(A_a^x-B_a^x)\psi$.
Its squared norm is at most $\delta$.
\end{proof}

\subsubsection{Miscellaneous distance properties}

We now state a few miscellaneous properties of our two distances.

\begin{proposition}[Triangle inequality for ``$\approx_{\delta}$"]
\label{prop:triangle-inequality-for-approx_delta}
Suppose $A_1 = \{(A_1)^x_a\}, \ldots, A_{k+1} = \{(A_{k+1})^x_a\}$ is a set of matrices such that
\begin{equation*}
(A_i)^x_a \approx_{\delta_i} (A_{i+1})^x_a
\end{equation*}
for all $i \in [k]$. Then
\begin{equation*}
(A_1)^x_a \approx_{(\sqrt{\delta_1}+\cdots+\sqrt{\delta_k})^2} (A_{k+1})^x_a.
\end{equation*}
\end{proposition}
\begin{proof}
Apply the direct-sum triangle inequality to the telescoping sum
\[
 ((A_1)^x_a-(A_{k+1})^x_a)\psi
 =\sum_{i=1}^k((A_i)^x_a-(A_{i+1})^x_a)\psi.
\]
Its norm is at most $\sum_i\sqrt{\delta_i}$; squaring proves the claim.
\end{proof}

We note that \Cref{prop:triangle-inequality-for-approx_delta}
contrasts with the triangle inequality for ``$\approx_{\delta}$'' when applied to numbers,
i.e.\ \Cref{eq:triangle-inequality-for-numbers}, because the parameter in
\Cref{def:approx_delta} is the square of the underlying norm.

\begin{lemma}[Disagreement through a projective measurement]
\label{lem:projective-bridge}
Let $\ket\psi\in\calH\otimes\calH$ be permutation invariant, let
$P=\{P_a^x\}$ be projective measurements, and let $C=\{C_a^x\}$ be
sub-measurements with the same outcome set.  If, on the same question
distribution,
\[
 P_a^x\otimes I\simeq_c I\otimes C_a^x,
 \qquad P_a^x\otimes I\simeq_\delta I\otimes P_a^x,
\]
then
\begin{equation}\label{eq:self-improvement-disagreement}
 \E_x\sum_{a\ne b}\bra\psi C_a^x\otimes C_b^x\ket\psi
 \leq(2\sqrt c+\sqrt{2\delta})^2.
\end{equation}
\end{lemma}
\begin{proof}
Suppress $x$ and take the direct sum over the question and $a\ne b$.
Use the vectors
\[
 v_{ab}=(\sqrt{C_a}\otimes\sqrt{C_b})\psi,\quad
 z_{ab}=(\sqrt{C_a}P_b\otimes\sqrt{C_b})\psi,\quad
 w_{ab}=(\sqrt{C_a}\otimes\sqrt{C_b}P_b)\psi.
\]
Since $\sum_{a\ne b}C_a\leq I$ and $P_b$ is a projection,
\[
 \|v-z\|^2\leq\E_x\sum_b
 \bra\psi(I-P_b)\otimes C_b\ket\psi\leq c.
\]
The squared coefficient of the $b$-th mirror difference in $z-w$ is
$(\sum_{a\ne b}C_a)\otimes C_b\leq I$.  Thus
\Cref{prop:simeq-to-approx,lem:direct-sum-estimates} give
$\|z-w\|^2\leq2\delta$.
Finally, $P_bC_bP_b\leq P_b$ and permutation invariance give
\[
 \|w\|^2\leq\E_x\sum_{a\ne b}
 \bra\psi C_a\otimes P_b\ket\psi\leq c.
\]
The triangle inequality bounds $\|v\|$ by $2\sqrt c+\sqrt{2\delta}$.
\end{proof}

Post-processing cannot increase inconsistency, even for
sub-measurements: $f(a)\ne f(a')$ implies $a\ne a'$, so
\[
 \E_x\sum_{b\ne b'}\bra\psi
 A^x_{[f(a)=b]}\otimes B^x_{[f(a)=b']}\ket\psi
 \leq\E_x\sum_{a\ne a'}\bra\psi A_a^x\otimes B_{a'}^x\ket\psi.
\]

The following fact is useful for translating between statements about consistency and closeness between sub-measurements.
\begin{proposition}
\label{prop:cons-sub-meas}
Let $\{A^x_a\}$ be a sub-measurement and let $\{B^x_a\}$ be a measurement such that on average over $x$,
\[
	A^x_a \ot I \simeq_\gamma I \ot B^x_a \,.
\]
Then the following hold
	\begin{gather}
		A^x_a \ot I \approx_{\gamma} A^x_a \ot B^x_a \approx_{\gamma} A^x \ot B^x_a,
		\label{eq:closeness5}
	\end{gather}
	where $A^x = \sum_a A^x_a$.
	As a result, by \Cref{prop:triangle-inequality-for-approx_delta},
	\begin{equation*}
	A^x_a \ot I \approx_{4\gamma} A^x \ot B^x_a
	\end{equation*}
\end{proposition}
\begin{proof}
The positive contractions
$X_a^x=A_a^x\otimes(I-B_a^x)$ and
$Y_a^x=(A^x-A_a^x)\otimes B_a^x$ are the differences in the two
approximations.  Since $X^2\leq X$ for every positive contraction,
\[
 \E_x\sum_a\bra\psi(X_a^x)^2\ket\psi
 \leq\E_x\sum_{a\ne b}\bra\psi A_a^x\otimes B_b^x\ket\psi\leq\gamma,
\]
and the same calculation with $Y_a^x$ gives the second bound.
The final assertion follows from the root-distance triangle inequality.
\end{proof}

\begin{proposition}  \label{prop:switch-sandwich}
  Suppose $\{A^x_a\}$ is a projective sub-measurement
  satisfying
  \begin{equation}
    A^x_a \ot I \approx_{\delta} I \ot A^x_a. \label{eq:Aapproxd}
  \end{equation}
  Then for any $0 \leq B \leq I$, the following holds:
  \begin{equation}
    \E_{\bx}\sum_a \bra{\psi} A^{\bx}_a B A^{\bx}_a \ot I \ket{\psi}
    \approx_{2\sqrt{\delta}} \E_{\bx} \sum_a \bra{\psi} B
    \ot A^{\bx}_a
    \ket{\psi} \approx_{\sqrt{\delta}} \E_{\bx} \sum_a \bra{\psi}
    BA^{\bx}_a \ot I \ket{\psi} \label{eq:switch-sandwich} \end{equation}

\end{proposition}
\begin{proof}
In the direct sum over $x,a$, put
$v_a^x=(A_a^x\otimes I)\psi$ and
$w_a^x=(I\otimes A_a^x)\psi$.  Both norms are at most one and
$\|v-w\|\leq\sqrt\delta$.  Apply
\Cref{eq:direct-sum-quadratic-form} with $M=B\otimes I$.
Projectivity identifies its two quadratic forms with the first two
expressions in \Cref{eq:switch-sandwich}, proving their
$2\sqrt\delta$ bound.

For the second bound, set $S^x=\sum_aA_a^x$.  Orthogonality gives
\begin{align*}
 &\sum_a(A_a^x\otimes I-I\otimes A_a^x)^2
       -(S^x\otimes I-I\otimes S^x)^2\\
 &\qquad=2\sum_{a\ne b}A_a^x\otimes A_b^x\geq0.
\end{align*}
Hence $\E_x\|(S^x\otimes I-I\otimes S^x)\psi\|^2\leq\delta$.
Cauchy--Schwarz with $(B\otimes I)\psi$ now bounds the difference
between the last two expressions in \Cref{eq:switch-sandwich} by
$\sqrt\delta$.
\end{proof}


\subsubsection{Strong self-consistency}

An important property of a sub-measurement~$A = \{A^x_a\}$
is that if both provers measure using~$A$,
then they receive the same outcome.
It seems natural to study this using \emph{self-consistency} of~$A$, i.e. the number~$\delta$ such that
\begin{equation*}
A^x_a \ot I \simeq_{\delta} I \ot A^x_a.
\end{equation*}
However, when~$A$ is a sub-measurement,
being $\delta$-self consistent only bounds the probability that
both provers return outcomes and those outcomes differ.  It does not
bound the probability that one prover succeeds while the other
returns no outcome.
This motivates defining the following stronger notion of self-consistency.

\begin{definition}[Strong self consistency]\label{def:strong-self-consistency}
Let $\ket{\psi}$ be a permutation-invariant state in $\calH \ot \calH$
and let $A = \{A^x_a\}$ be a sub-measurement acting on~$\calH$.
Then $A$ is \emph{$\delta$-strongly self consistent} if
\begin{equation*}
\E_{\bx} \sum_a \bra{\psi} A^{\bx}_a \otimes A^{\bx}_a \ket{\psi}
\geq
\bra{\psi} A \otimes I \ket{\psi} -  \delta.
\end{equation*}
\end{definition}

Strong self-consistency bounds ordinary inconsistency because
\[
 \E_x\sum_{a\ne b}\bra\psi A_a^x\otimes A_b^x\ket\psi
 \leq\bra\psi A\otimes I\ket\psi
       -\E_x\sum_a\bra\psi A_a^x\otimes A_a^x\ket\psi.
\]
For measurements this is an equality.  Its relation to the
state-dependent distance is the following.

\begin{proposition}\label{prop:two-notions-of-self-consistency}
Let $A=\{A_a^x\}$ be a sub-measurement, and let $\ket\psi$
be a permutation-invariant state. If
\begin{equation*}
\E_{\bx} \sum_a \bra{\psi} A^{\bx}_a \otimes A^{\bx}_a \ket{\psi} \geq \bra{\psi} A \otimes I \ket{\psi} - \delta
\end{equation*}
then $A^x_a \otimes I \approx_{2\delta} I \otimes A^x_a$.
This is an ``if and only if" if $A$ is projective.
\end{proposition}
\begin{proof}
Expanding the squares and using permutation invariance gives
\begin{align}
 &\E_x\sum_a\|(A_a^x\otimes I-I\otimes A_a^x)\psi\|^2\nonumber\\
 &=2\E_x\sum_a\bra\psi
       (A_a^x)^2\otimes I-A_a^x\otimes A_a^x\ket\psi\nonumber\\
 &\leq2\left(\bra\psi A\otimes I\ket\psi
       -\E_x\sum_a\bra\psi A_a^x\otimes A_a^x\ket\psi\right)
 \leq2\delta.\label{eq:self-consistency-projective-step}
\end{align}
The inequality comparing the squared distance with twice the
strong self-consistency defect is an equality when $A$ is projective,
which proves the converse as well.
\end{proof}

Hence, we may also refer to the condition $A^x_a \ot I \approx_{2\delta} I \ot A^x_a$
as ``strong self-consistency'' if~$A$ is projective.

We will also use the following property of strongly self-consistent
sub-measurements after evaluating their outcomes.

\begin{proposition}\label{prop:two-notions-of-self-consistency-after-evaluation}
Let $A=\{A_a^x\}$ be a sub-measurement on a
permutation-invariant state $\ket\psi$, and suppose that
\begin{equation*}
\E_{\bx} \sum_a \bra{\psi} A^{\bx}_a \otimes A^{\bx}_a \ket{\psi} \geq  \bra{\psi} A \otimes I \ket{\psi} - \delta.
\end{equation*}
Then for any function~$f$,
$A^x_{[f(a)=b]} \otimes I \approx_{2\delta} I \otimes A^x_{[f(a)=b]}$.
\end{proposition}
\begin{proof}
Write $C_b^x=A^x_{[f(a)=b]}$.  This sub-measurement has the same
total effect as $A^x$, and positivity gives
\[
 \E_x\sum_b\bra\psi C_b^x\otimes C_b^x\ket\psi
 =\E_x\sum_{a,a':f(a)=f(a')}\bra\psi A_a^x\otimes A_{a'}^x\ket\psi
 \geq\E_x\sum_a\bra\psi A_a^x\otimes A_a^x\ket\psi.
\]
Thus its strong self-consistency defect is at most $\delta$.
\Cref{prop:two-notions-of-self-consistency} yields
\begin{equation}\label{eq:evaluated-self-consistency-bound}
 \E_x\sum_b\|(C_b^x\otimes I-I\otimes C_b^x)\psi\|^2\leq2\delta.
\end{equation}
The same proof allows the map $f$ to depend on the question $x$.
\end{proof}


\section{Making measurements projective}
\label{sec:making-measurements-projective}

Projective measurements are substantially easier to manipulate than general
POVMs.  In particular, for projective measurements the consistency and
state-dependent distance agree exactly up to the normalization in
\Cref{prop:simeq-to-approx}:
\begin{equation*}
A^x_a \ot I \simeq_{\delta} I \ot B^x_a
\qquad\Longleftrightarrow\qquad
A^x_a \ot I \approx_{2\delta} I \ot B^x_a.
\end{equation*}

The conceptual parameter for orthogonalization is the state-dependent
idempotence defect $\sum_a\varphi(M_a-M_a^2)$.  The next lemma rounds
this defect linearly.  It is the finite-dimensional case of the
orthogonalization theorem of de~la~Salle~\cite{deLaSalle22}; we give the
argument to keep the final rounding self-contained.  The induction
itself uses simultaneous dilation and compression, as proved in
\Cref{sec:self-improvement-dilation}.

\subsection{State-dependent orthogonalization}\label{sec:orthogonalization}

\begin{lemma}[State-dependent orthogonalization]
\label{lem:state-dependent-orthogonalization}
Let $\calH$ be $d$-dimensional, let $\varphi$ be a state on
$B(\calH)$, and let $M=\{M_a\}$ be a finite-outcome measurement.  Put
\begin{equation}\label{eq:orthogonalization-defect}
 \Delta
 :=1-\varphi\left(\sum_aM_a^2\right)
 =\sum_a\varphi(M_a-M_a^2).
\end{equation}
Then there is a projective measurement $P=\{P_a\}$ on $\calH$ such that
\begin{equation}\label{eq:linear-orthogonalization}
 \sum_a\varphi\bigl((M_a-P_a)^2\bigr)\leq9\Delta.
\end{equation}
\end{lemma}

\begin{proof}
For each $a$, choose a spectral decomposition
\begin{equation*}
 M_a=\sum_{j=1}^d\lambda_{a,j}
       \ket{v_{a,j}}\!\bra{v_{a,j}},
 \qquad 0\leq\lambda_{a,j}\leq1,
\end{equation*}
and define
\begin{equation*}
 r_{a,j}=\varphi\bigl(\ket{v_{a,j}}\!\bra{v_{a,j}}\bigr),
 \qquad w_{a,j}=\lambda_{a,j}r_{a,j}.
\end{equation*}
Taking the ordinary trace in $\sum_aM_a=I$ gives
\begin{equation*}
 \sum_{a,j}\lambda_{a,j}=d.
\end{equation*}
For nonnegative weights, the sum of the $d$ largest weights is the
maximum of $\sum_{a,j}x_{a,j}w_{a,j}$ over
$0\leq x_{a,j}\leq1$ and $\sum_{a,j}x_{a,j}=d$: moving mass from a
smaller weight to a larger one can only increase the objective.
The fractional choice $x_{a,j}=\lambda_{a,j}$ is feasible.  Select the
$d$ largest $w_{a,j}$, breaking ties arbitrarily, and let $Q_a$ project
onto the selected eigenvectors belonging to outcome $a$.  Then
\begin{equation}\label{eq:global-rank-allocation}
 [Q_a,M_a]=0,
 \qquad \sum_a\operatorname{rank}(Q_a)=d,
\end{equation}
and
\begin{align}
 s:=\sum_a\varphi(Q_aM_a)
 &=\sum_{(a,j)\text{ selected}}w_{a,j}\nonumber\\
 &\geq\sum_{a,j}\lambda_{a,j}w_{a,j}
 =\sum_a\varphi(M_a^2)
 =1-\Delta.\label{eq:selected-overlap}
\end{align}
The $Q_a$ need not be mutually orthogonal; the point of the selection
is that their total rank is exactly $d$.

Let
\begin{equation*}
 \mathcal K=\bigoplus_a Q_a\calH,
\end{equation*}
where the summands are abstract orthogonal copies.  By
\Cref{eq:global-rank-allocation}, $\dim\mathcal K=d=\dim\calH$.  Let
$E_a$ be the coordinate projection onto the $a$-th summand and define
\begin{equation*}
 X:\calH\longrightarrow\mathcal K,
 \qquad
 X\ket v=\bigl(Q_aM_a^{1/2}\ket v\bigr)_a.
\end{equation*}
Since $Q_a$ commutes with $M_a$,
\begin{equation*}
 X^\dagger X=\sum_aQ_aM_a=:S\leq\sum_aM_a=I.
\end{equation*}
Write $\abs{X}=S^{1/2}=(X^\dagger X)^{1/2}$.  Extend the partial isometry in the polar
decomposition of $X$ to a unitary $U:\calH\to\mathcal K$, so that
$X=U\abs{X}$.  Such an extension exists even when $X$ is singular because
its domain and codomain have the same finite dimension.  Set
\begin{equation*}
 P_a=U^\dagger E_aU.
\end{equation*}
Then the $P_a$ are mutually orthogonal projections summing to $I$, and
\begin{equation}\label{eq:polar-coordinate-identity}
 \abs{X}P_a\abs{X}=X^\dagger E_aX=Q_aM_a.
\end{equation}

For an operator family $Y=(Y_a)_a$, write
\begin{equation*}
 \|Y\|_{\varphi,2}^2=\sum_a\varphi(Y_a^\dagger Y_a).
\end{equation*}
This is a Hilbert-space seminorm, so Minkowski's inequality applies.
By \Cref{eq:polar-coordinate-identity}, there is an exact decomposition
\begin{equation}\label{eq:orthogonalization-decomposition}
 M_a-P_a
 =(I-Q_a)M_a+(\abs{X}-I)P_a\abs{X}+P_a(\abs{X}-I).
\end{equation}
We bound the three families on the right separately.  First, using
$[Q_a,M_a]=0$, $M_a^2\leq M_a$, and
\Cref{eq:selected-overlap},
\begin{align*}
 \big\|\big((I-Q_a)M_a\big)_a\big\|_{\varphi,2}^2
 &=\sum_a
 \varphi\left(
 \big((I-Q_a)M_a\big)^\dagger
 \big((I-Q_a)M_a\big)
 \right)\\
 &=\sum_a
 \varphi\left(
 M_a(I-Q_a)^2M_a
 \right)\\
 &=\sum_a
 \varphi\left(
 M_a(I-Q_a)M_a
 \right)\\
 &=\sum_a
 \varphi\left(
 (I-Q_a)M_a^2
 \right)\\
 &\leq\sum_a
 \varphi\left(
 (I-Q_a)M_a
 \right)\\
 &=\sum_a\varphi(M_a)
 -\sum_a\varphi(Q_aM_a)\\
 &=\varphi\left(\sum_aM_a\right)-s\\
 &=\varphi(I)-s\\
 &=1-s\\
 &\leq\Delta.
\end{align*}
Second, $0\leq\abs{X}\leq I$ implies $(I-\abs{X})^2\leq I-\abs{X}^2$.  If
$C_a=\abs{X}P_a\abs{X}=Q_aM_a$, then
\begin{align*}
 \big\|\big((\abs{X}-I)P_a\abs{X}\big)_a\big\|_{\varphi,2}^2
 &=\sum_a\varphi\bigl(\abs{X}P_a(I-\abs{X})^2P_a\abs{X}\bigr)\\
 &\leq\sum_a\varphi\bigl(\abs{X}P_a(I-\abs{X}^2)P_a\abs{X}\bigr)\\
 &=\sum_a\varphi(C_a-C_a^2)\\
 &=\sum_a\varphi\bigl(Q_a(M_a-M_a^2)\bigr)\leq\Delta.
\end{align*}
The last inequality is valid because $Q_a$ commutes with the positive
operator $M_a-M_a^2$.  Finally, using $\sum_aP_a=I$,
\begin{align*}
 \big\|\big(P_a(\abs{X}-I)\big)_a\big\|_{\varphi,2}^2
 &=\varphi((I-\abs{X})^2)\\
 &\leq\varphi(I-\abs{X}^2)=1-s\leq\Delta.
\end{align*}
Minkowski's inequality in \Cref{eq:orthogonalization-decomposition}
now gives
\begin{equation*}
 \left(\sum_a\varphi((M_a-P_a)^2)\right)^{1/2}
 \leq3\sqrt\Delta,
\end{equation*}
which proves \Cref{eq:linear-orthogonalization}.
\end{proof}

We record the consequence used for the final rounding in
\Cref{thm:main-formal}.  It does not require a symmetric state or
identical measurements on the two tensor factors.

\begin{lemma}[Consistent measurements]
\label{lem:orthonormalization-main-lemma}
Let $\calH_{\mathrm A}$ and $\calH_{\mathrm B}$ be finite-dimensional,
let $\ket\psi\in\calH_{\mathrm A}\ot\calH_{\mathrm B}$ be a state, and
let $A=\{A_a\}$ and $B=\{B_a\}$ be finite-outcome measurements with the
same outcome set.  If
\begin{equation*}
 A_a\ot I\simeq_\zeta I\ot B_a,
\end{equation*}
then there is a projective measurement $P=\{P_a\}$ on
$\calH_{\mathrm A}$ such that
\begin{equation}\label{eq:consistent-measurement-orthogonalization}
 A_a\ot I\approx_{18\zeta}P_a\ot I.
\end{equation}
\end{lemma}

\begin{proof}
Let
\begin{equation*}
 e=1-\sum_a\bra\psi A_a\ot B_a\ket\psi
\end{equation*}
be the actual inconsistency.  By
\Cref{prop:simeq-for-measurements}, $0\leq e\leq\zeta$.  Direct-sum
Cauchy--Schwarz and $\sum_aB_a^2\leq I$ give
\begin{align*}
 (1-e)^2
 &\leq
 \left(\sum_a\bra\psi A_a^2\ot I\ket\psi\right)
 \left(\sum_a\bra\psi I\ot B_a^2\ket\psi\right)\\
 &\leq\sum_a\bra\psi A_a^2\ot I\ket\psi.
\end{align*}
Apply \Cref{lem:state-dependent-orthogonalization} to $A$ and the state
$\varphi(X)=\bra\psi X\ot I\ket\psi$.  Its defect satisfies
\begin{equation*}
 \Delta_A\leq1-(1-e)^2=2e-e^2\leq2\zeta,
\end{equation*}
so the resulting projective measurement obeys
\begin{equation*}
 \sum_a\|(A_a-P_a)\ot I\ket\psi\|^2
 \leq9(2e-e^2)\leq18\zeta.
\end{equation*}
\end{proof}

\section{The main induction step}
\label{sec:induction}

We will now carry out the main inductive argument.
The inductive hypothesis is stated as follows.

\begin{theorem}[Main induction]\label{thm:main-induction}
Let $m\geq1$, and let $(\psi,A,B,L)$ be an $(\eps,\delta,\gamma)$-good symmetric strategy
for the $(m,q,d)$ low individual degree test.  There is a measurement
$G\in\polymeas{m}{q}{d}$ such that, on average over
$\bu\sim\F_q^m$,
\[
  A^u_a\otimes I\simeq_{\sigma}I\otimes G_{[g(u)=a]},
\]
where, if $d=0$,
\[
 \sigma=\delta+\sqrt{8m\eps},
\]
whereas, if $d\geq1$,
\[
 \sigma=1000d^2m^4\left(
 \eps^{1/32}+\delta^{1/32}+(d/q)^{1/32}
 +\gamma^{1/8}\right).
\]
\end{theorem}

Comparing with our main theorem, \Cref{thm:main-formal},
\Cref{thm:main-induction} produces a measurement which is consistent with~$A$,
but it is not projective or self-consistent.
In addition, the strategy is assumed to be symmetric.
We correct these deficiencies in the following proof.

\begin{proof}[Proof of~\Cref{thm:main-formal} assuming \Cref{thm:main-induction}]
Let
\[
 \theta=\frac dq,\qquad
 K=K_{m,d},\qquad
 R=K\bigl(\eps^{1/64}+\theta^{1/64}\bigr),
\]
so the error in \Cref{thm:main-formal} is \(\nu=21000R\).
If \(\nu\geq1\), fix any \(g_0\in\polyfunc{m}{q}{d}\) and let both
players use the projective measurement supported on \(g_0\).  Its
self-inconsistency is zero, and each point inconsistency is at most
one, so the theorem is immediate.  We therefore assume
\begin{equation}\label{eq:nontrivial-main-error}
 R<\frac1{21000}.
\end{equation}
In particular, \(\eps,\theta<1\), and hence \(d<q\).

Call the strategy in \Cref{thm:main-formal} the original strategy.
Because the three subtests are chosen with equal probability, the
original strategy is \((3\eps,3\eps,3\eps)\)-good.  We symmetrize it
by a direct-sum construction, without first padding the local spaces
to a common dimension.  Let
\[
 \mathcal K=\calH_{\mathrm A}\oplus\calH_{\mathrm B},
\]
and let \(\iota_{\mathrm A}\) and \(\iota_{\mathrm B}\) denote the
canonical embeddings of the two summands into~\(\mathcal K\).  Define
\[
 \ket{\psi_{\rm sym}}
 =\frac1{\sqrt2}\left(
   (\iota_{\mathrm A}\otimes\iota_{\mathrm B})\ket\psi
   +(\iota_{\mathrm B}\otimes\iota_{\mathrm A})
      \ket{\psi_{\rm swap}}\right)\in\mathcal K\otimes\mathcal K,
\]
where \(\ket{\psi_{\rm swap}}\in
\calH_{\mathrm B}\otimes\calH_{\mathrm A}\) is obtained by swapping
the two registers of~\(\ket\psi\).  Componentwise, set
\[
 A_{\rm sym}=A^{\mathrm A}\oplus A^{\mathrm B},\qquad
 B_{\rm sym}=B^{\mathrm A}\oplus B^{\mathrm B},\qquad
 L_{\rm sym}=L^{\mathrm A}\oplus L^{\mathrm B}.
\]
This is a symmetric \((3\eps,3\eps,3\eps)\)-good strategy.

The two cases of \Cref{thm:main-induction} imply that there is a
measurement \(G\in\polymeas{m}{q}{d}\) such that, on average over~\(u\),
\begin{equation}\label{eq:short-applied-induction}
 (A_{\rm sym})^u_a\otimes I
 \simeq_\sigma I\otimes G_{[g(u)=a]},
 \qquad
 \sigma=10000K\bigl(\eps^{1/32}+\theta^{1/32}\bigr).
\end{equation}
For \(d=0\), this bounds \(3\eps+\sqrt{24m\eps}\); for \(d\geq1\),
it follows by substituting the three errors \(3\eps\) into the stated
induction bound and using \(\eps<1\).
Let \(\Pi_w=\iota_w\iota_w^\dagger\) for
\(w\in\{\mathrm A,\mathrm B\}\), and pinch \(G\) to the two summands:
\[
 \overline G_g=\Pi_{\mathrm A}G_g\Pi_{\mathrm A}
               +\Pi_{\mathrm B}G_g\Pi_{\mathrm B}.
\]
The operators \(\overline G_g\) form a measurement.  Replacing \(G\)
by \(\overline G\) does not change the inconsistency in
\eqref{eq:short-applied-induction}: \(A_{\rm sym}\) is block diagonal,
and the two terms of \(\ket{\psi_{\rm sym}}\) occupy orthogonal
direct-sum sectors.  Thus \eqref{eq:short-applied-induction} holds with
\(\overline G\) in place of~\(G\).

The symmetrized strategy also satisfies
\[
 (A_{\rm sym})^u_a\otimes I
 \simeq_{3\eps}I\otimes(A_{\rm sym})^u_a.
\]
Apply \Cref{lem:projective-bridge} to the projective point
measurement $(A_{\rm sym})^u$ and the evaluated measurement
$\overline G_{[g(u)=a]}$ on the permutation-invariant state
$\ket{\psi_{\rm sym}}$.  Together with
\eqref{eq:short-applied-induction}, this gives
\begin{equation}\label{eq:short-evaluation-self-consistency}
 \overline G_{[g(u)=a]}\otimes I
 \simeq_\tau I\otimes\overline G_{[g(u)=a]},
 \qquad
 \tau=(2\sqrt\sigma+\sqrt{6\eps})^2.
\end{equation}
For distinct individual-degree-\(d\) polynomials \(g,h\),
Schwartz--Zippel gives
\(\Pr_{\bu}[g(\bu)=h(\bu)]\leq md/q=m\theta\).  If
\(w_{g,h}=\bra{\psi_{\rm sym}}\overline G_g\otimes\overline G_h
\ket{\psi_{\rm sym}}\), then
\[
 \sum_{g\ne h}\Pr_{\bu}[g(\bu)=h(\bu)]w_{g,h}
 \leq m\theta\sum_{g\ne h}w_{g,h}\leq m\theta.
\]
Expanding \eqref{eq:short-evaluation-self-consistency} in the
polynomial outcomes therefore gives
\begin{equation}\label{eq:short-symmetric-polynomial-consistency}
 \sum_{g\ne h}
 \bra{\psi_{\rm sym}}\overline G_g\otimes\overline G_h
 \ket{\psi_{\rm sym}}
 \leq\zeta,
 \qquad
 \zeta=\tau+m\theta.
\end{equation}

Define the diagonal blocks
\[
 G^w_g=\iota_w^\dagger\overline G_g\iota_w,
 \qquad w\in\{\mathrm A,\mathrm B\}.
\]
They are measurements on the original local spaces.  Since
\(\overline G\) is block diagonal, the left-hand side of
\eqref{eq:short-symmetric-polynomial-consistency} is exactly
\[
 \sum_{g\ne h}\bra\psi
 G^{\mathrm A}_g\otimes G^{\mathrm B}_h\ket\psi.
\]
Indeed, the two direct-sum sectors give the same sum after exchanging
the indices $g$ and~$h$.
Consequently,
\begin{equation}\label{eq:short-polynomial-consistency}
 G^{\mathrm A}_g\otimes I
 \simeq_\zeta I\otimes G^{\mathrm B}_g.
\end{equation}
Moreover, decomposing the inconsistency in
\eqref{eq:short-applied-induction} into its two orthogonal sectors
shows that
\begin{align}
 A^{\mathrm A,u}_a\otimes I
 &\simeq_{2\sigma}
 I\otimes G^{\mathrm B}_{[g(u)=a]},\label{eq:short-point-cons-a}\\
 G^{\mathrm A}_{[g(u)=a]}\otimes I
 &\simeq_{2\sigma}
 I\otimes A^{\mathrm B,u}_a.\label{eq:short-point-cons-b}
\end{align}
Indeed, the symmetrized inconsistency is one half the sum of the two
displayed original-strategy inconsistencies, so each is at most
\(2\sigma\).

Apply \Cref{lem:orthonormalization-main-lemma} to
\eqref{eq:short-polynomial-consistency} in both tensor orders.  This
gives projective measurements
\(Q^{\mathrm A},Q^{\mathrm B}\in\polymeas{m}{q}{d}\) such that
\begin{equation}\label{eq:short-rounding-error}
 G^{\mathrm A}_g\otimes I\approx_\alpha Q^{\mathrm A}_g\otimes I,
 \qquad
 I\otimes G^{\mathrm B}_g\approx_\alpha I\otimes Q^{\mathrm B}_g,
 \qquad
 \alpha=18\zeta.
\end{equation}
Together with \Cref{prop:simeq-to-approx}, the root-distance triangle
inequality gives
\[
 Q^{\mathrm A}_g\otimes I
 \approx_{(2\sqrt\alpha+\sqrt{2\zeta})^2}
 I\otimes Q^{\mathrm B}_g.
\]
The parameter on the right is \(98\zeta\).  Since both measurements
are projective, their actual inconsistency is at most \(49\zeta\).

For point consistency, apply
\Cref{prop:round-right-submeasurement} with the
question-dependent map \(f_u(g)=g(u)\) to
\Cref{eq:short-point-cons-a,eq:short-rounding-error}, and apply its
tensor-reversed form to
\Cref{eq:short-point-cons-b,eq:short-rounding-error}.  This gives
\begin{align}
 A^{\mathrm A,u}_a\otimes I
 &\simeq_{(\sqrt{2\sigma}+\sqrt\alpha)^2}
 I\otimes Q^{\mathrm B}_{[g(u)=a]},\label{eq:short-rounded-point-a}\\
 Q^{\mathrm A}_{[g(u)=a]}\otimes I
 &\simeq_{(\sqrt{2\sigma}+\sqrt\alpha)^2}
 I\otimes A^{\mathrm B,u}_a.\label{eq:short-rounded-point-b}
\end{align}

It remains to compare these errors with~$\nu$.  Since $K\geq m\geq1$
and $R<1/21000$, the definitions imply
\begin{equation}\label{eq:short-sigma-bound}
 \sigma\leq10000R^2,\qquad
 \eps\leq R^{64},\qquad m\theta\leq R^{64}.
\end{equation}
Consequently,
\[
 \sigma\leq\zeta\leq8\sigma+12\eps+m\theta
 \leq80013R^2\leq4R.
\]
Using $\alpha=18\zeta$, each rounded point inconsistency is at most
$(\sqrt{2\sigma}+\sqrt{18\zeta})^2\leq32\zeta$, and the
self-inconsistency is at most $49\zeta$.  All three are therefore
at most $196R\leq21000R=\nu$, as required.
\end{proof}

The remainder of this section is organized as follows:
first, in \Cref{sec:self-improvement-and-pasting}, we define the two main steps in the proof of \Cref{thm:main-induction}
known as self-improvement and pasting.
Following that, we prove \Cref{thm:main-induction} in \Cref{sec:proof-of-main-induction}.

\subsection{Self-improvement and pasting}\label{sec:self-improvement-and-pasting}

There are two main steps in the proof of \Cref{thm:main-induction}.
The first is self-improvement, which is stated as follows.

\begin{theorem}[Self-improvement with dilation]\label{thm:self-improvement-in-induction-section}
Let $(\psi,A,B,L)$ be an $(\eps,\delta,\gamma)$-good symmetric
strategy on $\calH\otimes\calH$ for the $(m,q,d)$ low individual
degree test.  Let $G\in\polymeas{m}{q}{d}$ be a measurement satisfying
\begin{enumerate}
 \item (Consistency with $A$):\label{item:si-inductive-G-consistency}
 on average over $\bu\sim\F_q^m$,
 \[
 A^u_a\otimes I\simeq_\nu I\otimes G_{[g(u)=a]}.
 \]
\end{enumerate}
Set
\[
 \zeta=200m\bigl(\sqrt\eps+\sqrt\delta+\sqrt{d/q}\bigr).
\]
There exist a finite-dimensional ancilla $\mathcal K$, a unit vector
$\ket0\in\mathcal K$, and a projective sub-measurement
$H\in\polysub{m}{q}{d}$ on $\calH\otimes\mathcal K$ with the
following properties.  Write $Jv=v\otimes\ket0$,
$\ket{\psi'}=(J\otimes J)\ket\psi$, and
$A'{}^u_a=A^u_a\otimes I_{\mathcal K}$; all relations below are on
$\ket{\psi'}$.
\begin{enumerate}
 \item (Completeness):\label{item:si-inductive-completeness}
 if $H=\sum_hH_h$, then
 \[
 \bra{\psi'}H\otimes I\ket{\psi'}\geq1-\nu-\zeta.
 \]
 \item (Consistency with $A'$):\label{item:si-inductive-A-consistency}
 on average over $\bu\sim\F_q^m$,
 \[
 A'{}^u_a\otimes I\simeq_\zeta I\otimes H_{[h(u)=a]}.
 \]
 \item (Strong self-consistency):\label{item:si-inductive-self}
 \[
 H_h\otimes I\approx_\zeta I\otimes H_h.
 \]
 \item (Boundedness):\label{item:si-inductive-boundedness}
 there is a positive contraction $Z$ on $\calH\otimes\mathcal K$ such that
 \[
 \bra{\psi'}Z\otimes(I-H)\ket{\psi'}\leq\zeta,
 \qquad
 Z\geq\E_{\bu}A'{}^{\bu}_{h(\bu)}
 \quad\text{for every }h\in\polyfunc{m}{q}{d}.
 \]
\end{enumerate}
For any finite family of such inputs on $\calH$, the same ancilla and
embedding $J$ can be used for all outputs.  Extending $A,B,L$ by the
identity on the ancilla preserves projectivity and all test success
probabilities.
\end{theorem}

For projective $H$, \Cref{prop:two-notions-of-self-consistency} identifies
\Cref{item:si-inductive-self} with a strong self-consistency defect at
most $\zeta/2$.  Self-improvement moves the seed measurement's
consistency error into incompleteness.  Its filtered output on the
original space need not be projective; simultaneous dilation supplies
the projective slices needed for pasting without a rounding loss.
After pasting, compression recovers a measurement on the original
space, as explained in \Cref{sec:self-improvement-dilation}.

The second main step is pasting, which is stated as follows.

\begin{theorem}[Pasting]\label{thm:ld-pasting-in-induction-section}
  Let $m,d\geq1$, and let $(\psi, A, B, L)$ be an
  $(\eps, \delta, \gamma)$-good symmetric strategy for the
  $(m+1,q,d)$ low individual degree test.
  Let $\{G^x\}_{x \in \F_q}$ denote a set of projective sub-measurements in $\polysub{m}{q}{d}$ with the following properties:
  \begin{enumerate}
  \item (Completeness): \label{item:ld-pasting-inductive-completeness}  If $G = \E_{\bx} \sum_g G^{\bx}_g$, then
	\begin{equation*}
	\bra{\psi} G \otimes I \ket{\psi} \geq 1 - \kappa.
	\end{equation*}
  \item (Consistency with~$A$): \label{item:ld-pasting-inductive-consistency} On average over $(\bu, \bx) \sim \F_q^{m+1}$,
	  \begin{equation*}
	  A^{u, x}_a \otimes I \simeq_{\zeta} I \otimes G^x_{[g(u)=a]}.
	  \end{equation*}
  \item (Strong self-consistency): \label{item:ld-pasting-inductive-self-consistency} On average over~$\bx \sim \F_q$,
	\begin{equation*}
		G^x_g \otimes I \approx_{\zeta} I \otimes G^x_g.
	\end{equation*}
  \item (Boundedness): \label{item:ld-pasting-inductive-boundedness} There exists a positive-semidefinite matrix $Z^x$ for each $x \in \F_q$ such that
	\begin{equation*}
		\E_{\bx} \bra{\psi} (I-G^{\bx})\otimes Z^{\bx} \ket{\psi} \leq \zeta
	\end{equation*}
	and for each $x \in \F_q$ and $g \in \polyfunc{m}{q}{d}$,
	\begin{equation*}
	Z^x \geq \E_{\bu} A^{\bu, x}_{g(\bu)}.
	\end{equation*}
  \end{enumerate}
Let
\[
 s=\eps^{1/8}+\delta^{1/8}+\gamma^{1/8}
   +\zeta^{1/8}+(d/q)^{1/8}.
\]
Then there is a measurement $H\in\polymeas{m+1}{q}{d}$ such that,
on average over $\bu\sim\F_q^{m+1}$,
\begin{equation*}
 A^u_a\otimes I\simeq_{\sigma}I\otimes H_{[h(u)=a]},
 \qquad
 \sigma=\frac{m+1}{m}\,\kappa
       +21(m+1)^2d^2\sqrt{m}\,s.
\end{equation*}
\end{theorem}

This is the specialization of \Cref{thm:ld-pasting} with the attempt
count $k=(m+1)d$: because $m,d\geq1$, one has $k\geq d+1$, and
$k/(k-d)=(m+1)/m$.  Thus $k$ is a choice internal to the construction,
not an auxiliary hypothesis of the induction theorem.

Intuitively, $\sigma$ should be thought of as being roughly~$\kappa$,
plus a small amount of error which does not depend on~$\kappa$.
Hence, pasting states that we can take a family of sub-measurements $\{G^x\}$ with incompleteness $\kappa$
and produce a pasted measurement~$H$ whose consistency error with~$A$ is roughly~$\kappa$,
plus a small amount of new error.
Thus, the overall inductive step looks as follows:
given a family of measurements with some inconsistency error,
we ``move'' the error into the incompleteness using self-improvement,
and then we paste the measurements together to form a single measurement whose error is roughly the same as the original error.

The fresh error $21(m+1)^2d^2\sqrt{m}\,s$ depends only on the small parameters
$\eps$, $\delta$, $\zeta$, $\gamma$, and $d/q$, not on~$\kappa$.

\subsection{\texorpdfstring{Proof of \Cref{thm:main-induction}}{Proof of the main induction}}\label{sec:proof-of-main-induction}

\begin{definition}
Let $x \in \F_q$.
For each line $\ell \in \F_q^m$,
we define $\mathrm{append}_x(\ell)$ to be the line in $\F_q^{m+1}$ containing every point $(u, x)$ such that $u \in \ell$.
In addition, for each function $f : \ell \rightarrow \F_q$,
we define $\mathrm{append}_x(f) : \mathrm{append}_x(\ell) \rightarrow \F_q$
such that for each $(u, x) \in \mathrm{append}_x(\ell)$, $\mathrm{append}_x(f)(u, x) = f(u)$.
\end{definition}

\begin{definition}[$x$-restricted low-degree strategy]
Let $(\psi, A, B, L)$ be a symmetric strategy for the
$(m+1,q,d)$-low individual degree test.
Given $x \in \F_q$, we define the \emph{$x$-restricted strategy}
$(\psi, A^x, B^x, L^x)$ for the $(m,q,d)$-low individual degree test as follows.
\begin{enumerate}
\item For each $u \in \F_q^m$, $(A^x)^u_a = A^{u, x}_a$.
\item For each axis parallel line $\ell \in \F_q^m$, $(B^x)^\ell_f = B^{\mathrm{append}_x(\ell)}_{\mathrm{append}_x(f)}$.
\item For each  line $\ell \in \F_q^m$, $(L^x)^\ell_f = L^{\mathrm{append}_x(\ell)}_{\mathrm{append}_x(f)}$.
\end{enumerate}
\end{definition}

The last item is a relabeling of a complete projective measurement:
$f\mapsto\mathrm{append}_x(f)$ is a bijection between all functions on
$\ell$ and all functions on $\mathrm{append}_x(\ell)$.  This is why the
relaxed diagonal answer alphabet above is useful.  With only the
nominal degree bounds, restricting a degree-$(m+1)d$ outcome would not
necessarily give a degree-$md$ outcome, and $L^x$ could fail to be
complete.

\begin{lemma}\label{lem:restricted-probabilities}
Let $(\psi, A, B, L)$ be an $(\eps, \delta, \gamma)$-good symmetric strategy
		for the $(m+1,q,d)$-low individual degree test.
For each $x \in \F_q$, let $(\psi, A^x, B^x, L^x)$ be the corresponding $x$-restricted strategy.
In addition, write $\eps_x$ for the probability that it fails the axis-parallel lines test,
$\delta_x$ for the probability it fails the self-consistency test,
and $\gamma_x$ for the probability it fails the diagonal lines test.
Then
\begin{equation*}
\E_{\bx} \eps_{\bx} \leq \left(\frac{m+1}{m}\right) \cdot\eps,
\quad
\E_{\bx} \delta_{\bx} \leq  \delta,
\quad
\E_{\bx} \gamma_{\bx} \leq \left(\frac{m+1}{m}\right) \cdot \gamma.
\end{equation*}
\end{lemma}
\begin{proof}
We begin with the self-consistency test.
Here, both provers are given a uniformly random point $(\bu, \bx)$,
they measure using $A^{\bu, \bx} = (A^{\bx})^{\bu}$, and they succeed if their outcomes are the same.
This is equivalent to performing the self-consistency test on the $\bx$-restricted strategy, averaged over $\bx$,
and so $\E_{\bx} \delta_{\bx} \leq  \delta$. (It is a ``$\leq$" rather than an ``$=$" because $\delta$ is just an upper-bound on the failure probability.)

Next, we consider the axis-parallel lines test.
Suppose the provers are sent the line $\bell$ and the point $(\bu, \bx) \in \bell$.
With probability $\frac{1}{m+1}$, $\bell$ is parallel to the $(m+1)$-st direction.
When it is not, then $\bell = \mathrm{append}_x(\bell')$, where $\bell'$ is an axis-parallel line in $\F_q^m$.
In this case, the points prover measures with $A^{\bu, \bx} = (A^\bx)^{\bu}$,
the lines prover measures with $B^{\bell} = (B^{\bx})^{\bell'}$, and the two
outcomes fail the check exactly when the corresponding outcomes in the
$\bx$-restricted test fail it.  Hence,
\begin{align*}
\eps
&\geq \Pr_{\bell, \bu, \bx}[\text{the axis-parallel check fails}]\\
&\geq \left(\frac{m}{m+1}\right) \cdot
 \Pr[\text{the check fails}\mid \text{$\bell$ is not parallel to direction $m+1$}]\\
& =  \left(\frac{m}{m+1}\right) \cdot\E_{\bx} \eps_{\bx}.
\end{align*}

For the diagonal-lines test, condition on the sampled direction index
$i$ lying in $\{1,\ldots,m\}$, an event of probability $m/(m+1)$.
The ambient direction then has the form $(v',0)$, where $v'$ has exactly
the distribution of a diagonal direction in $\F_q^m$, and the ambient
line through $(u,x)$ is $\mathrm{append}_x(\ell')$ for the corresponding
line $\ell'$ through~$u$.  By the definition of $L^x$, the conditional
failure probability is $\E_x\gamma_x$.  Consequently,
\[
 \gamma\geq\frac{m}{m+1}\E_x\gamma_x,
\]
which is the remaining claim.
\end{proof}

\begin{proof}[Proof of~\Cref{thm:main-induction}]
We first dispose of the degree-zero case directly.  Suppose $d=0$.
For each $a\in\F_q$, let $g_a$ be the constant polynomial with value~$a$,
and define
\[
 G_{g_a}=\E_{\bv\sim\F_q^m}A^{\bv}_a.
\]
These operators form a measurement indexed by all polynomials in
$\polyfunc{m}{q}{0}$.

For this proof, sample $(\bu,\bv)$ by drawing
$\bu\sim\F_q^m$, $\bi\sim\{1,\ldots,m\}$, and $\bt\sim\F_q$
independently and setting $\bv=\bu+\bt e_{\bi}$; denote this edge
distribution by~$C$.  Let $\ell$ be the axis-parallel line through them
in the sampled direction.  Write $B^\ell_a$ for the outcome of $B^\ell$
corresponding to the constant function with value~$a$.  The two
point--line test relations, permutation invariance for the reversed
second relation, and \Cref{prop:simeq-to-approx} give
\[
 A^{\bu}_a\otimes I\approx_{2\eps}I\otimes B^\ell_a
 \approx_{2\eps}A^{\bv}_a\otimes I.
\]
Therefore \Cref{prop:triangle-inequality-for-approx_delta} gives
\[
 \E_{(\bu,\bv)\sim C}\sum_a
 \left\|(A^{\bu}_a-A^{\bv}_a)\otimes I\ket\psi\right\|^2
 \leq8\eps.
\]
Applying \Cref{lem:local-to-global} to each outcome and summing over~$a$
shows that
\begin{equation}\label{eq:d-zero-global-variance}
 D:=\E_{\bu,\bv}\sum_a
 \left\|(A^{\bu}_a-A^{\bv}_a)\otimes I\ket\psi\right\|^2
 \leq8m\eps,
\end{equation}
where $\bu$ and $\bv$ are now independent and uniform.

The inconsistency of $G$ with~$A$ is
\[
 1-\E_{\bu,\bv}\sum_a
 \bra\psi A^{\bu}_a\otimes A^{\bv}_a\ket\psi.
\]
The self-consistency test bounds the same expression with $\bv=\bu$
by~$\delta$.  The difference is at most $\sqrt D$ by Cauchy--Schwarz:
\begin{align*}
 &\left|\E_{\bu,\bv}\sum_a
 \bra\psi A^{\bu}_a\otimes(A^{\bu}_a-A^{\bv}_a)\ket\psi\right|\\
 &\quad\leq
 \left(\E_{\bu}\sum_a\bra\psi(A^{\bu}_a)^2\otimes I\ket\psi\right)^{1/2}
 \left(\E_{\bu,\bv}\sum_a\bra\psi I\otimes
 (A^{\bu}_a-A^{\bv}_a)^2\ket\psi\right)^{1/2}
 =\sqrt D.
\end{align*}
Here the first factor is one because $A^{\bu}$ is projective, and
the second factor equals $\sqrt D$ by symmetry.  Hence
\[
 A^{\bu}_a\otimes I\simeq_{\delta+\sqrt{8m\eps}}I\otimes G_{g_a}.
\]
This is the asserted degree-zero bound.

We may therefore assume $d\geq1$.
Consistency errors are at most one, so the result is immediate if one
of $\eps,\delta,\gamma,d/q$ is at least one.  Hence all four parameters
may now be assumed strictly less than one.
We induct on $m$.  For $m=1$, there is a single axis-parallel line
$\ell=\F_q$.  Taking $G=B^\ell$ gives
\[
 A^u_a\otimes I\simeq_{\eps}I\otimes G_{[g(u)=a]}.
\]
This is enough because $\eps\leq\eps^{1/32}\leq\sigma$.

Suppose the assertion is known in dimension $m$, and consider an
$(\eps,\delta,\gamma)$-good strategy in dimension $m+1$.  Let
$(\eps_x,\delta_x,\gamma_x)$ be the failure probabilities of its
$x$-restricted strategy.  By \Cref{lem:restricted-probabilities},
\[
 \E_x\eps_x\leq\frac{m+1}{m}\eps,
 \qquad \E_x\delta_x\leq\delta,
 \qquad \E_x\gamma_x\leq\frac{m+1}{m}\gamma.
\]
Let $p=1/32$, $r=1/8$, $\theta=d/q$, $c=(m+1)/m$, and
\[
 S=\eps^p+\delta^p+\theta^p+\gamma^r.
\]
The induction hypothesis on each restricted strategy gives a
measurement with error
\[
 \sigma_x=1000d^2m^4
 \bigl(\eps_x^p+\delta_x^p+\theta^p+\gamma_x^r\bigr).
\]
By concavity and \Cref{lem:restricted-probabilities},
\begin{equation}\label{eq:average-induction-error}
 \bar\sigma:=\E_x\sigma_x
 \leq1000d^2m^4c^rS.
\end{equation}

Apply \Cref{thm:self-improvement-in-induction-section} to these
measurements, using a common ancilla and embedding $J$ for all
slices.  Extend the ambient strategy by the identity on this ancilla
and use the symmetric state $\ket{\psi'}=(J\otimes J)\ket\psi$.
By concavity of the square root, the output errors satisfy
\[
 \E_x\zeta_x
 \leq200(m+1)\bigl(\sqrt\eps+\sqrt\delta+\sqrt\theta\bigr)
 \leq\bar\zeta,
 \qquad
 \bar\zeta=3000(m+1)
 \bigl(\eps^{1/4}+\delta^{1/4}+\theta^{1/4}\bigr)
 \leq3000(m+1)S.
\]
Here we retain the larger budget $\bar\zeta$ to simplify the scalar
recurrence; the parameters $\eps,\delta,\theta$ are less than one.
The projective slices on this common enlarged space satisfy all four
hypotheses of \Cref{thm:ld-pasting-in-induction-section}, with
\[
 \kappa=\bar\sigma+\bar\zeta,\qquad\zeta=\bar\zeta.
\]
For boundedness we use permutation invariance to exchange the tensor
factors.  The other three slice hypotheses retain the same tensor
orientation.  The extended ambient strategy has the original test
errors $\eps,\delta,\gamma$.

The pasting error also contains
\[
 s=\eps^{1/8}+\delta^{1/8}+\gamma^{1/8}
   +\bar\zeta^{1/8}+\theta^{1/8}.
\]
Subadditivity and $3000^{1/8}<3$ give
\[
 \bar\zeta^{1/8}
 \leq3(m+1)^{1/8}
 \bigl(\eps^p+\delta^p+\theta^p\bigr),
 \qquad
 s\leq4(m+1)^{1/8}S.
\]
The pasting theorem therefore gives
\begin{align*}
 \sigma_{m+1}
 &\leq c(\bar\sigma+\bar\zeta)
       +21d^2(m+1)^2\sqrt{m}\,s\\
 &\leq1000d^2m^4c^{1+r}S
       +3000c(m+1)S
       +84d^2\sqrt{m}\,(m+1)^{2+1/8}S.
\end{align*}
If $T=\{T_h\}$ is the pasted measurement on the enlarged local space,
set $\widehat T_h=J^\dagger T_hJ$.  These positive operators sum to
$I$, and
\[
 \bra\psi A^u_a\otimes\widehat T_h\ket\psi
 =\bra{\psi'}(A^u_a\otimes I_{\mathcal K})\otimes T_h\ket{\psi'}.
\]
Thus $\widehat T$ has the same consistency bound on the original
space.  Compression is applied only to individual measurement
effects; no multiplicativity or preservation of projectivity is
needed.  The displayed scalar recurrence is therefore the whole
induction bookkeeping.
By concavity, $c^r\leq1+r(c-1)=1+1/(8m)$; also
$(m+1)^r\leq2\sqrt{m}$.  Hence the gap between the desired
coefficient and the inherited term is at least
\[
 1000d^2(m+1)\left(\frac{23}{8}m^2+3m+1\right).
\]
After division by $d^2(m+1)$, the two fresh coefficients are at most
\[
  3000\frac{m+1}{m}+84\sqrt{m}(m+1)^{1+r}
  \leq6000+168m(m+1)
  \leq1000\left(\frac{23}{8}m^2+3m+1\right),\qquad(m\geq1).
\]
They are therefore absorbed by the gap, proving
\[
 \sigma_{m+1}\leq1000d^2(m+1)^4S.
\]
This is exactly the asserted positive-degree bound in dimension $m+1$.
\end{proof}

\section{Expansion in the hypercube graph}
\label{sec:expansion}

\begin{definition}[Hypercube graph]
The \emph{hypercube graph} $C=(V,E)$ has vertex set $V=\F_q^m$,
with an edge between points that disagree in at most one coordinate
(including a self-loop at every vertex).
A random edge $(\bu,\bv)\sim C$ is sampled by drawing
$\bu\sim\F_q^m$, $\bi\sim\{1,\ldots,m\}$, and $\bx\sim\F_q$
independently and setting $\bv=\bu+\bx e_{\bi}$.
\end{definition}

\subsection{Local and global variance}

In this subsection, $\ket\psi$ is a vector, not necessarily normalized,
in $\calH_A\otimes\calH_B$, and $A^u$ is Hermitian for every
$u\in\F_q^m$.

\begin{definition}\label{def:local-and-variance}
The \emph{local} and \emph{global variance} of~$A$ on $\ket\psi$ are
\begin{align*}
 \mathbf{Var}_{\mathrm{local}}(A,\psi)
 &:=\frac12\E_{(\bu,\bv)\sim C}
   \bra\psi(A^{\bu}-A^{\bv})^2\otimes I\ket\psi,\\
 \mathbf{Var}_{\mathrm{global}}(A,\psi)
 &:=\frac12\E_{\bu,\bv\sim\F_q^m}
   \bra\psi(A^{\bu}-A^{\bv})^2\otimes I\ket\psi,
\end{align*}
where $\bu,\bv$ are independent in the second line.
\end{definition}

\begin{lemma}[Product Poincar\'e inequality]\label{lem:local-to-global}
For any Hilbert space $\mathcal{K}$ and function $X:\F_q^m\to\mathcal{K}$,
\[
 \E_{\bu,\bv\sim\F_q^m}\|X(\bu)-X(\bv)\|^2
 \leq m\,\E_{(\bu,\bv)\sim C}\|X(\bu)-X(\bv)\|^2,
\]
where $\bu,\bv$ are independent on the left.  In particular,
\[
 \mathbf{Var}_{\mathrm{global}}(A,\psi)
 \leq m\,\mathbf{Var}_{\mathrm{local}}(A,\psi).
\]
\end{lemma}

\begin{proof}
Equip the space of $\mathcal{K}$-valued functions on $\F_q^m$ with the inner product
$\langle X,Y\rangle=\E_u\langle X(u),Y(u)\rangle$.
For each coordinate $i$, let $E_i$ average that coordinate while
fixing the others.  These are commuting orthogonal projections, and
$E=E_1\cdots E_m$ averages all coordinates.  The telescoping identity
\[
 I-E=\sum_{i=1}^m E_1\cdots E_{i-1}(I-E_i)
 \leq\sum_{i=1}^m(I-E_i)
\]
holds because each summand is a product of commuting projections,
hence is positive and bounded above by $I-E_i$.  Taking its quadratic
form at $X$ gives
\[
 \E_u\|X(u)-\E_vX(v)\|^2
 \leq\sum_{i=1}^m\E_u\|X(u)-(E_iX)(u)\|^2.
\]
The left side is half the mean squared difference between $X$ at
two independent uniform points.  Conditioned on all coordinates
except $i$, the $i$-th term on the right is half the mean squared
difference between two independent choices of coordinate $i$.
Their sum is therefore $m/2$ times the mean squared difference
along a random edge of $C$.  This proves the Hilbert-space-valued
inequality.  Taking $X(u)=(A^u\otimes I)\ket\psi$ gives the operator
formulation, without requiring $\ket\psi$ to be normalized.
\end{proof}

\section{Global variance of the points measurements}
\label{sec:variance}
Throughout this section, $(\psi, A, B, L)$ will denote a fixed $(\eps, \delta, \gamma)$-good symmetric strategy for the $(m,q,d)$-low individual degree test.

\begin{lemma}\label{lem:generalize-b}
Let $G \in \polysub{m}{q}{d}$. Then
\begin{equation*}
B^\ell_{[f(u) = g(u)]} \otimes (G_g)^{1/2} \approx_{d/q} B^\ell_{g|_\ell} \otimes (G_g)^{1/2}
\end{equation*}
on the axis-parallel lines test distribution.
\end{lemma}
\begin{proof}
On an axis-parallel line, both $f$ and $g|_\ell$ have univariate
degree at most~$d$.  Thus, when they are distinct, they agree at a
uniform point of~$\ell$ with probability at most $d/q$.
We now bound the squared distance:
\begin{align*}
&~\E_{\bu, \bell} \sum_{g \in \polyfunc{m}{q}{d}} \Vert (B^{\bell}_{[f(\bu) = g(\bu)]}  - B^{\bell}_{g|_{\bell}}) \otimes (G_g)^{1/2} \ket{\psi}\Vert^2\\
=&~\E_{\bu, \bell} \sum_{g \in \polyfunc{m}{q}{d}} \bra{\psi} \bigg(\sum_{f : f \neq g|_{\bell}} \bone[f(\bu) = g(\bu)] \cdot B^{\bell}_f\bigg)^2 \otimes G_g \ket{\psi}\\
\leq&~\E_{\bu, \bell} \sum_{g \in \polyfunc{m}{q}{d}} \bra{\psi} \bigg(\sum_{f : f \neq g|_{\bell}} \bone[f(\bu) = g(\bu)] \cdot B^{\bell}_f\bigg) \otimes G_g \ket{\psi} \\
= &~\E_{\bell} \sum_{g \in \polyfunc{m}{q}{d}} \sum_{f : f \neq g|_{\bell}} \bra{\psi}  B^{\bell}_f \otimes G_g \ket{\psi} \cdot \left( \E_{\bu} \bone[f(\bu) = g(\bu)]\right)\\
\leq &~\E_{\bell} \sum_{g \in \polyfunc{m}{q}{d}} \sum_{f : f \neq g|_{\bell}} \bra{\psi}  B^{\bell}_f \otimes G_g \ket{\psi} \cdot \frac{d}{q}\tag{by Schwartz-Zippel}\\
\leq &~\frac{d}{q}.\qedhere
\end{align*}
\end{proof}

\begin{lemma}\label{lem:local-variance-of-points}
Let $G \in \polysub{m}{q}{d}$.  Then
\begin{equation}\label{eq:local-variance-of-points-equation}
A^u_{g(u)} \otimes (G_g)^{1/2} \approx_{24\cdot(\eps + \delta + \frac{d}{q})} A^v_{g(v)} \otimes (G_g)^{1/2}
\end{equation}
on the distribution $(\bu, \bv) \sim C$.
\end{lemma}
\begin{proof}
Let $\bu$ and $\bell$ be distributed as in the axis-parallel lines test,
and sample $\bv \sim \bell$.
Then $\bv$ and $\bell$ are also distributed as in the axis-parallel lines test.
As a result,
\begin{align*}
A^u_{g(u)} \otimes (G_g)^{1/2}
&\approx_{2\delta} I \otimes (G_g)^{1/2} A^u_{g(u)}\tag{by \Cref{prop:simeq-to-approx}}\\
&\approx_{2\eps} B^\ell_{[f(u) = g(u)]} \otimes (G_g)^{1/2}\tag{by \Cref{prop:simeq-to-approx}}\\
&\approx_{\frac{d}{q}} B^\ell_{g|_{\ell}} \otimes (G_g)^{1/2}\tag{by \Cref{lem:generalize-b}}\\
&\approx_{\frac{d}{q}} B^\ell_{[f(v) = g(v)]} \otimes (G_g)^{1/2}\tag{by \Cref{lem:generalize-b}}\\
&\approx_{2\eps} I \otimes (G_g)^{1/2} A^v_{g(v)} \tag{by \Cref{prop:simeq-to-approx}}\\
&\approx_{2\delta} A^v_{g(v)} \otimes (G_g)^{1/2}.\tag{by \Cref{prop:simeq-to-approx}}
\end{align*}
Steps 1, 2, 5, and 6 also use
\Cref{rem:good-strat-characterization,prop:cab-approx-delta}; in
steps 2 and 5, \Cref{prop:approx-swap-invariance} first puts the
point--line approximation on the required tensor factor.
The lemma now follows from \Cref{prop:triangle-inequality-for-approx_delta}.
\end{proof}

We note that \Cref{eq:local-variance-of-points-equation} is equivalent to the statement that
\begin{equation}\label{eq:equivalent-local-variance}
\sum_{g \in \polyfunc{m}{q}{d}}\E_{(\bu, \bv) \sim C} \bra{\psi} (A^{\bu}_{g(\bu)}  - A^{\bv}_{g(\bv)})^2 \otimes G_g \ket{\psi}
\leq 24\left(\eps + \delta + \frac{d}{q}\right).
\end{equation}
This can be viewed as a form of local variance for the points measurements.
We now derive the corresponding expression for the global variance of the points measurements.

\begin{lemma}\label{lem:global-variance-of-points}
Let $G \in \polysub{m}{q}{d}$.  Then
\begin{equation}\label{eq:global-variance-of-points-equation}
A^u_{g(u)} \otimes (G_g)^{1/2} \approx_{24m\cdot(\eps + \delta + \frac{d}{q})} A^v_{g(v)} \otimes (G_g)^{1/2}
\end{equation}
on the distribution $\bu, \bv \sim \F_q^m$.
\end{lemma}
\begin{proof}
Apply \Cref{lem:local-to-global} to the single Hilbert-space-valued
function
\[
 X(u)=\bigoplus_{g\in\polyfunc{m}{q}{d}}
       (A^u_{g(u)}\otimes\sqrt{G_g})\ket\psi.
\]
The direct-sum norm adds the squared norms of its components, so
\begin{equation}\label{eq:global-variance-target}
 \begin{aligned}
 &\E_{\bu,\bv}\sum_g\bra\psi
       (A^{\bu}_{g(\bu)}-A^{\bv}_{g(\bv)})^2\otimes G_g\ket\psi\\
 &\qquad=\E_{\bu,\bv}\|X(\bu)-X(\bv)\|^2
 \leq m\,\E_{(\bu,\bv)\sim C}\|X(\bu)-X(\bv)\|^2\\
 &\qquad\leq24m\left(\eps+\delta+\frac dq\right),
 \end{aligned}
\end{equation}
where the unqualified points $\bu,\bv$ are independent and uniform,
and the last inequality is \eqref{eq:equivalent-local-variance}.
\end{proof}

\section{Self-improvement}
\label{sec:self-improvement}

Throughout this section, $(\psi,A,B,L)$ is an
$(\eps,\delta,\gamma)$-good symmetric strategy for the $(m,q,d)$ low
individual degree test.  We first construct a sub-measurement on the
original local space.  Its SDP certificate satisfies the
operator inequality $Z^2\leq H\leq I$, which supplies strong
self-consistency directly.  We then dilate any finite family of these
sub-measurements simultaneously to projective sub-measurements on one
common enlarged space.  After pasting, compression returns a
measurement to the original space.

For the non-projective output we state strong self-consistency as a
defect in the diagonal overlap, following
\Cref{def:strong-self-consistency}.  Upon dilation, the squared mirror
error is exactly twice this defect; these two error conventions will
be kept separate.

\begin{lemma}[Self-improvement with non-projective output]\label{lem:self-improvement-helper}
Let $G \in \polymeas{m}{q}{d}$ be a measurement with the following property:
\begin{enumerate}
  \item (Consistency with~$A$): \label{item:self-improvement-G-consistency} On average over $\bu \sim \F_q^{m}$,
	  \begin{equation*}
	  A^{u}_a \otimes I \simeq_{\nu} I \otimes G_{[g(u)=a]}.
	  \end{equation*}
	  \end{enumerate}
Let
\begin{equation*}
\zeta_0 = 100m\cdot \Big(\eps^{1/2} + \delta^{1/2} + (d/q)^{1/2}\Big).
\end{equation*}
Then there exists $H \in \polysub{m}{q}{d}$ with the following properties:
\begin{enumerate}
    \item(Completeness):  \label{item:self-improvement-completeness} If $H =  \sum_h H_h$, then
    \begin{equation*}
    \bra{\psi} H \otimes I \ket{\psi} \geq (1-\nu)-\zeta_0.
    \end{equation*}
    \item(Consistency with~$A$):\label{item:self-improvement-A-consistency} On average over $\bu \sim \F_q^m$,
        \begin{equation*}
            A^u_a \otimes I \simeq_{\zeta_0} I \otimes H_{[h(u) = a]}.
        \end{equation*}
    \item(Strong self-consistency): \label{item:self-improvement-self}
	\begin{equation*}
		\sum_{h} \bra{\psi} H_h \ot H_h \ket{\psi} \geq \bra{\psi} H \ot I \ket{\psi} - \zeta_0.
	\end{equation*}
    \item(Boundedness): \label{item:self-improvement-boundedness} There exists a positive-semidefinite matrix~$Z\leq I$ such that
    \begin{equation*}
        \bra{\psi} Z \otimes I \ket{\psi} -\E_{\bu} \sum_a \bra{\psi}  A^{\bu}_{a} \otimes H_{[h(\bu)=a]} \ket{\psi} \leq \zeta_0
    \end{equation*}
    and for each $h \in \polyfunc{m}{q}{d}$,
	\begin{equation*}
	Z \geq \left(\E_{\bu} A^{\bu}_{h(\bu)}\right).
	\end{equation*}
\end{enumerate}
\end{lemma}

\subsection{A semidefinite program for measurements}

To prove \Cref{lem:self-improvement-helper}, we use the primal--dual
semidefinite program from the original proof.  Introduce the notation
\begin{equation*}
A_g = \E_{\bu \sim \F_q^m} A^{\bu}_{g(\bu)}.
\end{equation*}
The primal and dual programs are
\begin{align}
 \sup_{T_g\geq0,\,\sum_gT_g\leq I}
   &\quad\sum_g\Tr(T_gA_g),\label{eq:primal-objective}\\
 \inf_{Z\geq A_g\,\forall g}
   &\quad\Tr(Z).\label{eq:dual-objective}
\end{align}

\begin{lemma}[SDP certificate]\label{lem:sdp}
The programs \eqref{eq:primal-objective} and
\eqref{eq:dual-objective} have the same attained optimum.  They admit
optimal solutions $T=\{T_g\}$ and $Z$ for which $T$ is a measurement,
$Z\geq A_g$ for every $g$, and
\begin{equation}\label{eq:complementary-slackness}
 T_g(Z-A_g)=0\qquad\text{for every }g\in\polyfunc{m}{q}{d}.
\end{equation}
In particular, $Z\geq0$.
\end{lemma}

\begin{proof}
The block-diagonal canonical SDP has one block for each polynomial and
one slack block for $I-\sum_gT_g$.  Its primal constraint is that the
sum of these blocks is $I$, and its dual constraints are $Z\geq A_g$
for each polynomial block and $Z\geq0$ for the slack block.  Thus its
objective values are exactly those of the displayed programs.
Both sides are strictly feasible: for $M=|\polyfunc{m}{q}{d}|$, take
$T_g=I/(2M)$ with slack $I/2$, and take $Z=2I$ since $0\leq A_g\leq I$.
Finite-dimensional SDP strong duality gives attained optima and
complementary slackness.  The polynomial blocks of that relation give
$T_g(Z-A_g)=0$.

We may take the optimal polynomial family to be complete.  Indeed,
assign its unused positive operator $R=I-\sum_gT_g$ to any fixed
polynomial outcome.  Since $A_g\geq0$, this cannot decrease the primal
objective, and optimality implies that it leaves the objective equal
to the optimum.  This produces a measurement; equality of primal and
dual values again gives the displayed complementary-slackness
relations.  The dual constraints give $Z\geq A_g\geq0$.
\end{proof}

\subsection{The filtered measurement and its certificate}

\begin{proof}[Proof of \Cref{lem:self-improvement-helper}]
Write $\langle X\rangle=\bra\psi X\ket\psi$ and
$\omega(M)=\langle M\otimes I\rangle$.  Apply
\Cref{lem:sdp}, and put
\[
 P_g^u=A^u_{g(u)},\qquad
 H_g=\E_u P_g^uT_gP_g^u,\qquad H=\sum_gH_g,
 \qquad\eta=24m(\eps+\delta+d/q).
\]
Grouping $g$ by its evaluation at $u$ gives
\[
 \sum_gP_g^uT_gP_g^u
 =\sum_aA_a^u\Bigl(\sum_{g:g(u)=a}T_g\Bigr)A_a^u
 \leq\sum_aA_a^u=I.
\]
Thus $H$ is a sub-measurement.  Complementarity and its adjoint imply
\[
 Z^2=Z\Bigl(\sum_gT_g\Bigr)Z
     =\sum_gA_gT_gA_g,
 \qquad
 H-Z^2=\sum_g\E_u(P_g^u-A_g)T_g(P_g^u-A_g)\geq0.
\]
Consequently
\begin{equation}\label{eq:self-improvement-certificate-contraction}
 Z^2\leq H\leq I,\qquad 0\leq Z\leq I.
\end{equation}
The only variance input is \Cref{lem:global-variance-of-points}, in
its equivalent form
\begin{equation}\label{eq:self-improvement-variance-input}
 \E_{u,v}\sum_g\langle(P_g^u-P_g^v)^2\otimes T_g\rangle
 \leq\eta.
\end{equation}
All direct sums below include the probability weights of the displayed
questions.

\medskip\noindent\emph{Consistency.}
Let $c=(\sqrt{2\delta}+\sqrt\eta)^2$ and, in the direct sum over
$(u,v,g)$, set
\[
 x_{u,v,g}=((I-P_g^u)\otimes\sqrt{T_g}P_g^v)\ket\psi,
 \qquad
 z_{u,v,g}=((I-P_g^u)P_g^v\otimes\sqrt{T_g})\ket\psi.
\]
The inconsistency with the point measurements is
\[
 \E_u\sum_g\langle(I-P_g^u)\otimes H_g\rangle=\|x\|^2.
\]
For $D_a^v=A_a^v\otimes I-I\otimes A_a^v$, grouping by $g(v)=a$
gives
\begin{align*}
 \|x-z\|^2
 &=\E_{u,v}\sum_a\left\langle D_a^v
     \Bigl(\sum_{g:g(v)=a}(I-P_g^u)\otimes T_g\Bigr)D_a^v\right\rangle\\
 &\leq\E_v\sum_a\langle(D_a^v)^2\rangle\leq2\delta.
\end{align*}
Here each parenthesized operator is at most $I$, and the last step
uses point self-consistency and projectivity.  Since
$(I-P_g^u)P_g^v=(I-P_g^u)(P_g^v-P_g^u)$,
\Cref{eq:self-improvement-variance-input} gives $\|z\|^2\leq\eta$.
The triangle inequality therefore proves
\begin{equation}\label{eq:self-improvement-short-consistency}
 \E_u\sum_g\langle(I-P_g^u)\otimes H_g\rangle\leq c.
\end{equation}

\medskip\noindent\emph{Strong self-consistency.}
Put $h=\omega(H)$.  By $Z\geq A_g$, symmetry, and
\Cref{eq:self-improvement-short-consistency},
\[
 \langle H\otimes Z\rangle
 \geq\sum_g\langle H_g\otimes A_g\rangle\geq h-c.
\]
Moreover, \Cref{eq:self-improvement-certificate-contraction} gives
$H\geq Z^2\geq2Z-I$.  Hence
\begin{equation}\label{eq:self-improvement-total-mass}
 \langle H\otimes H\rangle\geq h-2c.
\end{equation}

Apply \Cref{lem:projective-bridge} to $P=A^u$ and
$C_a^u=\sum_{g:g(u)=a}H_g$, using
\Cref{eq:self-improvement-short-consistency}.
Let
\[
 \theta=\min\{1,md/q\},\qquad
 F=\sum_{g\ne g'}\langle H_g\otimes H_{g'}\rangle.
\]
By Schwartz--Zippel, distinct polynomials of individual degree at
most $d$ agree at a random point with probability at most $\theta$.  Splitting $F$ by agreement of the
point evaluations gives
\[
 F\leq(2\sqrt c+\sqrt{2\delta})^2+\theta F
   \leq(2\sqrt c+\sqrt{2\delta})^2+\theta,
\]
since $F\leq\langle H\otimes H\rangle\leq1$.  Subtracting this
from \Cref{eq:self-improvement-total-mass} yields
\begin{equation}\label{eq:self-improvement-short-self-consistency}
 \sum_g\langle H_g\otimes H_g\rangle\geq h-s,
 \qquad
 s=2c+(2\sqrt c+\sqrt{2\delta})^2+\theta.
\end{equation}

\medskip\noindent\emph{Completeness and boundedness.}
Since the seed $G$ is a measurement and has consistency error at
most $\nu$,
\begin{equation}\label{eq:self-improvement-seed-certificate}
 \omega(Z)=\sum_g\langle Z\otimes G_g\rangle
 \geq\sum_g\langle A_g\otimes G_g\rangle\geq1-\nu.
\end{equation}
Write $T_a^u=\sum_{g:g(u)=a}T_g$ and form the direct-sum vectors
\[
 t=(\sqrt{T_a^u}\otimes I)\ket\psi,\qquad
 x=(\sqrt{T_a^u}A_a^u\otimes I)\ket\psi,\qquad
 y=(\sqrt{T_a^u}\otimes A_a^u)\ket\psi.
\]
They satisfy
\[
 \|t\|=1,\quad \operatorname{Re}\langle t,x\rangle=\omega(Z),
 \quad\|x\|^2=h,\quad\langle t,y\rangle=\|y\|^2,
 \quad\|x-y\|\leq\sqrt{2\delta}.
\]
The second identity follows from $\sum_gT_gA_g=Z$, and the last
inequality follows by inserting $T_a^u\otimes I\leq I$ in each
point mirror error.  In particular, $\|t-2y\|=1$, and expansion
around $y$ gives
\begin{align*}
 \omega(Z)-h
 &=\operatorname{Re}\langle t-2y,x-y\rangle-\|x-y\|^2\\
 &\leq\sqrt{2\delta}.
\end{align*}
Together with \Cref{eq:self-improvement-seed-certificate}, this proves
\begin{equation}\label{eq:self-improvement-short-completeness}
 h\geq1-\nu-\sqrt{2\delta}.
\end{equation}
Also, by \Cref{eq:self-improvement-short-consistency},
\begin{equation}\label{eq:self-improvement-short-boundedness}
 \omega(Z)-\sum_g\langle A_g\otimes H_g\rangle
 \leq\omega(Z)-h+c\leq\sqrt{2\delta}+c.
\end{equation}
This is the helper's boundedness condition.  Since $Z\geq A_g$, it
also implies
\[
 \langle Z\otimes(I-H)\rangle
 \leq\omega(Z)-\sum_g\langle A_g\otimes H_g\rangle
 \leq\sqrt{2\delta}+c.
\]

\medskip\noindent\emph{The uniform error budget.}
The explicit estimates above imply
\[
 c\leq52m(\eps+\delta+d/q),\qquad
 s\leq525m(\eps+\delta+d/q).
\]
Indeed $c\leq2\eta+4\delta$, and
$s\leq10c+4\delta+\theta\leq20\eta+44\delta+md/q$.
Put $E=\sqrt\eps+\sqrt\delta+\sqrt{d/q}$, so
$\zeta_0=100mE$.  If $\zeta_0\geq1$, take instead $H_g=0$ for all
$g$ and $Z=I$; all four conclusions hold.  Otherwise
$E<1/(100m)$, and
\[
 c\leq52mE^2\leq\zeta_0,\qquad
 s\leq525mE^2\leq5.25E\leq\zeta_0,
 \qquad
 \sqrt{2\delta}+c\leq(\sqrt2+0.52)E\leq\zeta_0.
\]
The completeness loss is at most $\sqrt{2\delta}\leq\zeta_0$ as
well.  This proves the lemma.
\end{proof}

\subsection{Simultaneous dilation and compression}
\label{sec:self-improvement-dilation}

\begin{lemma}[Simultaneous dilation and compression]
\label{lem:simultaneous-dilation-compression}
Let $\{H^x\}_x$ be a finite family of sub-measurements on a
finite-dimensional space $\calH$, with one common finite outcome
set.  There are a finite ancillary space $\mathcal{K}$, an isometry
$J:\calH\longrightarrow\calH\otimes\mathcal{K}$ of the form
$Jv=v\otimes\ket0$, and projective sub-measurements $P^x$ such that
\[
 J^\dagger P_g^xJ=H_g^x.
\]
For a symmetric state $\ket\psi$, put
$\ket{\psi'}=(J\otimes J)\ket\psi$, and extend any original operator
$R$ to $R'=R\otimes I_{\mathcal{K}}$.  Slice masses, consistency with
original measurements, and correlations involving one original
operator and one slice effect are preserved.  Moreover,
\begin{equation}\label{eq:dilation-mirror-defect}
 \begin{split}
 &\sum_g\|(P_g^x\otimes I-I\otimes P_g^x)\ket{\psi'}\|^2\\
 &\qquad=2\left(\bra\psi H^x\otimes I\ket\psi
       -\sum_g\bra\psi H_g^x\otimes H_g^x\ket\psi\right),
 \qquad H^x=\sum_gH_g^x.
 \end{split}
\end{equation}
Finally, if $T=\{T_h\}$ is a measurement on the enlarged local
space, $\widehat T_h=J^\dagger T_hJ$ is a measurement on $\calH$
with the same consistency with the original point measurements.
\end{lemma}

\begin{proof}
Complete each sub-measurement to a measurement by
$M_g^x=H_g^x$ and $M_\bot^x=I-H^x$.  Take $\mathcal{K}$ with an
orthonormal basis indexed by these outcomes.  For each $x$,
\[
 V_xv=\sum_a\sqrt{M_a^x}\,v\otimes\ket a
\]
is an isometry.  The isometry from $J\calH$ to $V_x\calH$ given
by $Jv\mapsto V_xv$ extends to a unitary $U_x$ on
$\calH\otimes\mathcal{K}$.  Define
\[
 P_a^x=U_x^\dagger(I\otimes\ket a\!\bra a)U_x.
\]
These form a projective measurement, and
$J^\dagger P_a^xJ=M_a^x$.  Discard the failure outcome.

For original operators $R$ and slice outcomes,
\[
 \bra{\psi'}R'\otimes P_g^x\ket{\psi'}
 =\bra\psi R\otimes H_g^x\ket\psi,
 \qquad
 \bra{\psi'}P_g^x\otimes P_h^y\ket{\psi'}
 =\bra\psi H_g^x\otimes H_h^y\ket\psi.
\]
These identities prove the correlation claims.  Expanding the
squared mirror error and using projectivity and symmetry proves
\Cref{eq:dilation-mirror-defect}.  Domination inequalities are
preserved by extension: for example, $Z^x\geq\E_uA_{g(u)}^{u,x}$
implies $Z^{x\prime}\geq\E_uA_{g(u)}^{u,x\prime}$.  Boundedness
is preserved exactly, since
\[
 \bra{\psi'}Z^{x\prime}\otimes\Bigl(I-\sum_gP_g^x\Bigr)\ket{\psi'}
 =\bra\psi Z^x\otimes(I-H^x)\ket\psi.
\]
Finally, compression preserves positivity and
$\sum_h\widehat T_h=J^\dagger I J=I$, while
\[
 \bra{\psi'}A_a^{u\prime}\otimes T_h\ket{\psi'}
 =\bra\psi A_a^u\otimes\widehat T_h\ket\psi.
\]
This proves the last assertion.  No preservation of products or of
projectivity under compression is asserted or required.
\end{proof}

\begin{proof}[Proof of \Cref{thm:self-improvement-in-induction-section}]
Apply \Cref{lem:self-improvement-helper} on the original space and
then \Cref{lem:simultaneous-dilation-compression}.  Completeness,
point consistency, and boundedness are preserved with error
$\zeta_0$, whereas \Cref{eq:dilation-mirror-defect} gives squared
mirror error at most $2\zeta_0$.  All four conclusions therefore
hold with
\[
 \zeta=2\zeta_0=200m(\sqrt\eps+\sqrt\delta+\sqrt{d/q}).
\]
The extended strategy has the same tested correlations and remains
symmetric.  For a family of restricted strategies, apply the helper
to every slice first and then use the single common dilation of
\Cref{lem:simultaneous-dilation-compression}.  All bounds hold
slice by slice, so they can be averaged before applying pasting.
The pasted measurement can then be compressed to the original
local space.
\end{proof}

\section{Commutativity of the points measurements}
\label{sec:commutativity-points}

\begin{theorem}\label{thm:commutativity-points}
Let $(\psi, A, B, L)$ be an $(\eps, \delta, \gamma)$-good symmetric strategy for the
$(m,q,d)$ low individual degree test.
On average over independent and uniformly random $\bu, \bv \sim \F_q^m$,
\begin{equation*}
(A^{u}_a \cdot A^{v}_b) \otimes I \approx_{32\gamma m} (A^{v}_b\cdot A^{u}_a) \otimes I.
\end{equation*}
\end{theorem}
\begin{proof}
The strategy passes the diagonal lines test with probability at least $1-\gamma$.
Therefore, it passes the $m$-restricted diagonal lines test with probability at least $1- \gamma\cdot m$.
This means that
\begin{equation*}
A^{u}_a \otimes I \simeq_{\gamma \cdot m} I \otimes L^{\ell}_{[f(u)=a]},
\end{equation*}
on average over independent uniform $\bu,\bw\in\F_q^m$, with
$\bell=\{\bu+t\bw:t\in\F_q\}$.  This includes the degenerate line
$\{\bu\}$ when $\bw=0$.
By~\Cref{prop:simeq-to-approx}, this implies that
\begin{equation}\label{eq:point-diagonal-line-approx}
A^{u}_a \otimes I \approx_{2\cdot \gamma \cdot m} I \otimes L^{\ell}_{[f(u)=a]}.
\end{equation}

Let $\bu$ and $\bv$ be independent and uniformly random points in $\F_q^m$.
Set $\bell=\{\bu+t(\bv-\bu):t\in\F_q\}$; in particular,
$\bell=\{\bu\}$ when $\bu=\bv$.  Conditional on~$\bu$, the
direction $\bv-\bu$ is uniform in~$\F_q^m$, so the marginal on
$(\bu,\bell)$ is exactly the one in
\Cref{eq:point-diagonal-line-approx}.  The same holds for
$(\bv,\bell)$, using the uniform direction $\bu-\bv$.
As a result,
\begin{align*}
(A^{u}_a \cdot A^{v}_b) \otimes I
&\approx_{2\cdot \gamma \cdot m} A^{u}_a \otimes L^{\ell}_{[f(v)=b]} \tag{by \Cref{eq:point-diagonal-line-approx}}\\
&\approx_{2\cdot \gamma \cdot m} I \otimes (L^{\ell}_{[f(v)=b]} \cdot L^{\ell}_{[f'(u)=a]}) \tag{by \Cref{eq:point-diagonal-line-approx}}\\
&= I \otimes (L^{\ell}_{[f'(u)=a]} \cdot L^{\ell}_{[f(v)=b]}) \tag{because~$L$ is projective}\\
&\approx_{2\cdot \gamma \cdot m} A^{v}_b \otimes L^{\ell}_{[f'(u)=a]} \tag{by \Cref{eq:point-diagonal-line-approx}}\\
&\approx_{2\cdot \gamma \cdot m} (A^{v}_b \cdot A^{u}_a) \otimes I. \tag{by \Cref{eq:point-diagonal-line-approx}}
\end{align*}
Each approximation also uses \Cref{prop:cab-approx-delta} to insert
the displayed relation next to the remaining contraction factors;
its normalization hypotheses follow from projectivity of~$A$ and
$L$.
The theorem now follows from
\Cref{prop:triangle-inequality-for-approx_delta}.
\end{proof}

\section{Commutativity}
\label{sec:g-comm}

Let $(\bu, \bx)$ and $(\bv, \by)$ be sampled independently and uniformly at random from $\F_q^{m+1}$.
In \Cref{sec:G-commutes-after-evaluation},
we will show that the $G$ measurements approximately commute ``after evaluation";
namely, that $G^{\bx}_{[g(\bu)=a]}$ commutes with $G^{\by}_{[g(\bv)=b]}$.
Then, in \Cref{sec:G-commutes},
we will use this to show that $G^{\bx}_g$ approximately commutes with $G^{\by}_h$.
This is necessary if we wish to ``paste" together several~$G^x$ measurements at different points~$x \in \F_q$
to produce a single global measurement~$H \in \polysub{m+1}{q}{d}$, as we do in \Cref{sec:ld-pasting} below.

\subsection{\texorpdfstring{Commutativity of $G$ after evaluation}{Commutativity after evaluation}}\label{sec:G-commutes-after-evaluation}

\begin{lemma}[Boundedness controls rejection]\label{lem:boundedness-rejection}
Let $F$ be a projection and $\{R_i\}$ a sub-measurement on the
first local space, and let $Z\geq0$ and $0\leq\overline A_i\leq Z$
act on the second local space.  Then
\begin{equation}\label{eq:boundedness-rejection-operator}
 \sum_iFR_iF\otimes\overline A_i\leq F\otimes Z.
\end{equation}
In particular, suppose $F^x$ and $\{R_i^x\}$ are such families
indexed by a slice~$x$, independent of the evaluation question~$u$,
and that $A_i^{u,x}\geq0$ satisfies
$\E_uA_i^{u,x}\leq Z^x$ for every $x,i$.  If
$\E_x\bra\psi F^x\otimes Z^x\ket\psi\leq\zeta$, then
\begin{equation}\label{eq:boundedness-rejection-average}
 \E_{x,u}\sum_i
 \bra\psi F^xR_i^xF^x\otimes A_i^{u,x}\ket\psi\leq\zeta.
\end{equation}
\end{lemma}
\begin{proof}
Positivity and $\sum_iR_i\leq I$ give
\[
 \sum_iFR_iF\otimes\overline A_i
 \leq F\Bigl(\sum_iR_i\Bigr)F\otimes Z
 \leq F^2\otimes Z=F\otimes Z.
\]
For the averaged conclusion, first average over~$u$, using the
stated independence, then apply this operator inequality for each~$x$
and take the state expectation and average over~$x$.
\end{proof}

\begin{lemma}[Commutativity of~$G$ after evaluation]
  \label{lem:comm-data-processed-g}
  Let $m\geq1$, and let $(\psi, A, B, L)$ be an $(\eps, \delta, \gamma)$-good symmetric strategy for the $(m+1,q,d)$ low individual degree test.
  Let $\{G^x\} \in \polysub{m}{q}{d}$ be a collection of projective sub-measurements indexed by $x \in \F_q$ with the following properties:
  \begin{enumerate}
  \item (Consistency with~$A$): \label{item:data-processed-consistency} On average over $(\bu, \bx) \sim \F_q^{m+1}$,
	  \begin{equation*}
	  A^{u, x}_a \otimes I \simeq_{\zeta} I \otimes G^x_{[g(u)=a]}.
	  \end{equation*}
  \item (Strong self-consistency): \label{item:data-processed-self-consistency} On average over~$\bx \sim \F_q$,
	\begin{equation*}
		G^x_g \otimes I \approx_{\zeta} I \otimes G^x_g.
	\end{equation*}
  \item (Boundedness): \label{item:data-processed-boundedness} There exists a positive-semidefinite matrix $Z^x$ for each $x \in \F_q$ such that
	\begin{equation*}
		\E_{\bx} \bra{\psi} (I-G^{\bx})\otimes Z^{\bx} \ket{\psi} \leq \zeta
	\end{equation*}
	and for each $x \in \F_q$ and $g \in \polyfunc{m}{q}{d}$,
	\begin{equation*}
	 Z^x \geq \left(\E_{\bu} A^{\bu, x}_{g(\bu)}\right).
	\end{equation*}
  \end{enumerate}
  Let
  \begin{equation*}
  \nu = 26 m \cdot (\gamma^{1/2} + \zeta^{1/2}).
  \end{equation*}
  Then on average over independent and uniformly random $(\bu,\bx), (\bv,\by) \sim \F_q^{m+1}$,
  \begin{equation*}
    G^x_{[g(u) = a]} G^y_{[h(v) = b]} \ot I \approx_{\nu} G^y_{[h(v) = b]}
    G^x_{[g(u) = a]} \ot I.
  \end{equation*}
\end{lemma}

\begin{proof}
For fixed $u,v,x,y$, abbreviate
\[
 P_g=G^x_g,\quad Q_h=G^y_h,\quad
 P_a=G^x_{[g(u)=a]},\quad Q_b=G^y_{[h(v)=b]},
 \quad P=G^x,\quad Q=G^y,
\]
and put $A_a=A^{u,x}_a$ and $B_b=A^{v,y}_b$.  We write
\[
 \mathcal E(X_{a,b})
 =\E_{\bu,\bv,\bx,\by}\sum_{a,b}
   \bra\psi X_{a,b}\ket\psi.
\]
All approximation relations below are understood with respect to the
averages displayed in their applications.

Permutation invariance reads the consistency hypothesis in reverse.
Since $\{P_a\}_a$ and $\{Q_b\}_b$ are projective
sub-measurements with totals $P$ and $Q$, respectively,
\Cref{prop:cons-sub-meas} gives
\begin{equation}\label{eq:evaluated-g-consistency-processed}
 P_a\otimes I\approx_{4\zeta}P\otimes A_a,
 \qquad
 Q_b\otimes I\approx_{4\zeta}Q\otimes B_b.
\end{equation}
Moreover,
\Cref{item:data-processed-self-consistency,%
prop:two-notions-of-self-consistency,%
prop:two-notions-of-self-consistency-after-evaluation} imply
\begin{equation}\label{eq:evaluated-g-self-consistency-processed}
 P_a\otimes I\approx_\zeta I\otimes P_a,
 \qquad
 Q_b\otimes I\approx_\zeta I\otimes Q_b.
\end{equation}
Finally, tensor-swap invariance and
\Cref{thm:commutativity-points} give, in either direction,
\begin{equation}\label{eq:evaluated-point-commutativity}
 I\otimes A_aB_b
 \approx_{32\gamma(m+1)}I\otimes B_bA_a.
\end{equation}
Set $\kappa=\sqrt{32\gamma(m+1)}$.

We apply \Cref{lem:boundedness-rejection} to the two rejection
terms below.  First, for fixed $y,h$, let
\[
 R_h^y=\E_{\bu,\bx}\sum_a P_aG_h^yP_a.
\]
This is positive, independent of~$v$, and satisfies
\[
 \sum_hR_h^y
 =\E_{\bu,\bx}\sum_aP_aQP_a
 \leq\E_{\bu,\bx}\sum_aP_a\leq I.
\]
Expanding $Q_b$ according to $h(v)=b$ and using projectivity of
the point measurements, Cauchy--Schwarz with
factors $\sqrt{R_h^y}\otimes I$ and
$\sqrt{R_h^y}(I-Q)\otimes A^{v,y}_{h(v)}$ gives
\begin{align}
&\left|\mathcal E\bigl(P_aQ_bP_a(I-Q)\otimes B_b\bigr)\right|^2\nonumber\\
&\quad\leq
 \E_{\bv,\by}\sum_h
   \bra\psi(I-Q)R_h^y(I-Q)\otimes A^{v,y}_{h(v)}\ket\psi
 \leq\zeta.\label{eq:evaluated-g-first-leakage}
\end{align}
The omitted squared norm is at most one.  The last inequality is
\eqref{eq:boundedness-rejection-average} with failure projection
$I-G^y$, evaluation question~$v$, and the averaged domination and
slice boundedness in \Cref{item:data-processed-boundedness}.

For the second term, set
\[
 \widetilde R_g^x=\E_{\bv,\by}\sum_bQ_bP_gQ_b.
\]
These operators are positive and independent of~$u$, and
$\sum_g\widetilde R_g^x
=\E_{\bv,\by}\sum_bQ_bPQ_b\leq\E_{\by}G^y\leq I$.
Expanding $P_a$ according to $g(u)=a$, Cauchy--Schwarz with factors
$\sqrt{P_g}\otimes B_b$ and
$\sqrt{P_g}Q_b(I-P)\otimes A^{u,x}_{g(u)}$ gives
\begin{align}
&\left|\mathcal E\bigl(P_aQ_b(I-P)\otimes B_bA_a\bigr)\right|^2\nonumber\\
&\quad\leq
 \E_{\bu,\bx}\sum_g
   \bra\psi(I-P)\widetilde R_g^x(I-P)
   \otimes A^{u,x}_{g(u)}\ket\psi
 \leq\zeta.\label{eq:evaluated-g-second-leakage}
\end{align}
Here the omitted squared norm is at most
$\E_{\bu,\bv,\bx,\by}\sum_{g,b}
\bra\psi P_g\otimes B_b\ket\psi\leq1$.
Apply \eqref{eq:boundedness-rejection-average} with failure projection
$I-G^x$, evaluation question~$u$, and
\Cref{item:data-processed-boundedness} to obtain the last inequality.

We next verify once and for all the normalizations needed to apply
\Cref{prop:closeness-of-ip}.  With
$D=\sum_bQ_b\otimes B_b\leq I$, projectivity and
sub-measurement normalization give
\begin{gather}
 \sum_b\left(\sum_aP_aQ_bP_a\right)^2\leq I,
 \qquad
 \sum_a(P_a\otimes I)D^2(P_a\otimes I)\leq I,\nonumber\\
 D(P\otimes I)D\leq I,
 \qquad
 \sum_bPQ_bP\leq I,
 \qquad
 \sum_{a,b}P_aQ_bP_a\leq I.\label{eq:evaluated-g-normalizations}
\end{gather}
The first four bounds justify the four uses of
\eqref{eq:evaluated-g-consistency-processed} below.  The last one
justifies using \eqref{eq:evaluated-point-commutativity}, through
\Cref{prop:closeness-of-ip}, with coefficient $P_aQ_b\otimes I$.

Expanding the commutator, using projectivity, and exchanging
$(u,x,a)$ with $(v,y,b)$ gives
\[
 \mathcal C=2(S-T),
\]
where $\mathcal C$ is the commutator error in the conclusion and
\[
 S=\mathcal E(Q_bP_aQ_b\otimes I),
 \qquad
 T=\mathcal E(P_aQ_bP_aQ_b\otimes I).
\]
The preceding relations and estimates yield the following direct
comparison:
\begin{align*}
T
&\approx_{2\sqrt\zeta}
 \mathcal E(P_aQ_bP_aQ\otimes B_b)\\
&\approx_{\sqrt\zeta}
 \mathcal E(P_aQ_bP_a\otimes B_b)\\
&\approx_{2\sqrt\zeta}
 \mathcal E(P_aQ_bP\otimes B_bA_a)\\
&\approx_{\sqrt\zeta}
 \mathcal E(P_aQ_b\otimes B_bA_a)\\
&\approx_\kappa
 \mathcal E(P_aQ_b\otimes A_aB_b)\\
&\approx_{2\sqrt\zeta}
 \mathcal E(P_aQ_b\otimes B_b)\\
&\approx_{2\sqrt\zeta}
 \mathcal E(P_aQ_b\otimes I)\\
&\approx_{3\sqrt\zeta}S.
\end{align*}
The second and fourth transitions are
\eqref{eq:evaluated-g-first-leakage} and
\eqref{eq:evaluated-g-second-leakage}.  The fifth uses
\eqref{eq:evaluated-point-commutativity}.  The first, third, sixth,
and seventh use \eqref{eq:evaluated-g-consistency-processed}, in that order.
For the final transition, sum over~$a$ and apply
\Cref{prop:switch-sandwich} to the projective sub-measurement
$\{Q_b\}_b$ and the positive contraction~$P$, reading its two
approximations from right to left and using
\eqref{eq:evaluated-g-self-consistency-processed}.

Consequently,
\[
 |T-S|\leq13\sqrt\zeta+\kappa.
\]
Since $m\geq1$ implies $m+1\leq2m^2$, we conclude that
\begin{align*}
 \mathcal C
 &\leq26\sqrt\zeta+2\sqrt{32\gamma(m+1)}\\
 &\leq26\sqrt\zeta+16m\sqrt\gamma\\
 &\leq26m(\sqrt\gamma+\sqrt\zeta),
\end{align*}
as claimed.
\end{proof}

\subsection{\texorpdfstring{Commutativity of $G$}{Commutativity of the slice measurements}}\label{sec:G-commutes}

\begin{theorem}[Commutativity of~$G$]\label{thm:com-main}
  Let $m\geq1$, and let $(\psi, A, B, L)$ be an $(\eps, \delta, \gamma)$-good symmetric strategy for the $(m+1,q,d)$ low individual degree test.
  Let $\{G^x\}_{x \in \F_q}$ denote a set of projective sub-measurements in $\polysub{m}{q}{d}$ with the following properties:
  \begin{enumerate}
  \item (Consistency with~$A$): \label{item:commuting-consistency} On average over $(\bu, \bx) \sim \F_q^{m+1}$,
	  \begin{equation*}
	  A^{u, x}_a \otimes I \simeq_{\zeta} I \otimes G^x_{[g(u)=a]}.
	  \end{equation*}
  \item (Strong self-consistency): \label{item:commuting-self-consistency} On average over~$\bx \sim \F_q$,
	\begin{equation*}
		G^x_g \otimes I \approx_{\zeta} I \otimes G^x_g.
	\end{equation*}
  \item (Boundedness): \label{item:commuting-boundedness} There exists a positive-semidefinite matrix $Z^x$ for each $x \in \F_q$ such that
	\begin{equation*}
		\E_{\bx} \bra{\psi} (I-G^{\bx})\otimes Z^{\bx} \ket{\psi} \leq \zeta
	\end{equation*}
	and for each $x \in \F_q$ and $g \in \polyfunc{m}{q}{d}$,
	\begin{equation*}
	Z^x \geq \left(\E_{\bu} A^{\bu, x}_{g(\bu)}\right).
	\end{equation*}
  \end{enumerate}
  Let
  \begin{equation*}
  \nu = 30m \cdot \left(\gamma^{1/4}+ \zeta^{1/4} + (d/q)^{1/4}  \right).
  \end{equation*}
  Then on average over independent and uniformly random $(\bu,\bx), (\bv,\by) \sim \F_q^{m+1}$,
  \begin{equation*}
    G^{x}_{g} G^{y}_h \ot I \approx_{\nu} G^y_h G^x_g \ot I
  \end{equation*}
\end{theorem}

\begin{proof}
We first give a general lifting argument for projective
sub-measurements $P=\{P_i\}$ and $Q=\{Q_j\}$ on~$\calH$.
Write $X^{\mathrm L}=X\otimes I$ and
$X^{\mathrm R}=I\otimes X$, and define
\[
 \mathcal T(P)=\sum_i(P_i^{\mathrm L}-P_i^{\mathrm R})^2,
 \qquad \delta_P=\bra\psi\mathcal T(P)\ket\psi,
 \qquad
 C(P,Q)=\sum_{i,j}\|[P_i^{\mathrm L},Q_j^{\mathrm L}]\ket\psi\|^2.
\]
Suppose that evaluation maps $f_u$ satisfy
$\Pr_{\bu}[f_{\bu}(i)=f_{\bu}(k)]\leq\theta<1$ for $i\ne k$,
and set $P^u_a=\sum_{i:f_u(i)=a}P_i$.  Then
\begin{equation}\label{eq:commutation-disagreement-operator}
 (1-\theta)\mathcal T(P)
 \leq\E_{\bu}\mathcal T(P^{\bu})\leq\mathcal T(P).
\end{equation}
To verify this even when $P$ is incomplete, adjoin the failure
projection $P_\perp=I-\sum_iP_i$ solely to decompose the tensor
product into orthogonal blocks $P_i\otimes P_k$.
On unequal ordinary labels, the two operators
$\mathcal T(P)$ and $\E_{\bu}\mathcal T(P^{\bu})$ act by
$2$ and $2\Pr_{\bu}[f_{\bu}(i)\ne f_{\bu}(k)]$, respectively.
On an ordinary/failure block both act by~$1$; on equal labels and
on the failure/failure block both vanish.  This proves the operator
inequalities block by block.

Put
\[
 D(P,Q)=\sum_j\bra\psi Q_j^{\mathrm L}
                   \mathcal T(P)Q_j^{\mathrm L}\ket\psi.
\]
The identity
\[
 (P_i^{\mathrm L}-P_i^{\mathrm R})Q_j^{\mathrm L}\ket\psi
 -[P_i^{\mathrm L},Q_j^{\mathrm L}]\ket\psi
 =Q_j^{\mathrm L}(P_i^{\mathrm L}-P_i^{\mathrm R})\ket\psi
\]
and $\sum_jQ_j\leq I$ give, by the triangle inequality in the
direct sum over~$i,j$,
\begin{equation}\label{eq:commutation-mirror-comparison}
 |\sqrt{D(P,Q)}-\sqrt{C(P,Q)}|\leq\sqrt{\delta_P}.
\end{equation}
Applying \Cref{eq:commutation-disagreement-operator} inside~$D$
and then \Cref{eq:commutation-mirror-comparison} before and after
evaluation yields
\begin{align}
 \sqrt{C(P,Q)}
 &\leq\sqrt{D(P,Q)}+\sqrt{\delta_P}\nonumber\\
 &\leq\frac{\sqrt{\E_{\bu}C(P^{\bu},Q)}}{\sqrt{1-\theta}}
  +\left(1+\frac1{\sqrt{1-\theta}}\right)\sqrt{\delta_P}.
 \label{eq:commutation-one-sided-lifting}
\end{align}
Here Minkowski and
$\E_{\bu}\delta_{P^{\bu}}\leq\delta_P$ justify the last line.
The same argument includes any additional outer averages, by
including their probability weights in the direct sums.

Now apply this to $P=G^x$ and $Q=G^y$, with all quantities averaged
over independent uniform $x,y$.  Schwartz--Zippel gives
$\theta=md/q$ for the maps $g\mapsto g(u)$.
Both averaged mirror errors are at most~$\zeta$, and
\Cref{lem:comm-data-processed-g} gives
\begin{equation}\label{eq:com-main-evaluated-commutativity}
 \E_{\bu,\bv,\bx,\by}C((G^{\bx})^{\bu},(G^{\by})^{\bv})
 \leq\chi,
 \qquad \chi=26m(\sqrt\gamma+\sqrt\zeta).
\end{equation}
Apply \Cref{eq:commutation-one-sided-lifting} once to each family,
using $C(P,Q)=C(Q,P)$.  If $\theta<1$, this gives
\begin{equation}\label{eq:commutation-two-sided-lifting}
 \sqrt{\E_{\bx,\by}C(G^{\bx},G^{\by})}
 \leq\frac{\sqrt\chi}{1-\theta}
 +\left(1+\frac1{\sqrt{1-\theta}}\right)^2\sqrt\zeta.
\end{equation}
In particular, if $\theta\leq\tfrac12$, then
$(r+s)^2\leq2r^2+2s^2$ and
$2(1+\sqrt2)^4<68$ imply
\begin{equation}\label{eq:commutation-code-distance-bound}
 \E_{\bx,\by}C(G^{\bx},G^{\by})
 \leq8\chi+68\zeta
 =208m(\sqrt\gamma+\sqrt\zeta)+68\zeta.
\end{equation}

It remains to compare with the stated budget~$\nu$.
Sub-measurement normalization gives $C(P,Q)\leq4$ universally,
since each of
$\sum_{i,j}\|P_i^{\mathrm L}Q_j^{\mathrm L}\ket\psi\|^2$ and
$\sum_{i,j}\|Q_j^{\mathrm L}P_i^{\mathrm L}\ket\psi\|^2$
is at most~$1$.  Thus $\nu\geq4$ is immediate.
Otherwise, writing $a=\gamma^{1/4}$ and $b=\zeta^{1/4}$, we have
\[
 a,b<\frac2{15m},\qquad
 \theta=\frac{md}{q}
 <m\left(\frac2{15m}\right)^4<\frac12.
\]
Consequently, \Cref{eq:commutation-code-distance-bound} gives
\begin{align*}
 \E_{\bx,\by}C(G^{\bx},G^{\by})
 &\leq208m(a^2+b^2)+68b^4\\
 &\leq\frac{416}{15}a
   +\left(\frac{416}{15}+68\left(\frac2{15}\right)^3\right)b\\
 &\leq30m(a+b)\leq\nu.
\end{align*}
This is exactly the desired state-dependent commutation bound.
\end{proof}

\section{Pasting}\label{sec:ld-pasting}

A single positive contraction controls a whole sequence of independent
random slice measurements.  Its powers give both the cost of moving a
successful outcome to the start and the change in marginal success
probabilities.  Interpolation and a first-moment bound then finish the
argument.

\begin{theorem}[Pasting, quantitative form]\label{thm:ld-pasting}
Let $m\geq1$, $d\geq0$, and let $(\psi,A,B,L)$ be an
$(\eps,\delta,\gamma)$-good symmetric strategy for the $(m+1,q,d)$ test.
Suppose that the projective slice sub-measurements $\{G^x\}$ satisfy
the four completeness, consistency, strong self-consistency, and
boundedness conditions in \Cref{thm:ld-pasting-in-induction-section},
interpreted also when $d=0$.
For every integer $k\geq d+1$, there is a measurement
$H\in\polymeas{m+1}{q}{d}$ such that
\begin{equation}\label{eq:new-pasting-conclusion}
 A^u_a\otimes I\simeq_\sigma I\otimes H_{[h(u)=a]},
 \qquad
 \sigma=\frac{k}{k-d}\kappa+21k^2\sqrt m\,s,
\end{equation}
where the relation is averaged over $u\in\F_q^{m+1}$ and
\[
 s=\eps^{1/8}+\delta^{1/8}+\gamma^{1/8}
   +\zeta^{1/8}+(d/q)^{1/8}.
\]
\end{theorem}

We may assume $\eps,\delta,\gamma,\zeta,d/q<1$: otherwise the error
bound is at least one and a fixed polynomial measurement suffices.
Throughout this section write
\[
 \omega(X)=\bra\psi X\otimes I\ket\psi,\qquad
 G^x=\sum_gG_g^x,\qquad G=\E_xG^x,\qquad
 p=\omega(G),\qquad\kappa_0=1-p\leq\kappa.
\]
We use $B^u$ for the line measurement in the last coordinate, as in
\Cref{not:conditioned-on-last-direction}.  Conditioning the
axis-parallel test on that direction gives error at most $(m+1)\eps$.
Together with point self-consistency and
\Cref{prop:simeq-to-approx,prop:triangle-inequality-for-approx_delta},
this gives
\begin{equation}\label{eq:ld-abcon}
 I\otimes A^{u,x}_a\approx_D I\otimes B^u_{[f(x)=a]},
 \qquad D=8m\eps+4\delta.
\end{equation}
Symmetry and \Cref{prop:triangle-sub} then imply
\begin{equation}\label{eq:ld-gbcon}
 G^x_{[g(u)=a]}\otimes I\simeq_b I\otimes B^u_{[f(x)=a]},
 \qquad b=\zeta+\sqrt D.
\end{equation}
Finally, \Cref{thm:com-main} gives
\begin{equation}\label{eq:quote-com-main}
 \E_{x,y}\sum_{g,h}
 \|([G_g^x,G_h^y]\otimes I)\ket\psi\|^2\leq\chi,
 \qquad
 \chi=30m\bigl(\gamma^{1/4}+\zeta^{1/4}+(d/q)^{1/4}\bigr).
\end{equation}

\subsection{A positive contraction for random measurement words}
\label{sec:hat-consistency}

Set $Q_x=I-G^x$, and complete each slice measurement by
$P_g^x=G_g^x$ and $P_\bot^x=Q_x$.  Indices for $P$ include $\bot$;
indices for $G_g^x$ do not.  Define
\[
 \Pi_x=\sum_aP_a^x\otimes P_a^x,
 \qquad K=\E_x\Pi_x.
\]
Each $\Pi_x$ is a projection, so $0\leq K\leq I$.

\begin{lemma}[Completion]\label{lem:pasting-completion}
The total slice projectors and completed measurements satisfy
\begin{align}
 \E_x\|(G^x\otimes I-I\otimes G^x)\ket\psi\|^2
 &\leq\zeta,\label{eq:pasting-total-mirror}\\
 \E_x\sum_a\|(P_a^x\otimes I-I\otimes P_a^x)\ket\psi\|^2
 &=2\bra\psi(I-K)\ket\psi\leq\alpha:=2\zeta.
 \label{eq:gselfconall}
\end{align}
In particular, $\E_x\|\psi-\Pi_x\psi\|^2\leq\zeta$.
\end{lemma}
\begin{proof}
For fixed $x$, put
$e_x=\sum_g\|(G_g^x\otimes I-I\otimes G_g^x)\ket\psi\|^2$.
Since $G^x\otimes G^x\geq\sum_gG_g^x\otimes G_g^x$,
expanding the squares gives
\[
 \|(G^x\otimes I-I\otimes G^x)\ket\psi\|^2\leq e_x.
\]
Its average is at most $\zeta$.  The mirror difference for $Q_x$
is the negative of that for $G^x$, so the completed error is at
most $2e_x$.  For a complete projective measurement that error
is exactly $2\bra\psi(I-\Pi_x)\ket\psi
=2\|\psi-\Pi_x\psi\|^2$.  Averaging proves all the claims.
\end{proof}

For independent uniform questions $y_1,\ldots,y_j$ and their
outcomes $a_1,\ldots,a_j$, write
\[
 W=P_{a_j}^{y_j}\cdots P_{a_1}^{y_1}.
\]
Below, $\E\sum_W$ means averaging these questions and summing their
outcomes; for $j=0$ there is the single word $W=I$.  Successive
outcome summation gives both $\sum_WW^\dagger W=I$ and
$\sum_WWW^\dagger=I$ for every question tuple.  Independence of the
questions gives the exact operator identity
\[
 \E\sum_W W\otimes W=K^j.
\]
Set $M_W=W\otimes I-I\otimes W^\dagger$.  Expanding the square
gives the operator identity
\begin{equation}\label{eq:random-word-energy}
 \E\sum_W M_W^\dagger M_W
 =2(I-K^j)\leq2j(I-K).
\end{equation}
Both diagonal terms sum to $I$, and the cross terms average to
$K^j$ and $(K^j)^\dagger=K^j$.  The inequality follows from
$I-K^j\leq j(I-K)$.  In particular, this identity holds on every
vector, and \Cref{lem:pasting-completion} gives
\begin{equation}\label{eq:random-word-mirror}
 \E\sum_W\|M_W\ket\psi\|^2
 =2\bra\psi(I-K^j)\ket\psi\leq j\alpha.
\end{equation}

\begin{lemma}[Stability of marginal probabilities]\label{lem:sequential-bot-marginal}
Let $\rho$ be Alice's reduced state and
$\Phi(\rho)=\E_x\sum_aP_a^x\rho P_a^x$.
For every effect $0\leq M\leq I$ and integer $j\geq0$,
\begin{equation}\label{eq:pasting-marginal-stability}
 |\tr(M\Phi^j(\rho))-\tr(M\rho)|\leq\sqrt{j\alpha}.
\end{equation}
\end{lemma}
\begin{proof}
In the weighted direct sum over words, put
$U_W=(W\otimes I)\psi$ and $V_W=(I\otimes W^\dagger)\psi$.
Both vectors have norm one.  Their $M\otimes I$ expectations are
$\tr(M\Phi^j(\rho))$ and $\tr(M\rho)$, respectively; the second
identity uses $\E\sum_WWW^\dagger=I$.
Because the norms agree, replacing $M$ by $M-I/2$ preserves their
expectation difference.  Cauchy--Schwarz and
\eqref{eq:random-word-mirror} bound its modulus by
$\tfrac12\|U-V\|(\|U\|+\|V\|)\leq\sqrt{j\alpha}$.
\end{proof}

\begin{lemma}[Moving a successful outcome]\label{lem:commute-g-half-sandwich}
For a length-$j$ random word $W$ and an independent uniform $x$, let
\[
 D_j=\E_x\E\sum_W\sum_g
       \|([G_g^x,W]\otimes I)\psi\|^2.
\]
Then both of the following bounds hold:
\begin{align}
 D_j&\leq2j(\sqrt\chi+2\sqrt\zeta)^2,
       \label{eq:random-word-commutation-sharp}\\
 D_j&\leq2j(\sqrt\chi+4\sqrt\zeta).
       \label{eq:random-word-commutation}
\end{align}
\end{lemma}
\begin{proof}
Write $P=G_g^x\otimes I$ and $\Delta_h^y=G_h^y\otimes I-I\otimes G_h^y$.
The operator inequality
\[
 \sum_h(\Delta_h^y)^2-(I-\Pi_y)
 =\sum_{h\ne h'}G_h^y\otimes G_{h'}^y\geq0
\]
controls the energy
$\beta=\E_x\sum_g\bra\psi P(I-K)P\ket\psi$.
Indeed,
\[
 \beta\leq\E_{x,y}\sum_{g,h}\|\Delta_h^yP\psi\|^2,
 \qquad
 \Delta_h^yP\psi
 =([G_h^y,G_g^x]\otimes I)\psi+P\Delta_h^y\psi.
\]
The direct-sum triangle inequality, fine commutation, and
$\sum_g(G_g^x)^2=G^x\leq I$ give
$\sqrt\beta\leq\sqrt\chi+\sqrt\zeta$; also $\beta\leq p\leq1$.

Since $P$ commutes with $I\otimes W^\dagger$,
\[
 ([G_g^x,W]\otimes I)\psi=PM_W\psi-M_WP\psi.
\]
The two direct-sum norms on the right are at most
$\sqrt{2j\zeta}$ and $\sqrt{2j\beta}$, respectively, by
\eqref{eq:random-word-energy}, \Cref{lem:pasting-completion}, and
sub-measurement normalization.  Hence
\[
 \sqrt{D_j}\leq\sqrt{2j}(\sqrt\zeta+\sqrt\beta),
\]
which proves \eqref{eq:random-word-commutation-sharp}.
For the second bound, use $\zeta<1$ and
$\beta\leq\min\{1,(\sqrt\chi+\sqrt\zeta)^2\}
\leq\sqrt\chi+\sqrt\zeta$ to obtain
$(\sqrt\zeta+\sqrt\beta)^2\leq\beta+3\sqrt\zeta
\leq\sqrt\chi+4\sqrt\zeta$.
\end{proof}

For coarse error estimates we use
\begin{equation}\label{eq:pasting-root-budgets}
 \sqrt\chi+4\sqrt\zeta\leq(\sqrt{30}+4)\sqrt m\,s,
 \qquad b\leq3\sqrt m\,s,
 \qquad\sqrt D\leq3\sqrt m\,s.
\end{equation}
These follow by subadditivity of the square root and the standing
assumption that all error parameters are below one.

\subsection{Interpolation and consistency}
\label{sec:ld-sandwiching}

For independent uniform $x_1,\ldots,x_k$ and completed outcomes
$a_1,\ldots,a_k$, the sequential effect is
\begin{equation}\label{eq:sequential-povm}
 E_a^x=P_{a_1}^{x_1}\cdots P_{a_k}^{x_k}
             P_{a_k}^{x_k}\cdots P_{a_1}^{x_1}.
\end{equation}
It is positive, and summing the outcomes from the inside out gives
$\sum_aE_a^x=I$.  Scan the successful indices in order, retaining
only the first success at each distinct slice coordinate.  If
$j_1<\cdots<j_{d+1}$ are the first $d+1$ retained indices, output
\begin{equation}\label{eq:first-success-interpolation}
 h_{x,a}(u,t)=\sum_{\ell=1}^{d+1}a_{j_\ell}(u)
       \prod_{\substack{v=1\\v\ne\ell}}^{d+1}
       \frac{t-x_{j_v}}{x_{j_\ell}-x_{j_v}}.
\end{equation}
Otherwise output nothing.  The denominators are nonzero, and this
polynomial has individual degree at most $d$ in every coordinate.
It is the unique such polynomial extending the selected slices.
In particular, no condition is imposed on later successful outcomes.
Define the resulting sub-measurement by
\begin{equation}\label{eq:first-success-pasted-measurement}
 T_{k,h}=\E_x\sum_{a:\,h_{x,a}\text{ defined and }h_{x,a}=h}E_a^x,
 \qquad T_k=\sum_hT_{k,h}.
\end{equation}
This is a post-processing of a measurement, hence a sub-measurement.
The construction is defined even for $k>q$.
Set $L_k=\sum_{j=0}^{k-1}\sqrt j$, so $L_1=0$ and
$L_k\leq k(k-1)/2$.

\begin{lemma}[Consistency]\label{lem:pasting-consistency}
Let $\eta_k$ be the inconsistency of $\{T_{k,h}\}$ with $A$.
Then
\begin{equation}\label{eq:pasting-consistency-budget}
 \eta_k\leq2kb+2k(k-1)(\sqrt\chi+4\sqrt\zeta)+\sqrt D
       \leq19k^2\sqrt m\,s.
\end{equation}
\end{lemma}
\begin{proof}
Fix $i$, sum out later outcomes, and write
$W=P_{a_{i-1}}^{x_{i-1}}\cdots P_{a_1}^{x_1}$,
$P=G_g^{x_i}$ for a successful outcome, and
$C=I-B^u_{[f(x_i)=g(u)]}$.  In the direct sum over $u,x,a_{<i},g$,
compare
\[
 v=(PW\otimes C^{1/2})\psi,\qquad
 v'=(WP\otimes C^{1/2})\psi.
\]
The squared norm of $v$ is the successful-mismatch probability
$S_i$ at position $i$.  Since $C$ is an effect,
\Cref{lem:commute-g-half-sandwich} gives
$\|v-v'\|\leq\sqrt{D_{i-1}}$.
Summing the earlier outcomes in $\|v'\|^2$ leaves the original
slice--line inconsistency, so $\|v'\|^2\leq b$ by
\eqref{eq:ld-gbcon}.  Hence
\[
 S_i\leq(\sqrt b+\sqrt{D_{i-1}})^2\leq2b+2D_{i-1}.
\]
If Bob's degree-$d$ polynomial $f$ differs from
$h_{x,a}(u,\cdot)$, it disagrees at some selected one of the $d+1$
distinct coordinates.  A union bound over all successful positions
bounds the line inconsistency by
\[
 \sum_{i=1}^kS_i
 \leq2kb+2k(k-1)(\sqrt\chi+4\sqrt\zeta).
\]
Evaluation at an independent uniform point and
\Cref{prop:triangle-sub,eq:ld-abcon} add at most $\sqrt D$.
Symmetry gives either tensor orientation.
Finally \eqref{eq:pasting-root-budgets} bounds this error by
\[
 \bigl[6k+2(\sqrt{30}+4)k(k-1)+3\bigr]
       \sqrt m\,s
 \leq19k^2\sqrt m\,s.
\]
\end{proof}

\subsection{Completeness from the expected number of successes}

\begin{lemma}[Completeness]\label{lem:first-success-completeness}
Put $c_k=\bone[d\geq1]\,k(k-1)/(2q)$.  Then
\begin{equation}\label{eq:first-success-completeness}
 \omega(T_k)\geq
 1-\frac{k}{k-d}\kappa_0
   -\frac{\sqrt\alpha L_k}{k-d}-c_k.
\end{equation}
\end{lemma}
\begin{proof}
Let $N$ count successful outcomes.  By
\Cref{lem:sequential-bot-marginal}, its $i$-th success probability is
at least $p-\sqrt{(i-1)\alpha}$, so
$\E N\geq kp-\sqrt\alpha L_k$.
The pointwise inequality
$N\leq d+(k-d)\bone[N\geq d+1]$ implies
$\Pr[N\geq d+1]\geq(\E N-d)/(k-d)$.
If the slice coordinates are distinct, $d+1$ successes suffice for
interpolation.  The probability of a coordinate collision is at
most $\binom{k}{2}/q$ by a union bound; when $d=0$, a single success
suffices and no collision can matter.  Subtracting this probability
gives the claim.  Only the questions are independent; no independence
between outcomes is used.
\end{proof}

\begin{proof}[Proof of \Cref{thm:ld-pasting}]
Choose any polynomial $h^*$ and complete $T_k$ by
$H_h=T_{k,h}+\bone[h=h^*](I-T_k)$.
Its inconsistency is at most $\eta_k+1-\omega(T_k)$: the additional
bad-outcome probability is bounded by the entire missing mass.
Write $\operatorname{Inc}(A,H)$ for this inconsistency.
The preceding lemmas give the explicit estimate
\begin{equation}\label{eq:positive-contraction-pasting}
 \begin{split}
 \operatorname{Inc}(A,H)\leq{}&\frac{k}{k-d}\kappa+2kb
       +2k(k-1)(\sqrt\chi+4\sqrt\zeta)\\
 &+\frac{\sqrt{2\zeta}\,L_k}{k-d}
       +\sqrt D+c_k.
 \end{split}
\end{equation}
By \eqref{eq:pasting-consistency-budget}, the total error is at most
\[
 \frac{k}{k-d}\kappa+19k^2\sqrt m\,s
       +\frac{\sqrt{2\zeta}L_k}{k-d}+c_k.
\]
Since $k-d\geq1$, the marginal term is at most
$k^2\sqrt m\,s$.  For $d\geq1$,
$c_k\leq k^2(d/q)/2\leq k^2\sqrt m\,s/2$; for $d=0$ it vanishes.
The fresh error is therefore at most $20.5k^2\sqrt m\,s$, which is
bounded by the stated $21k^2\sqrt m\,s$.
Choosing $k=(m+1)d$ when $m,d\geq1$ proves
\Cref{thm:ld-pasting-in-induction-section}.
For $d=0$ one may take $k=1$; then $L_k=c_k=0$, there is no
commutation term, and the uncompleted sub-measurement $T_1$ has
total effect exactly $G$.
\end{proof}

\subsection{The second moment and the original construction}
\label{sec:first-construction}

The induction uses the preceding construction for its stronger
error exponent.  The same estimates also certify the original
procedure: perform exactly $d+1$ slice sub-measurements at uniformly
distinct coordinates, and interpolate if they all succeed.
Call its sub-measurement $H^{(1)}$.

\begin{lemma}[Boundedness controls the second moment]\label{lem:pasting-second-moment}
If a global polynomial sub-measurement $\{T_h\}$ has total effect $T$
and inconsistency $\eta$ with $A$, then
\begin{equation}\label{eq:pasting-support}
 |\omega(T(I-G))|\leq\sqrt\eta+\sqrt\zeta,\qquad
 \omega(G-G^2)\leq p-\omega(T)+2\sqrt\eta+2\sqrt\zeta.
\end{equation}
\end{lemma}
\begin{proof}
Abbreviate $A_h=A^{u,x}_{h(u,x)}$ and split
\[
 \omega(T(I-G))=\E_{u,x}\sum_h\bra\psi T_hQ_x\otimes A_h\ket\psi
       +\E_{u,x}\sum_h\bra\psi T_hQ_x\otimes(I-A_h)\ket\psi.
\]
For the first term, Cauchy--Schwarz with vector factors
$\sqrt{T_h}\otimes\sqrt{A_h}$ and
$\sqrt{T_h}Q_x\otimes\sqrt{A_h}$ gives one squared norm at most one;
the other is bounded by \Cref{lem:boundedness-rejection}:
\[
 \E_{u,x}\sum_h\bra\psi Q_xT_hQ_x\otimes A_h\ket\psi
 \leq\zeta.
\]
Indeed, $\{T_h\}$ is independent of~$u$ and sums to at most~$I$,
each restriction $h|_x$ is an admissible slice polynomial, and hence
$\E_u A^{u,x}_{h(u,x)}\leq Z^x$.  The slice hypothesis
$\E_x\bra\psi Q_x\otimes Z^x\ket\psi\leq\zeta$ completes the
application of \eqref{eq:boundedness-rejection-average}.
Replacing $A_h$ by $I-A_h$ gives squared
norms at most $\eta$ and one for the second term.  This proves the
first inequality, without any norm bound or commutation assumption
on $Z^x$.  For the second, take expectations in
\[
 G^2-GT-TG+T=(G-T)^2+T-T^2\geq0
\]
and rearrange, using the real part of $\omega(T(I-G))$.
\end{proof}

\begin{corollary}[Pasting with exactly \(d+1\) slices]\label{cor:first-construction}
Suppose $d\geq1$ and $q>d$.  Set
\[
 S=\eps^{1/32}+\delta^{1/32}+\gamma^{1/32}
     +\zeta^{1/32}+(d/q)^{1/32}.
\]
Then
\[
 \omega(G-G^2)\leq40mdS,\qquad
 0\leq p-\omega(H^{(1)})\leq45md^2S.
\]
Completing $H^{(1)}$ gives inconsistency at most $\kappa+125md^2S$.
For $d=0$, the one-slice construction has total effect exactly $G$.
\end{corollary}
\begin{proof}
Write $V=\omega(G-G^2)$.  Apply \Cref{lem:pasting-second-moment}
to $T_k$, using its established consistency and completeness:
\begin{equation}\label{eq:direct-second-moment-tradeoff}
 V\leq\frac{d\kappa_0}{k-d}
  +\frac{\sqrt{2\zeta}L_k}{k-d}
  +\frac{k(k-1)}{2q}+2\sqrt{\eta_k}+2\sqrt\zeta.
\end{equation}
For explicit bounds put $R=s$.  If $R>1$, the claims are trivial;
otherwise $0<R\leq1$ and we choose $k=\lceil2dR^{-1/4}\rceil$.
Thus $2d\leq k\leq3dR^{-1/4}$.
Using $\eta_k\leq19k^2\sqrt m R$, $\kappa_0\leq1$,
$\sqrt{2\zeta}\leq2R$,
$L_k\leq k(k-1)/2$, and $d/q\leq R^8$,
the five terms in \eqref{eq:direct-second-moment-tradeoff} are at most
\[
 R^{1/4},\quad6dR^{3/4},\quad
 \tfrac92dR^{15/2},\quad6\sqrt{19}\,d m^{1/4}R^{1/4},\quad2R.
\]
Since $13.5+6\sqrt{19}<40$, their sum is at most
$40mdR^{1/4}\leq40mdS$.
This argument remains valid for $k>q$, since $T_k$ is defined for
every $k$ and the collision bound is still an upper bound.

To analyze the actual first construction, let
\[
 L=\E_{x,y}\sum_g\omega(G_g^xQ_yG_g^x)
\]
be the probability of a success followed immediately by rejection.
Two mirror replacements in the quadratic form of $G$ give
\[
 \left|\E_x\sum_g\omega(G_g^xGG_g^x)
       -\bra\psi G\otimes G\ket\psi\right|\leq2\sqrt\zeta.
\]
Indeed, the aggregate squared distance between
$(G_g^x\otimes I)\psi$ and $(I\otimes G_g^x)\psi$ is at most
$\zeta$, and each aggregate norm is at most one.
By \eqref{eq:pasting-total-mirror} and Jensen,
$|\bra\psi G\otimes G\ket\psi-\omega(G^2)|\leq\sqrt\zeta$.
Consequently $|L-V|\leq3\sqrt\zeta$.

In $d+1$ completed trials, if the first succeeds but some later
trial fails, there is an adjacent success-to-rejection transition.
The transition beginning at position $i$ has probability at most
$L+\sqrt{(i-1)\alpha}$ by
\eqref{eq:pasting-marginal-stability}, applied to its two-trial effect.
Summing these probabilities and changing independent coordinates to
uniformly distinct coordinates gives
\begin{equation}\label{eq:first-construction-completeness}
 p-\omega(H^{(1)})\leq
 dV+3d\sqrt\zeta+\sqrt{2\zeta}L_d
       +\frac{d(d+1)}{2q}.
\end{equation}
The distribution change costs at most the probability of a collision,
by conditioning the independent sample on having none.
Also $H^{(1)}\leq G$: for a fixed tuple the total successful effect
is bounded by its first success projector, and the first coordinate
is uniform under distinct sampling.  This proves the lower bound.
Using $L_d\leq d(d-1)/2$, the terms beyond $dV$ in
\eqref{eq:first-construction-completeness} are at most $5d^2S$,
proving the claimed upper bound.

For consistency, put $n=d+1$ and let $a>0$ be the probability that
$n$ independent uniform coordinates are distinct.  With only $n$
trials, $T_n$ outputs a polynomial exactly when every outcome
succeeds and every coordinate is distinct.  Hence
$T_{n,h}=aH_h^{(1)}$ for every $h$.  If $e$ denotes the
inconsistency of $H^{(1)}$ with $A$, then $ae=\eta_n$ and $e\leq1$,
so
\[
 e=\eta_n+(1-a)e
 \leq\eta_n+\frac{d(d+1)}{2q}
 \leq19(d+1)^2\sqrt m\,s+\frac{d(d+1)}{2q}
 \leq77md^2S.
\]
Completion adds at most $\kappa_0+45md^2S$, so the resulting
inconsistency is at most $\kappa+122md^2S\leq\kappa+125md^2S$.
\end{proof}

\bibliographystyle{alpha}
\newcommand{\etalchar}[1]{$^{#1}$}
\begin{thebibliography}{ALM{\etalchar{+}}98}

\bibitem[ALM{\etalchar{+}}98]{ALM+98}
Sanjeev Arora, Carsten Lund, Rajeev Motwani, Madhu Sudan, and Mario Szegedy.
\newblock Proof verification and the hardness of approximation problems.
\newblock {\em Journal of the ACM}, 45(3):501--555, 1998.

\bibitem[Ara02]{aravind2002simple}
PK~Aravind.
\newblock A simple demonstration of {B}ell's theorem involving two observers
  and no probabilities or inequalities.
\newblock {\em arXiv preprint quant-ph/0206070}, 2002.

\bibitem[AS98]{AS98}
Sanjeev Arora and Shmuel Safra.
\newblock Probabilistic checking of proofs: a new characterization of {NP}.
\newblock {\em Journal of the ACM}, 45(1):70--122, 1998.

\bibitem[BFL91]{BFL91}
L{\'a}szl{\'o} Babai, Lance Fortnow, and Carsten Lund.
\newblock Non-deterministic exponential time has two-prover interactive
  protocols.
\newblock {\em Computational complexity}, 1(1):3--40, 1991.

\bibitem[BLR93]{BLR93}
Manuel Blum, Michael Luby, and Ronitt Rubinfeld.
\newblock Self-testing/correcting with applications to numerical problems.
\newblock {\em Journal of computer and system sciences}, 47(3):549--595, 1993.

\bibitem[CMS17]{CMS17}
Alessandro Chiesa, Peter Manohar, and Igor Shinkar.
\newblock On axis-parallel tests for tensor product codes.
\newblock In {\em Proceedings of the 21st Annual International Workshop on
  Randomization and Computation}, 2017.

\bibitem[dlS22]{deLaSalle22}
Mikael de~la~Salle.
\newblock Orthogonalization of positive operator valued measures.
\newblock {\em Comptes Rendus Math\'ematique}, 360:549--560, 2022.
\newblock doi:10.5802/crmath.326.

\bibitem[FHS94]{FHS94}
Katalin Friedl, Zsolt Hatsagi, and Alexander Shen.
\newblock Low-degree tests.
\newblock In {\em Proceedings of the 5th Annual ACM-SIAM Symposium on Discrete
  Algorithms}, pages 57--64, 1994.

\bibitem[IV12]{IV12}
Tsuyoshi Ito and Thomas Vidick.
\newblock A multi-prover interactive proof for {NEXP} sound against entangled
  provers.
\newblock In {\em Proceedings of the 53rd Annual IEEE Symposium on Foundations
  of Computer Science}, pages 243--252, 2012.

\bibitem[JNV{\etalchar{+}}20]{JNV+20}
Zhengfeng Ji, Anand Natarajan, Thomas Vidick, John Wright, and Henry Yuen.
\newblock $\mathsf{MIP}^* = \mathsf{RE}$.
\newblock Technical report, arXiv:2001.04383, 2020.

\bibitem[KS08]{KS08}
Tali Kaufman and Madhu Sudan.
\newblock Algebraic property testing: the role of invariance.
\newblock In {\em Proceedings of the 40th Annual ACM Symposium on Theory of
  Computing}, pages 403--412, 2008.

\bibitem[KV11]{KV11}
Julia Kempe and Thomas Vidick.
\newblock Parallel repetition of entangled games.
\newblock In {\em Proceedings of the 43rd Annual ACM Symposium on Theory of
  Computing}, pages 353--362, 2011.

\bibitem[Mer90]{mermin1990simple}
David Mermin.
\newblock Simple unified form for the major no-hidden-variables theorems.
\newblock {\em Physical Review Letters}, 65(27):3373, 1990.

\bibitem[NV18a]{NV18a}
Anand Natarajan and Thomas Vidick.
\newblock Low-degree testing for quantum states, and a quantum entangled games
  {PCP}.
\newblock In {\em Proceedings of the 59th Annual IEEE Symposium on Foundations
  of Computer Science}, 2018.

\bibitem[NV18b]{NV18b}
Anand Natarajan and Thomas Vidick.
\newblock Two-player entangled games are {NP}-hard.
\newblock In {\em Proceedings of the 33rd Annual IEEE Conference on
  Computational Complexity}, 2018.

\bibitem[NW19]{NW19}
Anand Natarajan and John Wright.
\newblock $\mathsf{NEEXP} \subseteq \mathsf{MIP}^*$.
\newblock In {\em Proceedings of the 60th Annual IEEE Symposium on Foundations
  of Computer Science}, pages 510--518, 2019.

\bibitem[Per90]{peres1990incompatible}
Asher Peres.
\newblock Incompatible results of quantum measurements.
\newblock {\em Physics Letters A}, 151(3-4):107--108, 1990.

\bibitem[PS94]{PS94}
Alexander Polishchuk and Daniel Spielman.
\newblock Nearly-linear size holographic proofs.
\newblock In {\em Proceedings of the 26th Annual ACM Symposium on Theory of
  Computing}, pages 194--203, 1994.

\bibitem[RS97]{RS97}
Ran Raz and Shmuel Safra.
\newblock A sub-constant error-probability low-degree test, and a sub-constant
  error-probability {PCP} characterization of {NP}.
\newblock In {\em Proceedings of the 29th Annual ACM Symposium on Theory of
  Computing}, pages 475--484, 1997.

\bibitem[Sch80]{Sch80}
Jacob Schwartz.
\newblock Fast probabilistic algorithms for verification of polynomial
  identities.
\newblock {\em Journal of the ACM}, 27(4):701--717, 1980.

\bibitem[Sud11]{Sud11}
Madhu Sudan.
\newblock Guest column: testing linear properties: some general theme.
\newblock {\em ACM SIGACT News}, 42(1):59--80, 2011.

\bibitem[Vid11]{Vid11}
Thomas Vidick.
\newblock {\em The complexity of entangled games}.
\newblock PhD thesis, University of California, Berkeley, 2011.

\bibitem[Vid16]{Vid16}
Thomas Vidick.
\newblock Three-player entangled {XOR} games are {NP}-hard to approximate.
\newblock {\em SIAM Journal on Computing}, 45(3):1007--1063, 2016.

\bibitem[Zip79]{Zip79}
Richard Zippel.
\newblock Probabilistic algorithms for sparse polynomials.
\newblock In {\em Proceedings of the 2nd International Symposium on Symbolic
  and Algebraic Manipulation}, pages 216--226, 1979.

\end{thebibliography}

\end{document}
